\documentclass[a4paper,10pt,reqno]{amsart}

%------------------------------------------------
% Basic formatting
%------------------------------------------------
\usepackage[english]{babel}
\usepackage[
    a4paper,
    left=2.5cm,
    right=2.5cm,
    top=2.5cm,
    bottom=3cm]{geometry}
\usepackage{setspace}
\usepackage{fancyhdr}
\usepackage{csquotes}
\usepackage{enumitem}

%------------------------------------------------
% Section formatting
%------------------------------------------------
\usepackage{titlesec}

\titleformat{\section}
    {\normalfont\Large\bfseries}
    {\thesection.}
    {0.5em}
    {}
\titlespacing*{\section}
    {0pt}    % left indentation
    {1.5em}    % space before section
    {0pt}  % space after section

\titleformat{\subsection}
    {\normalfont\large\bfseries}
    {\bfseries\thesubsection.}
    {0.5em}
    {}
\titlespacing*{\subsection}
    {0pt}    % left indentation
    {0.5em}    % space before
    {-0.5em}  % space after

\titleformat{\paragraph}[runin]
    {\normalfont\normalsize\bfseries}
    {}
    {0pt}
    {}

\titlespacing*{\paragraph}
    {0pt}
    {0.5em}
    {1em}

%------------------------------------------------
% Mathematics
%------------------------------------------------
\usepackage{amsmath,amssymb,amsfonts,amsthm}
\usepackage{mathtools}
\allowdisplaybreaks
\usepackage{mathrsfs}
\usepackage{bm}

%------------------------------------------------
% Theorem environments
%------------------------------------------------

% Arguments:
% 1. Space above
% 2. Space below

\newtheoremstyle{myplain}
  {1em}        % Space above
  {0em}       % Space below
  {\itshape}     % Body font
  {}             % Indent
  {\bfseries}    % Head font
  {.}            % Punctuation
  {0.5em}        % Space after head
  {}             % Head spec

\newtheoremstyle{mydefinition}
  {1em}
  {0em}
  {\normalfont}
  {}
  {\bfseries}
  {.}
  {0.5em}
  {}

\newtheoremstyle{myremark}
  {1em}
  {0em}
  {\normalfont}
  {}
  {\itshape}
  {.}
  {0.5em}
  {}

% Theorems, lemmas, propositions, corollaries
\theoremstyle{myplain}
\newtheorem{theorem}{Theorem}[section]
\newtheorem{lemma}[theorem]{Lemma}
\newtheorem{proposition}[theorem]{Proposition}
\newtheorem{corollary}[theorem]{Corollary}

% Definitions, assumptions, examples
\theoremstyle{mydefinition}
\newtheorem{definition}[theorem]{Definition}
\newtheorem{assumption}[theorem]{Assumption}
\newtheorem{example}[theorem]{Example}

% Remarks
\theoremstyle{myremark}
\newtheorem{remark}[theorem]{Remark}

% Compact proof spacing
\makeatletter
\renewenvironment{proof}[1][\proofname]{%
  \par
  \pushQED{\qed}%
  \normalfont
  \topsep=0em
  \trivlist
  \item[\hskip\labelsep\itshape #1\@addpunct{.}]%
  \ignorespaces
}{%
  \popQED
  \endtrivlist
  \@endpefalse
}
\makeatother

%------------------------------------------------
% Graphics
%------------------------------------------------
\usepackage{graphicx}
\usepackage[dvipsnames]{xcolor}
\usepackage{tikz}
\usepackage{pgfplots}
\pgfplotsset{compat=1.18}

%------------------------------------------------
% Page numbering
%------------------------------------------------
\pagestyle{fancy}
\fancyhf{} % clear all header and footer fields

\fancyfoot[C]{\thepage} % centred page number at bottom
\setlength{\footskip}{1cm}

\renewcommand{\headrulewidth}{0pt}
\renewcommand{\footrulewidth}{0pt}

%------------------------------------------------
% Hyperlinks
%------------------------------------------------
\usepackage{hyperref}
\hypersetup{
    colorlinks=true,
    linkcolor=black,
    citecolor=black,
    urlcolor=black
}

%------------------------------------------------
% Bibliography
%------------------------------------------------
\usepackage[
    backend=biber,
    style=alphabetic
]{biblatex}
\addbibresource{references.bib}

%------------------------------------------------
% Paragraph formatting
%------------------------------------------------
\linespread{1.2}
\setlength{\parindent}{0pt}
\setlength{\parskip}{1em}

%------------------------------------------------
% Title page dimensions
%------------------------------------------------
\newlength{\logowidth}
\setlength{\logowidth}{70mm}

\newlength{\logopad}
\setlength{\logopad}{17.5mm}

\begin{document}

%------------------------------------------------
% Title page
%------------------------------------------------
\begin{titlepage}
    \pgfdeclareimage[width=\logowidth]{logo}{logo}
    \pgftext[at=\pgfpoint{-1.75cm+\logopad}{3cm-\logopad},
             top,left]
             {\pgfuseimage{logo}}

    \begin{center}
        \vspace*{\fill}

        {\large \scshape{MSc APPLIED MATHEMATICS RESEARCH PROJECT}}\\[0.8cm]

        \rule{\linewidth}{0.5mm}\\[0.4cm]

        {\LARGE\bfseries
        Functional Autoregressive Transformers as Markov Chains\par}

        \rule{\linewidth}{0.5mm}

        \vskip 2.5cm

        Author\\[3mm]
        {\large\bfseries Van Thuan Romoli}\\[1mm]
        {\normalsize CID: 06047418}

        \vskip 1.5cm

        Supervised by\\[3mm]
        {\large\bfseries Dr Nicolas Boullé}

        \vspace*{\fill}

        {\small\textit{The work contained in this thesis is my own work unless otherwise stated}}\\[-0.5mm]
        {\small
        Department of Mathematics
        $\cdot$
        Imperial College London
        $\cdot$
        2026}
    \end{center}
\end{titlepage}

%------------------------------------------------
% Table of contents formatting
%------------------------------------------------
\makeatletter
\def\l@section{\@tocline{1}{-2pt}{1pc}{1.5pc}{\bfseries}}
\def\l@subsection{\@tocline{2}{-4pt}{2.5pc}{2pc}{}}
\makeatother

%------------------------------------------------
% Enumerate formatting
%------------------------------------------------
\setlist[enumerate]{
    topsep=0pt,
    partopsep=0pt,
    parsep=0pt
}

% Your document starts here

\pagenumbering{roman}

\thispagestyle{empty}

\vspace*{\fill}

\begin{center}
    {\scshape Abstract}

\begin{minipage}{0.9\textwidth}

Many scientific machine learning problems involve sequences of function-valued observations. Examples include physical trajectories, spatial fields, demand curves, spectra, and continuous signals.
Existing operator-learning methods largely focus on deterministic prediction, leaving open questions about probabilistic autoregressive models for recursive generation in function space.
We develop a probabilistic autoregressive transformer for this setting and provide a rigorous framework for recursive generation. We show that the model defines a valid stochastic process which, under suitable conditions, converges to a unique invariant law. Furthermore, we study how well a fixed model performs when the training observations are dependent. Two complementary approaches provide generalisation and sample-complexity bounds under several sampling regimes, including independent trajectories, mixing sequences, and continuous functional dynamics. Overall, the thesis answers three questions: whether the model is well-defined, whether recursive generation is stable, and whether the performance of a fixed model can be assessed reliably from dependent data.
 \end{minipage}
\end{center}



\vspace*{\fill}


\clearpage

\tableofcontents
\thispagestyle{empty}

\clearpage

\pagenumbering{arabic}

\section{Introduction}\label{sec:Introduction}

\paragraph{Motivation.}
Many scientific data sets are sequences of functions: examples include trajectories of partial differential equations, spatial fields observed over time, demand curves, spectra, and continuously sampled signals. Treating each observation as a function preserves geometry and regularity that may be obscured by flattening a discretisation into a vector. Functional data analysis provides the statistical language for this viewpoint \cite{RamsaySilverman2005}, while functional autoregressive models and functional time-series methods provide probabilistic foundations for prediction under temporal dependence \cite{Bosq2000,HormannKokoszka2010,AueEtAl2015}. Alongside this classical theory, transformers provide an expressive finite-context mechanism for processing several recent observations. The present thesis studies how these perspectives can be combined and how transformers can be defined so that each functional context determines a conditional probability law for the next function and recursive sampling generates a stochastic process. This raises questions that are not present in a finite-dimensional implementation: the conditional output law must be mathematically admissible and measurable, the shift-and-append update must define a well-posed Markov chain, and its invariant behaviour and moments must be controlled. A separate statistical question is how the transition risk of a fixed model parameter can be estimated from dependent functional observations.

\paragraph{Literature review.}
Probabilistic sequence models describe sequential data through conditional distributions for future observations given their preceding context. Autoregressive transformers use self-attention to construct context-dependent representations of the preceding sequence \cite{VaswaniEtAl2017}. At scale, they exhibit in-context learning, where predictions adapt to information or examples supplied in the context without updating the model parameters \cite{BrownEtAl2020}. Theoretical accounts include Bayesian interpretations \cite{XieEtAl2022}, learnable function classes \cite{GargEtAl2022,AkyurekEtAl2023}, gradient-based and algorithmic interpretations \cite{VonOswaldEtAl2023,LiEtAl2023ICL,ChengEtAl2024}, and information-theoretic decompositions \cite{JeonEtAl2024}. Markovian analyses relate attention to context-conditioned Markov models \cite{IldizEtAl2024}, study induction heads on Markov-chain prompts \cite{EdelmanEtAl2024}, and, most closely, represent finite-context language generation as a finite-state Markov chain by taking the full token context as the state \cite{ZekriEtAl2025}. This state augmentation also underlies the present thesis, but the function-valued setting requires Hilbert-space admissibility and measurability arguments that do not arise in the same form in the finite-state construction.

In parallel, functional data analysis and functional time-series theory study function-valued observations under temporal dependence \cite{RamsaySilverman2005,Bosq2000,HormannKokoszka2010,AueEtAl2015}. Operator learning approximates maps between function spaces; for example, DeepONet, the Fourier neural operator, and neural-operator theory develop this perspective \cite{LuEtAl2021,LiEtAl2021FNO,KovachkiEtAl2023}. The mathematical, numerical foundations, and statistical theory have also been investigated \cite{BoulleTownsend2024,ReinhardtEtAl2026}. More recent developments have sought to combine operator learning with attention and in-context learning, enabled in part by continuum formulations of attention \cite{CalvelloEtAl2025}. Attention-based variants include Galerkin and neural operator transformers \cite{Cao2021,HaoEtAl2023}, while in-context operator learning uses functional demonstrations to infer operators without parameter updates, including continuum-transformer interpretations as operator gradient descent \cite{YangEtAl2023ICON,YangOsher2024,MishraEtAl2026}. However, probabilistic sequence modelling in function space raises a number of open questions concerning well-posedness, recursive generation, and long-term stochastic behaviour.

A separate literature develops concentration tools for dependent observations, including Marton-coupling bounded-difference inequalities \cite{Paulin2015} and transportation inequalities for dependent sequences \cite{DjelloutGuillinWu2004}. These provide general probabilistic machinery rather than model-specific generalisation guarantees. In the present setting, this leaves the statistical question of how dependence in the pretraining data affects transition-risk estimation, separately from the long-term behaviour of the recursively generated process.

\paragraph{Research questions.}
The thesis addresses four questions.
\begin{enumerate}[label=(Q\arabic*)]
    \item How can a functional autoregressive transformer be defined so that every context produces a valid conditional probability measure and the resulting shift-and-append update is a Borel Markov transition kernel?

    \item Which structural properties of the architecture are sufficient to establish existence and uniqueness of an invariant distribution, uniform geometric convergence in total variation, and finite invariant moments for recursive generation?

    \item For a fixed model parameter, how can the gap between empirical and population pretraining risk be controlled when the functional observations are dependent, and what explicit sample-complexity guarantees follow under zero empirical risk?

    \item How do coupling-based and metric-transport approaches differ in their assumptions and applicability across independent trajectories, overlapping windows, mixing chains, and continuous function-valued dynamics?
\end{enumerate}

\paragraph{Main contributions.}
The principal contributions are as follows.
\begin{enumerate}
    \item A functional autoregressive transformer is defined on a Hilbert space, with output heads lying in a fixed finite-dimensional approximation space, that produce basis coefficients specifying the conditional mean and covariance of the Gaussian next-function distribution.

    \item The conditional Gaussian measure is proved to exist uniquely for every context, to have common support, and to depend measurably on the context. Consequently, recursive generation is formulated rigorously as a time-homogeneous Markov chain on the full functional context space.

    \item Uniform conditional moments yield a geometric Foster--Lyapunov drift inequality, while the common approximation space and covariance floor yield a uniform lower bound for the output densities. After the original context has been shifted out, these bounds give a global multistep minorisation. A direct Doeblin argument then proves a unique invariant distribution, uniform geometric convergence in total variation from every initial distribution, support on the product approximation space, and a finite weighted second moment.

    \item For a fixed parameter value, a squared Wasserstein loss between reference and model transition kernels is introduced and shown to reduce exactly to the discrepancy between their conditional output laws. Under a uniform second-moment assumption, Marton coupling and dependent bounded differences give a high-probability generalisation bound and an explicit pretraining sample-complexity bound under zero empirical risk.

    \item A complementary generalisation theory is developed for metric-dependent functional trajectories. A path-space transportation inequality and Wasserstein stability of the reference and model kernels yield concentration for Lipschitz empirical losses and a second zero-risk sample-complexity bound. The architectural Lipschitz estimates required for this result are proved in the appendix, and the thesis distinguishes these pretraining assumptions from the ergodicity of the recursively generated inference chain.
\end{enumerate}

\paragraph{Organisation of the thesis.}
Section~\ref{sec:Preliminaries} introduces the mathematical tools used throughout and defines the functional autoregressive transformer. Section~\ref{sec:General state-space Markov chain framework for functional autoregressive transformers} establishes the Markov-chain formulation, Gaussian-kernel well-posedness, and the Chain's geometric ergodicity. Section~\ref{sec:Fixed-parameter pretraining generalisation} develops the transition loss, the Marton-coupling generalisation and sample-complexity results, the transportation-based metric theory, and the applicability of both approaches to different pretraining-data regimes. Section~\ref{sec:Conclusion and future work} summarises the results and outlines directions for extending the model and theory. Appendix~\ref{app:Lipschitz regularity} proves the global Lipschitz estimates used in the metric generalisation analysis.

\section{Preliminaries}\label{sec:Preliminaries}
This section collects the mathematical concepts required to define functional autoregressive transformers and to analyse the stochastic processes they induce. We first set notation and introduce the Galerkin-type approximation space \(S_d\). We then recall Gaussian measures on Hilbert spaces, including densities relative to finite-dimensional Lebesgue measure on \(S_d\). Next, we collect the general-state-space Markov-chain notions needed for drift, minorisation, and total-variation geometric ergodicity. The optimal-transport material is used for Gaussian comparison, the pretraining loss, metric dependence, and the Lipschitz regularity of the conditional output law. Finally, we define functional autoregressive models and the function-space transformer architecture studied throughout the thesis.

\subsection{Notation}\label{subsec:Notation}

We denote by \(\mathbb N=\{1,2,\ldots\}\) the set of positive integers and by \(\mathbb N_0=\{0,1,2,\ldots\}\) the set of non-negative integers. For \(N\in\mathbb N\), we write \([N]=\{1,\ldots,N\}\).

For vectors in \(\mathbb R^d\), \(\langle\cdot,\cdot\rangle_2\) and \(\|\cdot\|_2\) denote the Euclidean inner product and norm. The \(d\times d\) identity matrix is denoted by \(I_d\), and the transpose of a matrix \(A\) by \(A^\top\). For self-adjoint operators or symmetric matrices on the same space, \(A\preceq B\) means that \(B-A\) is non-negative.

Unless otherwise specified, let \(H\) denote a real separable Hilbert space with inner product \(\langle\cdot,\cdot\rangle_H\) and norm \(\|\cdot\|_H\). The zero element of \(H\) is denoted by \(0_H\), or simply \(0\) when no ambiguity can arise.

For a fixed context length \(n\in\mathbb N\), we denote the product space by \(E=H^n\). Elements of \(E\) are written using capital letters, for example \(F=(f_0,\ldots,f_{n-1})\), where each \(f_j\in H\). Unless otherwise specified, \(E\) is equipped with its product topology and corresponding Borel \(\sigma\)-algebra. Its zero element is denoted by \(0_E=(0_H,\ldots,0_H)\).

For a topological space \(X\), we denote its Borel \(\sigma\)-algebra by \(\mathcal B(X)\), and its set of Borel probability measures by \(\mathcal P(X)\). If \(X,d\) is a metric space, we denote the set of Borel probability measures with finite second moment by \(\mathcal P_2(X)=\{\mu\in\mathcal P(X):\int_X d(x,x_0)^2\,\mu(\mathrm dx)<\infty,\text{ for fixed }x_0\in X\}\). This definition does not depend on the choice of \(x_0\).

The Dirac probability measure concentrated at \(x\in X\) is denoted by \(\delta_x\). For a measurable set \(A\), its indicator function is denoted by \(\mathbf 1_A\). If \(T:X\to Y\) is a measurable map and \(\mu\in\mathcal P(X)\), the pushforward of \(\mu\) under \(T\) is denoted by \(T_\#\mu\), so that \((T_\#\mu)(A)=\mu(T^{-1}(A))\) for every \(A\in\mathcal B(Y)\).

For probability measures \(\mu,\nu\in\mathcal P(X)\), we write \(\mu\Rightarrow\nu\) to denote weak convergence. When \(X\) is equipped with a metric and \(\mu,\nu\in\mathcal P_2(X)\), the set of couplings of \(\mu\) and \(\nu\) is denoted by \(\Pi(\mu,\nu)\).

We denote by \(\mathcal L(H)\) the space of bounded linear operators from \(H\) to itself. For \(K\in\mathcal L(H)\), its adjoint is denoted by \(K^*\). The identity operator on \(H\) is denoted by \(I_H\). The space of self-adjoint, non-negative, trace-class operators on \(H\) is denoted by \(\mathcal L_1^+(H)\). For a trace-class operator \(K\), its trace norm is denoted by \(\|K\|_1\).

For normed spaces \(X\) and \(Y\), we denote by \(\mathcal L(X,Y)\) the space of bounded linear operators from \(X\) to \(Y\), and simplify the notation to \(\mathcal L(X)=\mathcal L(X,X)\). For \(T\in\mathcal L(X,Y)\), its operator norm is defined by \(\|T\|_{\mathrm{op}}=\sup_{\|x\|_X\leqslant1}\|Tx\|_Y=\sup_{x\neq0}\{{\|Tx\|_Y}/{\|x\|_X}\}. \) Consequently, \(\|Tx\|_Y\leqslant\|T\|_{\mathrm{op}}\|x\|_X\) for every \(x\in X\).

Random elements are denoted by capital letters when convenient, while their realisations are denoted by the corresponding lowercase letters. Probability measures are generally denoted by Greek letters such as \(\mu\), \(\nu\), and \(\pi\), while probability kernels are denoted by capital letters such as \(P\) and \(Q\). The notation ``a.s.'' stands for almost surely.



\subsection{Finite-dimensional approximation spaces}\label{subsec:Finite-dimensional approximation spaces}
The probabilistic output of the functional autoregressive transformer will be constructed in a fixed finite-dimensional subspace of a Hilbert space \(H\). Introducing such a space serves two purposes. First, it provides a finite-dimensional representation in which the conditional mean and covariance can be parameterised through coefficient vectors, making the output distribution computationally tractable while retaining a function-space interpretation. Second, requiring the same subspace to be used for every input context ensures that all conditional Gaussian measures are supported on a common space. This common support will be essential later when deriving a uniform Gaussian density lower bound and the minorisation condition used in the ergodicity analysis.

This construction is naturally interpreted through the Galerkin approximation framework, in which an infinite-dimensional problem is approximated on a sequence of finite-dimensional subspaces; see, for example, \cite{QuarteroniValli1994}. The definitions and identities used below follow from the standard theory of orthonormal systems and orthogonal projections in Hilbert spaces; see \cite[Chapter~5]{Brezis2011}.

Fix \(d\in\mathbb N\), \(d\leqslant\dim H\), and let \((e_1,\ldots,e_d)\subset H\) be an orthonormal family. We define the finite-dimensional approximation space with \(S_d=\operatorname{span}\{e_1,\ldots,e_d\}\subset H\).
Since \(S_d\) is finite-dimensional, it is a closed subspace of \(H\) and therefore admits a well-defined orthogonal projection. The dimension \(d\) controls the expressive resolution of the probabilistic output: increasing \(d\) enlarges the class of functions that can be represented, whereas components orthogonal to \(S_d\) cannot be produced by the model.

The choice of the orthonormal family is application dependent. For example, \(S_d\) may be generated by the first \(d\) elements of a Fourier or wavelet basis, by leading functional principal components, or by a globally learned collection of basis functions. The particular construction of the basis is not required for the analysis below. The essential structural assumption is that \(S_d\) is fixed once the model is specified and does not depend on the functional context \(F\). Consequently, the conditional mean and all covariance directions introduced later can be constrained to the same space \(S_d\), ensuring that every conditional Gaussian output has the same finite-dimensional support.

\begin{definition}[Synthesis and analysis operators]\label{def:Synthesis and analysis operators}
The \emph{synthesis operator} \(\mathsf E_d:\mathbb R^d\to S_d\subset H\) is defined by \(\mathsf E_d a=\sum_{j=1}^d a_j e_j\) for \(a=(a_1,\ldots,a_d)\in\mathbb R^d\). The corresponding \emph{analysis operator} \(\mathsf E_d^*:H\to\mathbb R^d\) is defined by \(\mathsf E_d^*f=(\langle f,e_1\rangle_H,\ldots,\langle f,e_d\rangle_H)\). We denote the orthogonal projection from \(H\) onto \(S_d\) by \(\Pi_{S_d}\).
\end{definition}

The synthesis operator identifies a vector of coefficients in \(\mathbb R^d\) with the corresponding function in \(S_d\), while the analysis operator extracts the coefficients of a function along the orthonormal family \((e_1,\ldots,e_d)\). The following proposition recalls the basic relations between these operators.

\begin{proposition}[Basic identities for the approximation space]\label{prop:Basic identities for the approximation space}
The synthesis and analysis operators satisfy \(\mathsf E_d^*\mathsf E_d=I_d\) and \(\mathsf E_d\mathsf E_d^*=\Pi_{S_d}\). Moreover, \(\mathsf E_d\) is an isometry, so that \(\|\mathsf E_d a\|_H=\|a\|_2\) for every \(a\in\mathbb R^d\), and \(\|\Pi_{S_d}\|_{\mathrm{op}}=1\).
\end{proposition}

\begin{proof}
We prove each statement in order. Let \(a=(a_1,\ldots,a_d)\in\mathbb R^d\). For each \(i\in\{1,\ldots,d\}\), orthonormality gives \((\mathsf E_d^*\mathsf E_d a)_i=\langle\sum_{j=1}^d a_j e_j,e_i\rangle_H=\sum_{j=1}^d a_j\langle e_j,e_i\rangle_H=a_i\). Hence \(\mathsf E_d^*\mathsf E_d a=a\) for every \(a\in\mathbb R^d\), and therefore \(\mathsf E_d^*\mathsf E_d=I_d\).

Now let \(f\in H\). By definition of the two operators, \(\mathsf E_d\mathsf E_d^*f=\sum_{j=1}^d\langle f,e_j\rangle_H e_j\). Since \((e_1,\ldots,e_d)\) is an orthonormal basis of \(S_d\), this is precisely the orthogonal projection of \(f\) onto \(S_d\). Thus \(\mathsf E_d\mathsf E_d^*=\Pi_{S_d}\).

To prove the isometry, \(\|\mathsf E_d a\|_H^2=\langle\sum_{i=1}^d a_i e_i,\sum_{j=1}^d a_j e_j\rangle_H=\sum_{i,j=1}^d a_i a_j\langle e_i,e_j\rangle_H=\sum_{j=1}^d a_j^2=\|a\|_2^2\), by orthogonality. Taking square roots yields \(\|\mathsf E_d a\|_H=\|a\|_2\).

Finally, since \(\Pi_{S_d}\) is an orthogonal projection, \(\|\Pi_{S_d}f\|_H\leqslant\|f\|_H\) for every \(f\in H\), so \(\|\Pi_{S_d}\|_{\mathrm{op}}\leqslant1\). On the other hand, for any nonzero \(f\in S_d\), \(\Pi_{S_d}f=f\), and therefore \(\|\Pi_{S_d}f\|_H/\|f\|_H=1\). Hence \(\|\Pi_{S_d}\|_{\mathrm{op}}=1\).
\end{proof}


\subsection{Gaussian random elements on Hilbert spaces}\label{subsec:Gaussian random elements on Hilbert spaces}

We recall the basic theory of Gaussian random elements on separable Hilbert spaces that will be used to define the probabilistic output of the functional autoregressive transformer in Section \ref{subsec:Functional autoregressive transformer}. For a comprehensive treatment, we refer to \cite{Bogachev1998}.
Let \(H\) be a real separable Hilbert space with inner product \(\langle\cdot,\cdot\rangle_H\) and norm \(\|\cdot\|_H\).

\begin{definition}[Gaussian random element]
An \(H\)-valued random element \(X\) is called \emph{Gaussian} if, for every \(h\in H\), the real-valued random variable \(\langle X,h\rangle_H\) is Gaussian. If \(X\) is Gaussian, then its law is called a \emph{Gaussian measure} on \(H\).
\end{definition}

The definition can be understood as an infinite-dimensional extension of the usual notion of a Gaussian random vector. Rather than requiring the random element itself to possess a density, which is generally not meaningful in infinite dimensions, we require every one-dimensional projection of the random element to be Gaussian.

\begin{example}[]
Let \(m,c\in H\) and let \(\xi\sim\mathcal N(0,1)\). Define \(X=m+\xi c\). Then \(X\) is an \(H\)-valued Gaussian random element. Indeed, for every \(h\in H\), \(\langle X,h\rangle_H=\langle m,h\rangle_H+\xi\langle c,h\rangle_H\), which is a real Gaussian random variable.

Although \(X\) takes values in the infinite-dimensional space \(H\), all of its randomness lies in the single direction \(c\). This example will later generalise to the finite-rank Gaussian construction used by the functional autoregressive transformer.
\end{example}

\begin{theorem}[Finite moments of Gaussian random elements]\label{thm:Finite moments of Gaussian random elements}
Let \(H\) be a real separable Hilbert space and let \(X\) be an \(H\)-valued Gaussian random element. Then, for every \(p\geqslant1\), \(\mathbb E\bigl[\|X\|_H^p\bigr]<\infty\).
\end{theorem}

\begin{proof}
By Fernique theorem \parencite[Theorem 2.8.5.]{Bogachev1998}, there exists \(\alpha>0\) such that \(\mathbb E\left[\exp\left(\alpha\|X\|_H^2\right)\right]<\infty\). Fix \(p\geqslant1\). Since exponential growth dominates polynomial growth, there exists a constant \(C_{p,\alpha}<\infty\) such that \(x^p\leqslant C_{p,\alpha}\exp(\alpha x^2)\) for every \(x\geqslant0\). Applying this inequality with \(x=\|X\|_H\) and taking expectations gives \(\mathbb E\bigl[\|X\|_H^p\bigr]\leqslant C_{p,\alpha}\mathbb E\left[\exp\left(\alpha\|X\|_H^2\right)\right]<\infty\). Hence every polynomial moment of \(\|X\|_H\) is finite.
\end{proof}

The preceding theorem guarantees the integrability required to define the mean and covariance of an \(H\)-valued Gaussian random element. We may therefore introduce the following definitions.

\begin{definition}[Mean and covariance operator]
Let \(X\) be a Gaussian random element. The \emph{mean} of \(X\) is the unique element \(m\in H\), denoted \(m=\mathbb E[X]\), such that, for every \(h\in H\), \(\langle m,h\rangle_H=\mathbb E[\langle X,h\rangle_H]\). The \emph{covariance operator} of \(X\) is the unique bounded linear operator \(C:H\to H\) characterised by \(\langle Ch,k\rangle_H=\mathbb E[\langle X-m,h\rangle_H\langle X-m,k\rangle_H]\) for every \(h,k\in H\).
\end{definition}

The mean and covariance operator are therefore the infinite-dimensional analogue of the mean vector and covariance matrix of a Gaussian random variable.

\begin{definition}[Rank-one operator]
For \(u,v\in H\), the rank-one operator \(u\otimes v:H\to H\) is defined by \(
(u\otimes v)h=\langle v,h\rangle_Hu\), for any \(h\in H\).
\end{definition}

\begin{example}[]
Consider as in the previous example the Gaussian random element \(X=m+\xi c\), where \(m,c\in H\) and \(\xi\sim\mathcal N(0,1)\). Since \(\mathbb E[\xi]=0\), the mean of \(X\) is \(m\). To identify its covariance operator, let \(h,k\in H\). Then
\[
\langle Ch,k\rangle_H=\mathbb E\bigl[\langle X-m,h\rangle_H\langle X-m,k\rangle_H\bigr]=\mathbb E[\xi^2]\langle c,h\rangle_H\langle c,k\rangle_H=\langle c,h\rangle_H\langle c,k\rangle_H = \langle (c\otimes c)h,k\rangle_H.
\]
Hence, the covariance operator of \(X\) is \(C=c\otimes c\). Thus, the function \(c\) determines the direction of variation, while its norm controls the magnitude of that variation.
\end{example}

\begin{definition}[Self-adjoint, non-negative, trace-class operators]
A bounded linear operator \(C:H\to H\) is called (i) \emph{self-adjoint} if \(\langle Ch,k\rangle_H=\langle h,Ck\rangle_H\) for every \(h,k\in H\), and (ii) \emph{non-negative} if \(\langle Ch,h\rangle_H\geqslant0\) for every \(h\in H\). A non-negative self-adjoint operator \(C\) is called (iii) \emph{trace class} if
\[
\operatorname{tr}(C)=\sum_{j=1}^{\infty}\langle Ce_j,e_j\rangle_H<\infty,
\]
where \((e_j)_{j\geqslant1}\) is any orthonormal basis of \(H\). We denote the collection of self-adjoint, non-negative, trace-class bounded linear operators on \(H\) by \(\mathcal L_1^+(H)\).
\end{definition}

The following result characterises the covariance operators of Gaussian measures on a separable Hilbert space.

\begin{theorem}[Covariance operator of a Gaussian random element]\label{thm:Covariance operator of a Gaussian random element}
Let \(X\) be an \(H\)-valued Gaussian random element with mean \(m\) and covariance operator \(C\). Then \(C\in\mathcal L_1^+(H)\) and \(\operatorname{tr}(C)=\mathbb E\bigl[\|X-m\|_H^2\bigr]\). Consequently, \(\mathbb E\bigl[\|X\|_H^2\bigr]=\|m\|_H^2+\operatorname{tr}(C)
\).
\end{theorem}

\begin{proof}
Let \(Y=X-m\). Then \(Y\) is a centred Gaussian random element with covariance operator \(C\). We first prove the covariance properties for the centred variable \(Y\), and then recover the final second-moment identity for \(X\).

We first show that \(C\) is self-adjoint by computing directly for \(h,k\in H\):
\[
\langle Ch,k\rangle_H=\mathbb E[\langle Y,h\rangle_H\langle Y,k\rangle_H]=\mathbb E[\langle Y,k\rangle_H\langle Y,h\rangle_H]=\langle Ck,h\rangle_H=\langle h,Ck\rangle_H.
\]
Next, \(C\) is non-negative since \(\langle Ch,h\rangle_H=\mathbb E[\langle Y,h\rangle_H^2]\geqslant0\). \newline
It remains to prove that \(C\) is trace class. Let \((e_j)_{j\geqslant1}\) be an orthonormal basis of \(H\). Then,
\[\operatorname{tr}(C)=\sum_{j=1}^{\infty}\langle Ce_j,e_j\rangle_H=\sum_{j=1}^{\infty}\mathbb E[\langle Y,e_j\rangle_H^2]=\mathbb E\left[\sum_{j=1}^{\infty}\langle Y,e_j\rangle_H^2\right]=\mathbb E[\|Y\|_H^2]=\mathbb E[\|X-m\|_H^2],
\]
where the first equality is the definition of the trace for a non-negative operator, the second follows from the definition of the covariance operator, the third follows from Tonelli's theorem since all terms are non-negative, the fourth follows from Parseval's identity, and the final equality uses \(Y=X-m\).

Finally, since \(X=m+Y\), we have \(\|X\|_H^2=\|m\|_H^2+2\langle m,Y\rangle_H+\|Y\|_H^2\). Taking expectations gives \(\mathbb E[\|X\|_H^2]=\|m\|_H^2+2\mathbb E[\langle m,Y\rangle_H]+\mathbb E[\|Y\|_H^2]\). Since \(Y\) is centred, \(\mathbb E[Y]=0\), and therefore \(\mathbb E[\langle m,Y\rangle_H]=0\). Using the trace identity above, we conclude that \(\mathbb E[\|X\|_H^2]=\|m\|_H^2+\operatorname{tr}(C)\).
\end{proof}

Thus, in infinite dimensions, not every bounded non-negative operator can serve as the covariance of an \(H\)-valued Gaussian random element. The trace-class condition ensures that the total variance of the random element is finite.

The converse is also true and provides the fundamental existence result used throughout this thesis. For a proof, we refer to Bogachev \parencite[Theorem 2.3.1.]{Bogachev1998}

\begin{theorem}[Existence and uniqueness of Gaussian random elements]\label{thm:Existence and uniqueness of Gaussian random elements}
Let \(m\in H\) and \(C\in\mathcal L_1^+(H)\). Then there exists an \(H\)-valued Gaussian random element \(X\) with mean \(m\) and covariance operator \(C\). Moreover, the distribution of \(X\) is uniquely determined by \(m\) and \(C\). In other words, if \(X_1\) and \(X_2\) are \(H\)-valued Gaussian random elements with the same mean \(m\) and covariance operator \(C\), then \(X_1\) and \(X_2\) have the same law. We denote this common law by \(\mathcal N(m,C)\).
\end{theorem}

\subsection{Gaussian measures supported on the approximation space}
\label{subsec:Gaussian measures supported on the approximation space}

The Gaussian-measure results above apply to probability measures on the full Hilbert space \(H\). The output distributions considered later have additional structure: their randomness is confined to the fixed finite-dimensional approximation space \(S_d\). It is therefore useful to recall explicitly how the general Hilbert-space Gaussian theory reduces to the familiar finite-dimensional Gaussian theory on \(S_d\).

This distinction will be important in the ergodicity analysis. Although the covariance operators used by the functional autoregressive transformer are degenerate when regarded as operators on the full infinite-dimensional space \(H\), their restrictions to \(S_d\) are uniformly positive definite. Consequently, the corresponding Gaussian measures admit ordinary finite-dimensional densities with respect to Lebesgue measure on \(S_d\). Uniform control of these densities will later provide the common minorisation measure for the autoregressive transition kernel.

Using the synthesis operator \(\mathsf E_d:\mathbb R^d\to S_d\), we identify \(S_d\) isometrically with \(\mathbb R^d\). Let \(\lambda_d\) denote Lebesgue measure on \(\mathbb R^d\), and define the corresponding \(d\)-dimensional Lebesgue measure on \(S_d\) by
\(
\lambda_{S_d}=(\mathsf E_d)_\#\lambda_d.
\)
This is a Lebesgue measure on the finite-dimensional space \(S_d\); no infinite-dimensional Lebesgue measure on \(H\) is being introduced.

\begin{proposition}[Gaussian density on the approximation space]
\label{prop:Gaussian density on the approximation space}
Let \(m\in S_d\), and let \(C\in\mathcal L_1^+(H)\) satisfy \(\operatorname{Ran}(C)\subseteq S_d\). Suppose that the restriction \(C|_{S_d}:S_d\to S_d\) is positive definite. Then the Gaussian measure \(\mathcal N(m,C)\) is supported on \(S_d\) and is absolutely continuous with respect to \(\lambda_{S_d}\). Its density is
\[
q_{m,C}(y)
=
(2\pi)^{-d/2}
\textstyle
\det_{S_d}(C)^{-1/2}
\exp\left(
-\frac12
\left\langle
(C|_{S_d})^{-1}(y-m),y-m
\right\rangle_H
\right),
\qquad y\in S_d,
\]
where \(\det_{S_d}(C)\) is the determinant of the matrix representing \(C|_{S_d}\) in the orthonormal basis \((e_1,\ldots,e_d)\).
\end{proposition}

\begin{proof}
Write \(m=\mathsf E_d\mu\) for some \(\mu\in\mathbb R^d\), and define
\(
\Sigma=\mathsf E_d^*C\mathsf E_d\in\mathbb R^{d\times d}.
\)
Since \(C|_{S_d}\) is positive definite and \(\mathsf E_d\) is an isometric isomorphism from \(\mathbb R^d\) onto \(S_d\), the matrix \(\Sigma\) is symmetric and positive definite.

If \(Z\sim\mathcal N(\mu,\Sigma)\) in \(\mathbb R^d\), then \(\mathsf E_dZ\) is an \(H\)-valued Gaussian random element with mean \(m\), covariance \(C\), and values in \(S_d\). Hence, \(
\mathcal N(m,C)
=
(\mathsf E_d)_\#
\mathcal N(\mu,\Sigma).
\)
The Euclidean Gaussian \(\mathcal N(\mu,\Sigma)\) has the usual density with respect to \(\lambda_d\). Transporting this density through the isometry \(\mathsf E_d\), together with the definition \(\lambda_{S_d}=(\mathsf E_d)_\#\lambda_d\), gives the result.
\end{proof}

The density is an ordinary finite-dimensional Gaussian density on \(S_d\). In particular, no claim is made that \(\mathcal N(m,C)\) possesses a Lebesgue density on the ambient infinite-dimensional space \(H\).

The next result extracts the property that will be needed in the minorisation argument. If the means remain in a fixed bounded set and the covariance operators are uniformly bounded above and below on \(S_d\), then all of the corresponding Gaussian densities are bounded uniformly away from zero on any fixed bounded subset of \(S_d\).

\begin{lemma}[]
\label{lem:Uniform lower bound for Gaussian densities on a fixed ball}
Let \(M,R_0<\infty\) and \(0<\underline\lambda\leqslant\overline\lambda<\infty\). Suppose that \(m\in S_d\) and \(C\in\mathcal L_1^+(H)\) satisfy
\(
\|m\|_H\leqslant M,
\)
\(
\operatorname{Ran}(C)\subseteq S_d,
\)
and
\(
\underline\lambda I_{S_d}
\preceq
C|_{S_d}
\preceq
\overline\lambda I_{S_d}.
\)
Then, for every \(y\in B_{S_d}(0,R_0)\),
\[
q_{m,C}(y)
\geqslant
(2\pi)^{-d/2}
\overline\lambda^{-d/2}
\exp\left(
-\frac{(R_0+M)^2}{2\underline\lambda}
\right)
>0.
\]
\end{lemma}

\begin{proof}
Let \(\lambda_1,\ldots,\lambda_d\) denote the eigenvalues of \(C|_{S_d}\). The upper spectral bound gives \(\lambda_j\leqslant\overline\lambda\) for every \(j\), and therefore
\(
\det_{S_d}(C)
=
\prod_{j=1}^d\lambda_j
\leqslant
\overline\lambda^d.
\)
Consequently, \(\det_{S_d}(C)^{-1/2}\geqslant\overline\lambda^{-d/2}\).

The lower spectral bound gives
\(
(C|_{S_d})^{-1}
\preceq
\underline\lambda^{-1}I_{S_d}.
\)
Hence, for \(y\in B_{S_d}(0,R_0)\),
\[
\begin{aligned}
\left\langle
(C|_{S_d})^{-1}(y-m),y-m
\right\rangle_H
&\leqslant
\underline\lambda^{-1}\|y-m\|_H^2\\
&\leqslant
\underline\lambda^{-1}
\bigl(\|y\|_H+\|m\|_H\bigr)^2\\
&\leqslant
\underline\lambda^{-1}(R_0+M)^2.
\end{aligned}
\]
Substituting the determinant and quadratic-form estimates into the density formula of Proposition~\ref{prop:Gaussian density on the approximation space} gives the stated lower bound.
\end{proof}


\subsection{Finite-rank Gaussian covariance construction}\label{subsec:Finite-rank Gaussian covariance construction}

We now introduce the finite-rank Gaussian covariance construction that will later be used in constructing the functional autoregressive transformer in Section \ref{subsec:Functional autoregressive transformer}.

\begin{definition}[Finite-rank operator]
A bounded linear operator \(C:H\to H\) is called \emph{finite rank} if it can be written as a finite sum of rank-one operators. In other words \(C\) is finite rank if we can write \(C=\sum_{k=1}^r u_k\otimes v_k\), for some \(u_1,\ldots,u_r,v_1,\ldots,v_r\in H\).
\end{definition}

Given functions \(c_1,\ldots,c_r\in H\), we define the operator \(C=\sum_{k=1}^r c_k\otimes c_k\). This construction is motivated by the need to parameterise covariance operators in a way that is both mathematically valid and computationally tractable. In infinite-dimensional spaces, a covariance operator must be in \(\mathcal L_1^+\). Rather than asking a model to output an arbitrary bounded operator and then imposing these constraints afterwards, we construct the covariance operator directly from a finite collection of functions.

The functions \(c_1,\ldots,c_r\) represent directions of variability in the random function. The resulting operator assigns variance only in the finite-dimensional subspace spanned by these directions. The following proposition verifies that this finite-rank construction automatically produces an admissible covariance operator.

\begin{proposition}[Finite-rank covariance]\label{prop:Finite-rank covariance}
Let \(c_1,\ldots,c_r\in H\), and define the finite-rank covariance by \(C=\sum_{k=1}^r c_k\otimes c_k\). Then \(C\in\mathcal L_1^+(H)\).
\end{proposition}

\begin{proof}
First, \(C\) is self-adjoint. For \(h,k\in H\), \(\langle Ch,k\rangle_H
=
\sum_{j=1}^r\langle h,c_j\rangle_H\langle c_j,k\rangle_H
=
\langle h,Ck\rangle_H\). Second, \(C\) is non-negative. For every \(h\in H\), \(\langle Ch,h\rangle_H=\sum_{j=1}^r|\langle h,c_j\rangle_H|^2\geqslant0\). Finally, \(C\) has finite rank because \(\operatorname{Ran}(C)\subseteq\operatorname{span}\{c_1,\ldots,c_r\}\). Every finite-rank operator is trace class. Therefore, \(C\in\mathcal L_1^+(H)\).
\end{proof}

Since this finite-rank covariance operator is admissible, we now introduce a factorisation that will be useful for sampling and for the later functional autoregressive transformer.

\begin{definition}[Covariance factor]
Let \(c_1,\ldots,c_r\in H\). We call a \emph{covariance factor} an operator \(B:\mathbb R^r\to H\) such that \(Ba=\sum_{k=1}^r a_kc_k\), where \(a=(a_1,\ldots,a_r)\in\mathbb R^r\).
\end{definition}

The operator \(B\) maps a vector of coefficients in \(\mathbb R^r\) to the corresponding linear combination of the functions \(c_1,\ldots,c_r\).

\begin{lemma}[Adjoint of the covariance factor]
Let \(B\) be a covariance factor. Then its adjoint is the operator \(B^*:H\to\mathbb R^r\) given by \(B^*h=\bigl(\langle h,c_1\rangle_H,\ldots,\langle h,c_r\rangle_H\bigr)\).
\end{lemma}

\begin{proof}
Let \(a=(a_1,\ldots,a_r)\in\mathbb R^r\) and \(h\in H\). Then
\[
\langle Ba,h\rangle_H=\left\langle\sum_{k=1}^ra_kc_k,h\right\rangle_H=\sum_{k=1}^ra_k\langle c_k,h\rangle_H=\left\langle a,\bigl(\langle h,c_1\rangle_H,\ldots,\langle h,c_r\rangle_H\bigr)\right\rangle_{\mathbb R^r}.
\]
The first equality uses the definition of \(B\), the second uses linearity of the inner product, and the third uses the symmetry of the real Hilbert-space inner product. By the definition of the adjoint, the vector in \(\mathbb R^r\) appearing on the right-hand side is \(B^*h\).
\end{proof}

\begin{proposition}[Finite-rank covariance factorisation]\label{prop:Finite-rank covariance factorisation}
Let \(c_1,\ldots,c_r\in H\) and \(B\) be a covariance factor. Then the finite-rank covariance satisfies \(C=BB^*\).
\end{proposition}

\begin{proof}
For every \(h\in H\),
\[
BB^*h=B\bigl(\langle h,c_1\rangle_H,\ldots,\langle h,c_r\rangle_H\bigr)=\sum_{k=1}^r\langle h,c_k\rangle_Hc_k=\left(\sum_{k=1}^r c_k\otimes c_k\right)h=Ch.
\]
The first equality uses the definition of the adjoint \(B^*\), the second uses the definition of \(B\), the third uses the definition of the rank-one operator \(c_k\otimes c_k\), and the final equality uses the definition of \(C\). Since this holds for every \(h\in H\), we conclude that \(C=BB^*\).
\end{proof}

The following result gives the probabilistic interpretation of this factorisation.

\begin{theorem}[Finite-rank Gaussian construction]\label{thm:Finite-rank Gaussian construction}
Let \(m\in H\), let \(c_1,\ldots,c_r\in H\), and define as above \(B:\mathbb R^r\to H\) and \(C:H\to H\). Let \(\xi=(\xi_1,\ldots,\xi_r)\sim\mathcal N(0,I_r)\), where \(\xi_1,\ldots,\xi_r\) are independent standard Gaussian random variables. Define the \(H\)-valued random element \(X=m+B\xi\). Then \(X\) is Gaussian with mean \(m\) and covariance operator \(C=BB^*\)
\end{theorem}

\begin{proof}
Let \(h\in H\). Since \(X=m+B\xi\), we have
\(\langle X,h\rangle_H=\langle m,h\rangle_H+\langle B\xi,h\rangle_H\). Using the definition of the adjoint, \(\langle B\xi,h\rangle_H=\langle \xi,B^*h\rangle_{\mathbb R^r}\). Since \(\xi\sim\mathcal N(0,I_r)\), the random variable \(\langle \xi,B^*h\rangle_{\mathbb R^r}\) is a real Gaussian random variable. Therefore, for every \(h\in H\), \(\langle X,h\rangle_H\) is Gaussian, and hence \(X\) is an \(H\)-valued Gaussian random element.

Since \(\mathbb E[\xi]=0\), the mean is easily computed \(\mathbb E[X]=m+\mathbb E[B\xi]=m+B\mathbb E[\xi]=m\). It remains to identify the covariance operator, but first a technical lemma.
\begin{lemma}\label{lem:Covariance identity for a standard Gaussian vector}
Let \(\xi\sim\mathcal N(0,I_r)\) and let \(a,b\in\mathbb R^r\). Then \(\mathbb E\bigl[\langle \xi,a\rangle_{\mathbb R^r}\langle \xi,b\rangle_{\mathbb R^r}\bigr]
=
\langle a,b\rangle_{\mathbb R^r}
\).
\end{lemma}

\begin{proof}
Writing \(a=(a_1,\ldots,a_r)\), \(b=(b_1,\ldots,b_r)\), and \(\xi=(\xi_1,\ldots,\xi_r)\), we have
\[
\mathbb E\bigl[\langle \xi,a\rangle_{\mathbb R^r}\langle \xi,b\rangle_{\mathbb R^r}\bigr]
=
\mathbb E\left[\sum_{i=1}^r\sum_{j=1}^r\xi_i\xi_j a_i b_j\right]
=
\sum_{i=1}^r\sum_{j=1}^r\mathbb E[\xi_i\xi_j]a_ib_j
=
\sum_{i=1}^ra_ib_i
=
\langle a,b\rangle_{\mathbb R^r}.
\]
The third equality uses that the components of \(\xi\) are independent standard Gaussian random variables, so \(\mathbb E[\xi_i\xi_j]=\delta_{ij}\).
\end{proof}

Let \(h,k\in H\) and let \(\mathscr C:H \to H\) defined by \(\langle \mathscr Ch,k\rangle = \mathbb E[\langle X-m,h\rangle_H\langle X-m,k\rangle_H]\) be the covariance operator. Then
\[
\langle \mathscr Ch,k\rangle =\mathbb E[\langle B\xi,h\rangle_H\langle B\xi,k\rangle_H]=\mathbb E[\langle \xi,B^*h\rangle_{\mathbb R^r}\langle \xi,B^*k\rangle_{\mathbb R^r}]=\langle B^*h,B^*k\rangle_{\mathbb R^r}=\langle BB^*h,k\rangle_H,
\]
where the first equality follows from \(X-m=B\xi\), the second equality follows from the definition of the adjoint operator, the third equality uses Lemma \ref{lem:Covariance identity for a standard Gaussian vector} with \(a=B^*h\) and \(b=B^*k\), and the final equality follows from the definition of the adjoint. Hence the covariance operator \(\mathscr C = BB^* = C\), therefore \(X\sim\mathcal N(m,C)\).
\end{proof}

The finite-rank factorisation has useful consequences, such as \(C\) being \(\mathcal L_1^+(H)\), and \(\operatorname{tr}(C)=\sum_{k=1}^r\|c_k\|_H^2\). Furthermore, this construction will be used later to parameterise the covariance output of the functional autoregressive transformer. Rather than learning an arbitrary operator on \(H\), the model learns a finite collection of functions \(c_1,\ldots,c_r\). These functions determine both the covariance operator and the random perturbation used to generate the next function.


\subsection{General state-space Markov chains}\label{subsec:General state-space Markov chains}

We recall the general state-space framework for Markov chains, following Meyn and Tweedie \cite{MeynTweedie2009}. In this thesis, the state space will be the product Hilbert space \(E=H^n\), where \(H\) is a separable Hilbert space. We equip \(E\) with its Borel \(\sigma\)-algebra \(\mathcal B(E)\).
The starting point for the construction of a Markov chain on \(E\) is a collection of one-step transition probabilities.

\begin{definition}[Markov transition kernel]\label{def:Markov transition kernel}
Let \(E\) be a measurable state space equipped with its Borel \(\sigma\)-algebra \(\mathcal B(E)\). A map \(P:E\times\mathcal B(E)\to[0,1]\) is called a \emph{Markov transition kernel} if: (i) for every fixed \(F\in E\), the map \(A\mapsto P(F,A)\) is a probability measure on \((E,\mathcal B(E))\), and (ii) for every fixed \(A\in\mathcal B(E)\), the map \(F\mapsto P(F,A)\) is Borel measurable.
\end{definition}
For a fixed state \(F\in E\), the probability measure \(P(F,\cdot)\) describes the distribution of the state after one transition from \(F\). The measurability condition ensures that these transition probabilities depend measurably on the current state.

\begin{example}[Deterministic transition kernel]
Let \(T:E\to E\) be a Borel measurable map and define \(P(F,A)=\mathbf 1_A(T(F))\). For each \(F\in E\), \(P(F,\cdot)\) is the Dirac probability measure concentrated at \(T(F)\). Moreover, for every \(A\in\mathcal B(E)\), \(F\mapsto P(F,A)=\mathbf 1_A(T(F))\) is Borel measurable. Hence \(P\) is a Markov transition kernel. This example shows that a Markov kernel need not describe a genuinely random transition: deterministic dynamical systems are included as a special case.
\end{example}

\begin{example}[Gaussian transition kernel]
Suppose that \(m:E\to H\) and \(C:E\to\mathcal L_1^+(H)\) are suitably measurable. A kernel of the form \(Q(F,\cdot)=\mathcal N(m(F),C(F))\) describes a transition whose next output is a Gaussian random function with context-dependent mean and covariance. The functional autoregressive transformer studied later will produce a kernel of precisely this form.
\end{example}

The following result, see \cite[Theorem 3.4.1.]{MeynTweedie2009}, shows that a transition probability kernel, together with an initial probability measure, determines a stochastic process.

\begin{theorem}[Existence of the Markov process]\label{thm:Existence of the Markov process}
Let \(\mu\) be a probability measure on \((E,\mathcal B(E))\), and let \(P\) be a transition probability kernel on \(E\). Then there exists a stochastic process \((F_t)_{t\in\mathbb N_0}\) with state space \(E\) and initial distribution \(\mu\) such that, for every \(t\in\mathbb N_0\) and every \(A_0,\ldots,A_t\in\mathcal B(E)\),
\begin{equation}\label{eq:Markov process}
\mathbb P_\mu(F_0\in A_0,\ldots,F_t\in A_t) = \int_{A_0}\mu(\mathrm dx_0) \int_{A_1}P(x_0,\mathrm dx_1) \cdots \int_{A_t}P(x_{t-1},\mathrm dx_t).
\end{equation}
\end{theorem}

\begin{definition}[Time-homogeneous Markov chain]
A stochastic process \((F_t)_{t\in\mathbb N_0}\) taking values in \(E\) is called a \emph{time-homogeneous Markov chain} with transition probability kernel \(P\) and initial distribution \(\mu\) if its finite-dimensional distributions satisfy \eqref{eq:Markov process}, for every \(t\in\mathbb N_0\) and \(A_0,\ldots,A_t\in\mathcal B(E)\).

Equivalently, the process satisfies the Markov property \(\mathbb P_\mu(F_{t+1}\in A\mid F_0,\ldots,F_t)=P(F_t,A)\) almost surely, for every \(t\in\mathbb N_0\) and \(A\in\mathcal B(E)\). Thus, conditionally on the current state, the distribution of the next state is independent of the preceding history.
\end{definition}

The one-step transition kernel describes only a single update of the context. To study the long-time behaviour of the Markov chain, we also need notation for repeated transitions and for the evolution of an arbitrary initial distribution under the kernel. This leads naturally to iterated kernels and to the notion of an invariant probability measure. These concepts will be used later to formulate convergence of the functional autoregressive transformer inference chain towards equilibrium.

\begin{definition}[Iterated kernels and action on measures]\label{def:Iterated kernels and action on measures}
For a Markov kernel \(P\), we define the initial kernel \(P^0(F,A)=\mathbf 1_A(F)\) and recursively
\(
P^{t+1}(F,A)=\int_E P^t(G,A)P(F,\mathrm dG).
\)
For \(\mu\in\mathcal P(E)\), the measure \(\mu P\) is defined by \((\mu P)(A)=\int_E P(F,A)\mu(\mathrm dF)\). A probability measure \(\pi\) is \emph{invariant} if \(\pi P=\pi\).
\end{definition}

To quantify convergence of these evolving probability measures, we require a metric on probability distributions. Total variation measures the largest difference that two probability measures assign to the same measurable event.

\begin{definition}[Total variation distance]\label{def:Total variation distance}
For \(\mu,\nu\in\mathcal P(E)\),
\(
d_{\mathrm{TV}}(\mu,\nu)=\sup_{A\in\mathcal B(E)}|\mu(A)-\nu(A)|.
\)
Note that \(0\leqslant d_{\mathrm{TV}}(\mu,\nu)\leqslant1\).
\end{definition}

\subsection{Drift, small sets, and geometric ergodicity}\label{subsec:Drift small sets and geometric ergodicity}

Having introduced invariant measures and total variation distance, we now turn to conditions that can be used to establish convergence towards equilibrium. In a general state space, the existence of an invariant probability measure and convergence to it do not follow automatically from the definition of a Markov kernel. Additional structure is required to control both the behaviour of the chain far from equilibrium and its ability to mix once it enters a suitable region of the state space.

Two standard mechanisms are particularly useful for this purpose. A \emph{Lyapunov drift condition} controls the growth of the chain by showing that, on average, a suitable non-negative function decreases outside a bounded region. A \emph{minorisation condition} provides a common probabilistic component in the transition laws and therefore quantifies the ability of the chain to forget its current state. Together, drift and minorisation form a classical route to geometric ergodicity for general-state-space Markov chains; see, for example, \cite{MeynTweedie2009}.

For the functional autoregressive transformer considered later, these two mechanisms arise naturally from the architecture. The deterministic shift of the context window, together with a uniform bound on the second moment of the newly generated function, will yield a geometric drift inequality. The fixed approximation space and the nondegenerate Gaussian covariance floor will provide a common lower bound for the conditional output densities, which will in turn lead to a multistep minorisation of the context kernel. We therefore introduce the corresponding notions below.

\begin{definition}[Lyapunov function and drift]\label{def:Lyapunov function and drift}
Any measurable map \(V:E\to[1,\infty)\) is called a \emph{Lyapunov function} and we write
\(
PV(F)=\int_EV(G)P(F,\mathrm dG).
\)
A geometric Foster--Lyapunov drift condition is an inequality of the form
\(
PV\leqslant\lambda V+b\mathbf 1_C
\)
for some \(\lambda\in(0,1)\), \(b<\infty\), and measurable set \(C\).
\end{definition}

\begin{definition}[Small set and minorisation]\label{def:Small set and minorisation}
A measurable set \(C\subseteq E\) is called \((m,\varepsilon,\nu)\)-small if there exist \(m\in\mathbb N\), \(\varepsilon>0\), and \(\nu\in\mathcal P(E)\) such that
\(
P^m(F,A)\geqslant\varepsilon\nu(A)
\)
for every \(F\in C\) and \(A\in\mathcal B(E)\). If \(C=E\), this is a global \(m\)-step Doeblin condition.
\end{definition}

\begin{definition}[Uniform geometric ergodicity]\label{def:Uniform geometric ergodicity}
A Markov kernel \(P\) is called \emph{uniformly geometrically ergodic in total variation} if it admits an invariant probability measure \(\pi\) and there exist \(M_{\mathrm{TV}}<\infty\) and \(q_{\mathrm{TV}}\in(0,1)\) such that
\(
\sup_{F\in E}d_{\mathrm{TV}}\bigl(P^t(F,\cdot),\pi\bigr)\leqslant M_{\mathrm{TV}}q_{\mathrm{TV}}^t,\)
for all \(t\in\mathbb N_0.
\)
\end{definition}

The equivalence between uniform ergodicity and a global smallness condition is discussed in Meyn and Tweedie \cite[Theorem~16.0.2]{MeynTweedie2009}. For the functional autoregressive transformer, we will only require the following elementary consequence of a global Doeblin condition. It shows that whenever every transition law contains the same probabilistic component, the Markov kernel is contractive in total variation.

\begin{lemma}[Doeblin contraction in total variation]\label{lem:Doeblin contraction in total variation}
Let \(P\) be a Markov kernel on \(E\). Suppose there exist \(\varepsilon\in(0,1]\) and \(\nu\in\mathcal P(E)\) such that \(P(F,A)\geqslant\varepsilon\nu(A)\) for every \(F\in E\) and \(A\in\mathcal B(E)\). Then, for every \(\mu_1,\mu_2\in\mathcal P(E)\), \(d_{\mathrm{TV}}(\mu_1P,\mu_2P)\leqslant(1-\varepsilon)d_{\mathrm{TV}}(\mu_1,\mu_2)\).
\end{lemma}

\begin{proof}
If \(\varepsilon=1\), then \(P(F,A)\geqslant\nu(A)\) for every \(A\in\mathcal B(E)\). Since both \(P(F,\cdot)\) and \(\nu\) are probability measures, equality of their total masses implies \(P(F,\cdot)=\nu(\cdot)\) for every \(F\in E\). Hence \(\mu_1P=\mu_2P=\nu\), and the result follows immediately.

Suppose now that \(\varepsilon<1\). Define \(R(F,A)=[P(F,A)-\varepsilon\nu(A)]/(1-\varepsilon)\) for \(F\in E\) and \(A\in\mathcal B(E)\). The minorisation assumption gives \(R(F,A)\geqslant0\), while \(R(F,E)=[P(F,E)-\varepsilon\nu(E)]/(1-\varepsilon)=1\). Moreover, \(F\mapsto R(F,A)\) is measurable for every \(A\in\mathcal B(E)\), so \(R\) is a Markov kernel. Then, by construction \(P(F,\cdot)=\varepsilon\nu(\cdot)+(1-\varepsilon)R(F,\cdot)\).

Therefore, for \(i\in\{1,2\}\), \(\mu_iP=\varepsilon\nu+(1-\varepsilon)\mu_iR\), and hence \(\mu_1P-\mu_2P=(1-\varepsilon)(\mu_1R-\mu_2R)\). It follows that \(d_{\mathrm{TV}}(\mu_1P,\mu_2P)=(1-\varepsilon)d_{\mathrm{TV}}(\mu_1R,\mu_2R)\). Since Markov kernels are non-expansive in total variation, \(d_{\mathrm{TV}}(\mu_1R,\mu_2R)\leqslant d_{\mathrm{TV}}(\mu_1,\mu_2)\). Combining these two inequalities concludes the proof.
\end{proof}

\begin{remark}
If \(P\) satisfies a global \(m\)-step Doeblin condition \(P^m(F,A)\geqslant\varepsilon\nu(A)\), then Lemma~\ref{lem:Doeblin contraction in total variation} may be applied to the Markov kernel \(P^m\). Hence \(d_{\mathrm{TV}}(\mu_1P^m,\mu_2P^m)\leqslant(1-\varepsilon)d_{\mathrm{TV}}(\mu_1,\mu_2)\).
\end{remark}


\subsection{Optimal transport}\label{subsec:Optimal transport}

We recall the basic notions of optimal transport, following Villani \cite{Villani2009}, that will be used for Gaussian comparison, the pretraining loss, metric dependence, and regularity estimates for the Gaussian output kernel. Throughout this subsection, let \(H\) be a Polish metric space with metric \(d\).

The central idea of optimal transport is to compare probability measures by considering all possible ways of coupling them. A coupling specifies how samples from one probability measure are paired with samples from another, and the optimal transport problem seeks the pairing with minimal expected cost.

\begin{definition}[Coupling]
Let \(\mu,\nu\in\mathcal P(H)\). A probability measure \(\gamma\in\mathcal P(H\times H)\) is called a \emph{coupling} of \(\mu\) and \(\nu\) if its first marginal is \(\mu\) and its second marginal is \(\nu\). The set of all couplings of \(\mu\) and \(\nu\) is denoted by \(\Pi(\mu,\nu)\).
\end{definition}

\begin{example}[Couplings of Dirac measures]
Let \(y,z\in H\). There is only one coupling of \(\delta_y\) and \(\delta_z\), namely \(\gamma=\delta_{(y,z)}\). Thus, in the special case of point masses, there is no freedom in choosing a coupling: the transport cost is simply the distance between the two points.
\end{example}

\begin{definition}[Wasserstein spaces]
For \(p\geqslant1\), the Wasserstein space \(\mathcal P_p(H)\) consists of probability measures whose \(p\)-th moment is finite,
\[
\mathcal P_p(H)=\left\{\mu\in\mathcal P(H):\int_H d(h,h_0)^p\,\mu(\mathrm dh)<\infty\right\},
\]
where \(h_0\in H\) is arbitrary. This definition does not depend on the choice of \(h_0\).
\end{definition}

\begin{definition}[Wasserstein distance]
For \(\mu,\nu\in\mathcal P_p(H)\), the \(p\)-Wasserstein distance is defined by
\[
W_p(\mu,\nu)=\left[\inf_{\gamma\in\Pi(\mu,\nu)}\int_{H\times H}d(y,z)^p\,\gamma(\mathrm dy,\mathrm dz)\right]^{1/p}.
\]
\end{definition}
The Wasserstein distance measures the minimal average cost of transporting the probability mass of \(\mu\) into the probability mass of \(\nu\). The coupling \(\gamma\) describes a possible transport plan, and the infimum selects the optimal plan. In this thesis, we will primarily use the case \(p=2\). If \(Y\) and \(Z\) are random elements on a common probability space with laws \(\mu\) and \(\nu\), then their joint distribution is an admissible coupling. Therefore, \(W_2^2(\mu,\nu)\leqslant\mathbb E[d(Y,Z)^2]\). Equality holds whenever the joint distribution of \((Y,Z)\) is an optimal coupling. A particularly useful special case is the Wasserstein distance between a probability measure and a Dirac measure.

\begin{example}\label{ex:Dirac optimal transport}
For \(h_0\in H\) and \(\mu\in\mathcal P_2(H)\), \(W_2^2(\mu,\delta_{h_0})=\int_H d(h,h_0)^2\,\mu(\mathrm dh)\).
There is only one possible coupling between \(\mu\) and \(\delta_{h_0}\), namely the measure that pairs every point \(h\) with \(h_0\). Substituting this coupling into the definition of the Wasserstein distance gives the result.
\end{example}

Next, we give the standard completeness theorem for Wasserstein spaces; see \cite[Theorem~6.18]{Villani2009}.

\begin{theorem}[Completeness of the Wasserstein space]\label{thm:Completeness of the Wasserstein space}
If \((H,d)\) is complete and separable, then \((\mathcal P_2(H),W_2)\) is also a complete and separable metric space.
\end{theorem}

From now on, we specialise to the case where the underlying space is a separable Hilbert space and introduce the covariance metric appearing in the Wasserstein distance between Gaussian measures. Let \(H\) be a real separable Hilbert space and let \(C_1,C_2\in\mathcal L_1^+(H)\).

\begin{definition}[Bures--Wasserstein distance]
The Bures--Wasserstein distance between \(C_1\) and \(C_2\) is defined by
\[d_{\mathrm{BW}}^2(C_1,C_2)=\operatorname{tr}(C_1)+\operatorname{tr}(C_2)-2\operatorname{tr}((C_2^{1/2}C_1C_2^{1/2})^{1/2}).\]
\end{definition}

The Bures--Wasserstein distance measures the difference between covariance operators in the geometry induced by optimal transport. It plays the same role for covariance operators that the Euclidean distance between means plays for deterministic points.

\begin{theorem}[Wasserstein distance between Gaussian measures]\label{thm:Wasserstein distance between Gaussian measures}
Suppose that \(m_1,m_2\in H\) and \(C_1,C_2\in\mathcal L_1^+(H)\). Let \(\mu_1=\mathcal N(m_1,C_1)\) and \(\mu_2=\mathcal N(m_2,C_2)\). Then \(\mu_1,\mu_2\in\mathcal P_2(H)\), and
\[
W_2^2(\mu_1,\mu_2)=\|m_1-m_2\|_H^2+d_{\mathrm{BW}}^2(C_1,C_2).
\]
\end{theorem}

\begin{proof}
We separate the contribution of the means from the contribution of the covariance operators. Let \(Y_1\sim\mathcal N(m_1,C_1)\) and \(Y_2\sim\mathcal N(m_2,C_2)\). Define the centred Gaussian random elements \(Z_1=Y_1-m_1\) and \(Z_2=Y_2-m_2\), such that \(Z_1\sim\mathcal N(0,C_1)\) and \(Z_2\sim\mathcal N(0,C_2).\)

Consider any coupling of the centred Gaussian measures. That is, let \((Z_1,Z_2)\) be any pair with the above marginals. Then \((Y_1,Y_2)\) is a coupling of \(\mu_1\) and \(\mu_2\) with transport cost
\[\mathbb E\lVert Y_1-Y_2\rVert_H^2=\mathbb E\lVert m_1-m_2+Z_1-Z_2\rVert_H^2 = \mathbb E\lVert m_1-m_2\rVert_H^2+\mathbb E\lVert Z_1-Z_2\rVert_H^2+2\mathbb E\langle m_1-m_2,Z_1-Z_2\rangle_H.\]
The first term is deterministic, and the cross term vanishes as \(\mathbb E[Z_1]=\mathbb E[Z_2]=0\), hence
\[\mathbb E\lVert Y_1-Y_2\rVert_H^2=\lVert m_1-m_2\rVert_H^2+\mathbb E\lVert Z_1-Z_2\rVert_H^2.\]
Taking the infimum over all couplings of the centred measures gives
\[W_2^2(\mu_1,\mu_2)=\lVert m_1-m_2\rVert_H^2+W_2^2(\mathcal N(0,C_1),\mathcal N(0,C_2)).\]
The Wasserstein distance between centred Gaussian measures on a separable Hilbert space is given by \cite{Gelbrich1990}: \(W_2^2(\mathcal N(0,C_1),\mathcal N(0,C_2))=d_{\mathrm{BW}}^2(C_1,C_2)\), completing the proof.
\end{proof}

This decomposition is fundamental for the analysis of the functional autoregressive transformer. It separates the difference between two Gaussian output distributions into two contributions: the difference between their predicted means and the difference between their predicted covariance operators.

For the finite-rank covariance parameterisation used by the functional autoregressive transformer, it is useful to relate the Bures--Wasserstein distance to covariance factors.

\begin{definition}[Hilbert--Schmidt norm]\label{def:Hilbert--Schmidt norm}
Let \(U,V\) be Hilbert spaces and let \(B:U\to V\) be a bounded linear operator. If \((e_j)_{j\geqslant1}\) is an orthonormal basis of \(U\), the Hilbert--Schmidt norm of \(B\) is defined by \(\|B\|_{\mathrm{HS}}^2=\sum_{j=1}^{\infty}\|Be_j\|_V^2\), whenever this quantity is finite.
\end{definition}

\begin{proposition}[Coupling bound for covariance factors]\label{prop:Coupling bound for covariance factors}
Let \(B_1,B_2:\mathbb R^r\to H\) be bounded linear operators and define \(C_1=B_1B_1^*\) and \(C_2=B_2B_2^*\). Then \(d_{\mathrm{BW}}^2(C_1,C_2)\leqslant\|B_1-B_2\|_{\mathrm{HS}}^2\).
\end{proposition}

\begin{proof}
The idea is to construct an explicit coupling between the two Gaussian measures using the same Gaussian noise. Let \(\xi\sim\mathcal N(0,I_r)\), where \(\xi\) is a standard Gaussian vector in \(\mathbb R^r\), and define \(Y_1=B_1\xi\) and \(Y_2=B_2\xi\).

Since \(B_1:\mathbb R^r\to H\) is linear, \(Y_1\) is Gaussian with mean \(\mathbb E[Y_1]=\mathbb E[B_1\xi]=B_1\mathbb E[\xi]=0\) and covariance \(\mathbb E[(B_1\xi)\otimes(B_1\xi)] = B_1\,\mathbb E[\xi\otimes\xi]\,B_1^* =B_1I_rB_1^* = B_1B_1^* = C_1 \).
Hence \(Y_1\sim\mathcal N(0,C_1)\) and similarly \(Y_2\sim\mathcal N(0,C_2)\). It remains now to compute:
\[
\begin{aligned}
d_{\mathrm{BW}}^2(C_1,C_2) &=W_2^2\!\left(\mathcal N(0,C_1),\mathcal N(0,C_2)\right) \\
&\leqslant\mathbb E\lVert Y_1-Y_2\rVert_H^2 \\
&=
\mathbb E\lVert(B_1-B_2)\xi\rVert_H^2 \\
&=
\mathbb E\left\langle
\xi,(B_1-B_2)^*(B_1-B_2)\xi
\right\rangle_{\mathbb R^r}\\
&=
\operatorname{tr}\left(
(B_1-B_2)^*(B_1-B_2)\mathbb E[\xi\xi^\top]
\right)\\
&=
\operatorname{tr}\left(
(B_1-B_2)^*(B_1-B_2)
\right)\\
&=
\lVert B_1-B_2\rVert_{\mathrm{HS}}^2,
\end{aligned}
\]
where the first equality follows from the Gaussian Wasserstein formula for centred Gaussian measures, the second inequality follows from the definition of the Wasserstein distance since the joint law of \((Y_1,Y_2)\) is an admissible coupling, the third equality uses \(Y_1=B_1\xi\) and \(Y_2=B_2\xi\), the fourth equality follows from the identity \(\|A\xi\|_H^2=\langle \xi,A^*A\xi\rangle_{\mathbb R^r}\) with \(A=B_1-B_2\), the fifth equality uses \(\mathbb E\langle \xi,M\xi\rangle=\operatorname{tr}(M\mathbb E[\xi\xi^\top])\), the sixth equality follows from \(\mathbb E[\xi\xi^\top]=I_r\) since \(\xi\sim\mathcal N(0,I_r)\), and the final equality is the definition of the Hilbert--Schmidt norm.
\end{proof}

This inequality controls the Wasserstein stability of Gaussian output distributions directly through covariance factors. It is used in the Lipschitz analysis of the pretraining loss and in Appendix~\ref{app:Lipschitz regularity}.


\subsection{Functional autoregressive transformer}
\label{subsec:Functional autoregressive transformer}

We now introduce the functional autoregressive transformer studied throughout the thesis. The model combines the finite-context autoregressive structure of transformer-based sequence models with a function-space architecture inspired by neural operators and continuum attention \cite{VaswaniEtAl2017,KovachkiEtAl2023,CalvelloEtAl2025}. In contrast to a standard autoregressive language model, whose observations are tokens in a finite vocabulary, the observations considered here are functions in a Hilbert space. Functional autoregressive modelling is closely related to the classical theory of functional time series, where the evolution of function-valued observations is described through dependence on preceding functions \cite{Bosq2000,HormannKokoszka2010,AueEtAl2015}.

For the concrete architecture considered in this thesis, we now specify the underlying Hilbert space. Let \(d_x,d_f\in\mathbb N\), where \(d_x\) denotes the spatial dimension of the domain and \(d_f\) denotes the number of channels of each function-valued observation. Let \(U\subset\mathbb R^{d_x}\) be a non-empty bounded open set, and specialise to
\(
H=L^2(U;\mathbb R^{d_f}).
\) Clearly, \(H\) is a real separable Hilbert space.
We fix a context length \(n\in\mathbb N\) and define the context space
\(
E=H^n.
\) At time \(t\), the model receives the functional context
\(
F_t=(f_t,\ldots,f_{t+n-1})\in E.
\)
Given \(F_t\), the transformer determines a conditional probability distribution \(Q_\Theta(F_t,\cdot)\) for the next function \(f_{t+n}\). A function is sampled from this distribution and appended to the context, after which the oldest function is removed. The procedure is then repeated recursively. This is the direct function-space analogue of finite-context autoregressive generation \cite{BrownEtAl2020,ZekriEtAl2025}.

The main difficulty is therefore to construct \(Q_\Theta(F,\cdot)\) in a way that is both expressive and mathematically well defined on function space. We use the fixed approximation space \(S_d\) introduced in Section~\ref{subsec:Finite-dimensional approximation spaces} and define the conditional output as a Gaussian measure whose mean and covariance are parameterised through coefficient vectors in \(\mathbb R^d\), which are synthesised in \(S_d\). The transformer architecture described below maps the functional context \(F\) to these coefficients.

\paragraph{Functional embeddings and bounded projection.}

As in a standard transformer, the input functions are first embedded into a latent representation. Let \(d_{\mathrm{model}}\in\mathbb N\) denote the latent channel dimension and define the latent Hilbert space \(\mathcal H=L^2(U;\mathbb R^{d_{\mathrm{model}}})\). Each function-valued observation \(f_j\in H\) is lifted to a vector-valued latent function through a bounded linear operator \(\mathsf T_{\mathrm{in}}:H\to\mathcal H\).

To distinguish the position of each function within the context window, we introduce fixed positional embeddings \(p_0,\ldots,p_{n-1}\in\mathcal H\). The \(j\)-th function is first lifted and combined with its positional embedding according to \(\widetilde v_j^{(0)}=\mathsf T_{\mathrm{in}}f_j+p_j\).
The attention mechanism introduced later involves bilinear interactions between query and key functions. To keep the latent representations uniformly bounded throughout the transformer, we project the functional tokens onto a fixed closed ball in \(\mathcal H\).

For \(R>0\), define the metric projection \(\Pi_R:\mathcal H\to\mathcal H\) onto the closed ball of radius \(R\) by
\[
\Pi_R(v)=
\begin{cases}
v, & \text{if } \|v\|_{\mathcal H}\leqslant R,\\[4pt]
{Rv/}{\|v\|_{\mathcal H}}, & \text{if } \|v\|_{\mathcal H}>R.
\end{cases}
\]
The initial latent token is defined by \(v_j^{(0)}=\Pi_R(\widetilde v_j^{(0)})\). Consequently, \(\|v_j^{(0)}\|_{\mathcal H}\leqslant R\) for every context position \(j\).
The embedded functional context is therefore \(V^{(0)}=(v_0^{(0)},\ldots,v_{n-1}^{(0)})\in\mathcal H^n\). This construction is analogous to token and positional embeddings in a conventional transformer, with the important difference that each token representation is itself a function.
The projection will also be applied after each residual update. We write \(\Pi_R^{(n)}:\mathcal H^n\to\mathcal H^n\) for the corresponding componentwise projection. As a result, every functional token entering an attention layer has norm at most \(R\).

\paragraph{Function-valued self-attention.}

We use a function-valued version of self-attention motivated by the continuum-attention framework of \cite{CalvelloEtAl2025}. The context index remains discrete, while each element of the sequence is itself a function.
Consider first a single attention head. Let \(d_q,d_v\in\mathbb N\) denote the query-key and value channel dimensions, respectively, and define \(\mathcal H_q=L^2(U;\mathbb R^{d_q})\) and \(\mathcal H_v=L^2(U;\mathbb R^{d_v})\). Let \(\mathsf Q:\mathcal H\to\mathcal H_q\), \(\mathsf K:\mathcal H\to\mathcal H_q\), and \(\mathsf V:\mathcal H\to\mathcal H_v\) be bounded linear operators, referred to as the query, key, and value operators, respectively.

Let \(V^{(\ell)}=(v_0^{(\ell)},\ldots,v_{n-1}^{(\ell)})\in\mathcal H^n\) denote the input to layer \(\ell\). By construction, every token entering the attention mechanism satisfies \(\|v_j^{(\ell)}\|_{\mathcal H}\leqslant R\). For the first layer this follows from the projection applied to the initial embeddings, while for subsequent layers it follows from the projections applied after the residual updates.

The attention score between positions \(i\) and \(j\) is defined by \(a_{ij}^{(\ell)}=\tau^{-1}\langle \mathsf Qv_i^{(\ell)},\mathsf Kv_j^{(\ell)}\rangle_{\mathcal H_q}\), where \(\tau>0\) is a temperature parameter.
The corresponding attention coefficients are \(\alpha_{ij}^{(\ell)}=\exp(a_{ij}^{(\ell)})/\sum_{k=0}^{n-1}\exp(a_{ik}^{(\ell)})\). For every \(i\), the coefficients satisfy \(\alpha_{ij}^{(\ell)}>0\) and \(\sum_{j=0}^{n-1}\alpha_{ij}^{(\ell)}=1\).
The attention output associated with position \(i\) is the function
\[\mathcal A(V^{(\ell)})_i=\sum_{j=0}^{n-1}\alpha_{ij}^{(\ell)}\mathsf Vv_j^{(\ell)}.\]
Thus, function-valued self-attention forms a weighted combination of transformed functions. The similarity between two functional tokens is measured through the Hilbert-space inner product of their query and key representations.

No causal mask is required when the architecture is used directly for next-function generation, since the model is provided only with functions belonging to the current context and not with the function being predicted. A causal mask could nevertheless be incorporated if several positions were predicted simultaneously during training.

\paragraph{Multi-head functional attention.}

As in finite-dimensional transformer architectures, several attention heads may be used in parallel. Let \(N_h\in\mathbb N\) denote the number of attention heads. For each \(h\in\{1,\ldots,N_h\}\), let \(\mathsf Q_h^{(\ell)}:\mathcal H\to\mathcal H_q\), \(\mathsf K_h^{(\ell)}:\mathcal H\to\mathcal H_q\), and \(\mathsf V_h^{(\ell)}:\mathcal H\to\mathcal H_v\) denote the query, key, and value operators of the \(h\)-th head in layer \(\ell\), and let \(\mathcal A_h^{(\ell)}\) denote the corresponding attention map. Since \(\mathcal A_h^{(\ell)}(V^{(\ell)})_i\in\mathcal H_v=L^2(U;\mathbb R^{d_v})\), the outputs of the \(N_h\) heads may be concatenated along the channel dimension to give an element of \(\mathcal H_{\mathrm{cat}}=L^2(U;\mathbb R^{N_h d_v})\). Let \(\mathsf W_O^{(\ell)}:\mathcal H_{\mathrm{cat}}\to\mathcal H\) be a bounded linear output operator. The multi-head attention output at position \(i\) is then defined by \(\mathcal A_{\mathrm{MHA}}^{(\ell)}(V^{(\ell)})_i=\mathsf W_O^{(\ell)}\operatorname{Concat}\bigl(\mathcal A_1^{(\ell)}(V^{(\ell)})_i,\ldots,\mathcal A_{N_h}^{(\ell)}(V^{(\ell)})_i\bigr)\). Thus \(\mathcal A_{\mathrm{MHA}}^{(\ell)}(V^{(\ell)})_i\in\mathcal H\) for each context position \(i\), and the output of multi-head attention is again an element of \(\mathcal H^n\).

\paragraph{Attention block and residual connection.}

The output of the multi-head attention block is combined with the original representation through a residual connection. We define \(Z^{(\ell)}=\Pi_R^{(n)}(V^{(\ell)}+\eta_\ell\mathcal A_{\mathrm{MHA}}^{(\ell)}(V^{(\ell)}))\), where \(\eta_\ell\geqslant0\) is a residual scaling parameter.
The projection is applied componentwise and ensures that every functional token in \(Z^{(\ell)}\) has norm at most \(R\). In particular, the hidden representations remain uniformly bounded after the attention update.

\paragraph{Feed-forward functional network.}

The attention block is followed by a feed-forward transformation acting pointwise in the continuous variable. Let \(d_{\mathrm{ff}}\in\mathbb N\) denote the hidden channel dimension of the feed-forward network. For each layer \(\ell\), let \(\mathsf W_1^{(\ell)}:\mathbb R^{d_{\mathrm{model}}}\to\mathbb R^{d_{\mathrm{ff}}}\) and \(\mathsf W_2^{(\ell)}:\mathbb R^{d_{\mathrm{ff}}}\to\mathbb R^{d_{\mathrm{model}}}\) be bounded linear maps, let \(b_1^{(\ell)}\in\mathbb R^{d_{\mathrm{ff}}}\) and \(b_2^{(\ell)}\in\mathbb R^{d_{\mathrm{model}}}\) be bias vectors, and let \(\sigma:\mathbb R\to\mathbb R\) be a Lipschitz nonlinear activation function applied componentwise. For \(z\in\mathcal H\), define the feed-forward functional network by \(\mathcal F_\ell(z)(x)=\mathsf W_2^{(\ell)}\sigma\!\left(\mathsf W_1^{(\ell)}z(x)+b_1^{(\ell)}\right)+b_2^{(\ell)}\). Since \(U\) is bounded and \(\sigma\) is Lipschitz, \(\mathcal F_\ell(z)\in\mathcal H\) whenever \(z\in\mathcal H\).

A second residual connection gives \(V^{(\ell+1)}=\Pi_R^{(n)}(Z^{(\ell)}+\xi_\ell\mathcal F_\ell^{(n)}(Z^{(\ell)}))\), where \(\xi_\ell\geqslant0\) is a residual scaling parameter and \(\mathcal F_\ell^{(n)}:\mathcal H^n\to\mathcal H^n\) denotes the componentwise application of \(\mathcal F_\ell\). The componentwise projection implies that \(\|v_j^{(\ell+1)}\|_{\mathcal H}\leqslant R\) for every context position \(j\). Hence the boundedness of the latent functional tokens is propagated from one transformer layer to the next. Repeating this procedure over \(L\) transformer layers gives the final latent context \(V^{(L)}=(v_0^{(L)},\ldots,v_{n-1}^{(L)})\in\mathcal H^n\).

\paragraph{Parameter constraints.}

The learnable linear operators appearing in the architecture are parameterised so that their operator norms are bounded. In particular, we restrict the model parameters so that \(\|\mathsf T_{\mathrm{in}}\|_{\mathrm{op}}\leqslant B_{\mathrm{in}}\), \(\|\mathsf Q_h^{(\ell)}\|_{\mathrm{op}}\leqslant B_Q^{(\ell,h)}\), and \(\|\mathsf K_h^{(\ell)}\|_{\mathrm{op}}\leqslant B_K^{(\ell,h)}\), and similarly \(\|\mathsf V_h^{(\ell)}\|_{\mathrm{op}}\leqslant B_V^{(\ell,h)}\), \(\|\mathsf W_O^{(\ell)}\|_{\mathrm{op}}\leqslant B_O^{(\ell)}\), and \(\|\mathsf W_i^{(\ell)}\|_{\mathrm{op}}\leqslant B_i^{(\ell)}\) for \(i\in\{1,2\}\). As assumed above, the nonlinear activation functions are Lipschitz. The same type of constraint is imposed on the learnable linear layers appearing in the mean and covariance heads introduced below.
These restrictions guarantee finite regularity constants and may be enforced through an explicit normalisation or reparameterisation of the learnable operators. The resulting bounds are derived in Appendix~\ref{app:Lipschitz regularity} and are used in the metric-dependence analysis of the pretraining loss.

\paragraph{Next-function representation.}

After \(L\) transformer layers, the functional context is represented by \(V^{(L)}=(v_0^{(L)},\ldots,v_{n-1}^{(L)})\in\mathcal H^n\). Since the model predicts the function following the complete context, we use the latent representation associated with the final context position. Define the readout operator \(\mathsf R_{\mathrm{out}}:\mathcal H^n\to\mathcal H\) by \(\mathsf R_{\mathrm{out}}(v_0,\ldots,v_{n-1})=v_{n-1}\), and set \(h_\Theta(F)=\mathsf R_{\mathrm{out}}(V^{(L)}(F))\).
The componentwise projections applied throughout the transformer ensure that \(\|h_\Theta(F)\|_{\mathcal H}\leqslant R\) for every \(F\in E\). Thus every output head receives an element of the same bounded subset of the latent function space.

\paragraph{Mean head.}

The conditional mean is parameterised through coefficient vectors in \(\mathbb R^d\), which are synthesised in the fixed approximation space \(S_d\). Let \(\widetilde{\mathsf T}_m:\mathcal H\to\mathbb R^d\) denote the mean coefficient head, which is assumed to be globally Lipschitz. Let \(\mu_0\in\mathbb R^d\) be a context-independent coefficient vector. For \(F\in E\), define \(\mu_\Theta(F)=\mu_0+\widetilde{\mathsf T}_m(h_\Theta(F))\in\mathbb R^d\).
The conditional mean function is obtained through the synthesis operator, so that \(m_\Theta(F)=\mathsf E_d\mu_\Theta(F)=\sum_{j=1}^d\mu_{j,\Theta}(F)e_j\). Consequently, \(m_\Theta(F)\in S_d\) for every \(F\in E\).
Because the latent representation is uniformly bounded and the coefficient head is Lipschitz, the family of conditional means is uniformly bounded in \(H\). The corresponding bound is established formally in Section~\ref{subsec:Well-posedness of Markov transition kernels}.

\paragraph{Covariance-factor heads.}

The covariance is parameterised in the same fixed approximation space, similar to the mean head. Fix a covariance-factor rank \(r\in\mathbb N\). For each \(k\in\{1,\ldots,r\}\), let \(\widetilde{\mathsf T}_{C,k}:\mathcal H\to\mathbb R^d\) be a globally Lipschitz coefficient head. Let \(\gamma_{0,k}\in\mathbb R^d\) be a context-independent coefficient vector. Define \(\gamma_{k,\Theta}(F)=\gamma_{0,k}+\widetilde{\mathsf T}_{C,k}(h_\Theta(F))\in\mathbb R^d\), and synthesise the corresponding functional covariance factor by \(c_{k,\Theta}(F)=\mathsf E_d\gamma_{k,\Theta}(F)\in S_d\).
Let \(G_\Theta(F)=[\gamma_{1,\Theta}(F)\ \cdots\ \gamma_{r,\Theta}(F)]\in\mathbb R^{d\times r}\) denote the matrix whose \(k\)-th column is \(\gamma_{k,\Theta}(F)\). As with the mean head, boundedness of \(h_\Theta(F)\) and Lipschitz continuity of the coefficient heads imply uniform bounds on the covariance factors. The explicit bounds are derived in Section~\ref{subsec:Well-posedness of Markov transition kernels}.

\paragraph{Conditional covariance.}

Fix a constant \(\sigma_0>0\). The conditional covariance operator is defined by \(C_\Theta(F)=\sigma_0^2\Pi_{S_d}+\sum_{k=1}^r c_{k,\Theta}(F)\otimes c_{k,\Theta}(F)\). Equivalently, under the identification of \(S_d\) with \(\mathbb R^d\) induced by the synthesis operator, the covariance is represented by the matrix \(\Sigma_\Theta(F)=\sigma_0^2I_d+G_\Theta(F)G_\Theta(F)^\top\).
The term \(\sigma_0^2\Pi_{S_d}\) supplies a fixed nondegenerate covariance floor on \(S_d\), while the covariance-factor term introduces a context-dependent and non-diagonal covariance structure. In particular, the covariance factors are not restricted to individual basis directions: correlations between the basis coefficients are represented through \(G_\Theta(F)G_\Theta(F)^\top\).
The formal verification that \(C_\Theta(F)\) is an admissible covariance operator, together with its uniform lower and upper spectral bounds, is given in Section~\ref{subsec:Well-posedness of Markov transition kernels}.

\paragraph{Gaussian output layer.}

For a context \(F\in E\), the conditional distribution of the next function is defined by \(Q_\Theta(F,\cdot)=\mathcal N(m_\Theta(F),C_\Theta(F))\).
Both the conditional mean and every stochastic direction belong to the same fixed approximation space \(S_d\). Hence all conditional output distributions are supported on the same subspace of \(H\), independently of the context \(F\). Moreover, the positive covariance floor makes the Gaussian distribution nondegenerate when it is regarded as a probability measure on \(S_d\).
This common-support structure is important for the subsequent Markov-chain analysis. In particular, all conditional Gaussian laws admit densities with respect to the same \(d\)-dimensional Lebesgue measure on \(S_d\), which will later allow a common minorising measure to be constructed.

\paragraph{Function generation.}

We now use the finite-rank Gaussian construction of Section \ref{subsec:Finite-rank Gaussian covariance construction} to obtain an explicit sampling representation of the conditional output distribution. Let \(\zeta_t=(\zeta_{t,1},\ldots,\zeta_{t,d})\sim\mathcal N(0,I_d)\) and \(\xi_t=(\xi_{t,1},\ldots,\xi_{t,r})\sim\mathcal N(0,I_r)\) be independent standard Gaussian vectors, and assume that they are independent of the current context \(F_t\). Given \(F_t=(f_t,\ldots,f_{t+n-1})\in E\), define the next function by \(f_{t+n}=m_\Theta(F_t)+\sigma_0\mathsf E_d\zeta_t+\sum_{k=1}^r\xi_{t,k}c_{k,\Theta}(F_t)\).

We verify that this random element has conditional distribution \(Q_\Theta(F_t,\cdot)\). Since \(\mathsf E_d\zeta_t=\sum_{j=1}^d\zeta_{t,j}e_j\), the random perturbation can be written as \(\sigma_0\mathsf E_d\zeta_t+\sum_{k=1}^r\xi_{t,k}c_{k,\Theta}(F_t)=\sum_{j=1}^d\zeta_{t,j}(\sigma_0e_j)+\sum_{k=1}^r\xi_{t,k}c_{k,\Theta}(F_t)\). Therefore, conditionally on \(F_t\), the generated function \(f_{t+n}\) can be written in the form considered in Theorem \ref{thm:Finite-rank Gaussian construction}, with the \(d+r\) covariance factors \(\sigma_0e_1,\ldots,\sigma_0e_d,c_{1,\Theta}(F_t),\ldots,c_{r,\Theta}(F_t)\).

Now let us observe that, since \(\zeta_t\) and \(\xi_t\) are independent standard Gaussian vectors, their concatenation \(\eta_t=(\zeta_{t,1},\ldots,\zeta_{t,d},\xi_{t,1},\ldots,\xi_{t,r})\) is a standard Gaussian vector in \(\mathbb R^{d+r}\). Hence, by Theorem \ref{thm:Finite-rank Gaussian construction}, conditionally on \(F_t\), the random element \(f_{t+n}\) is Gaussian with mean \(m_\Theta(F_t)\) and covariance operator \(\sum_{j=1}^d(\sigma_0e_j)\otimes(\sigma_0e_j)+\sum_{k=1}^rc_{k,\Theta}(F_t)\otimes c_{k,\Theta}(F_t)\).

Lastly, using the bilinearity of the rank-one operator, we have \((\sigma_0e_j)\otimes(\sigma_0e_j)=\sigma_0^2(e_j\otimes e_j)\), and therefore \(\sum_{j=1}^d(\sigma_0e_j)\otimes(\sigma_0e_j)=\sigma_0^2\sum_{j=1}^de_j\otimes e_j\). By the definition of the orthogonal projection onto \(S_d\), \(\Pi_{S_d}=\sum_{j=1}^de_j\otimes e_j\). Consequently, the covariance operator of \(f_{t+n}\), conditionally on \(F_t\), is \(\sigma_0^2\Pi_{S_d}+\sum_{k=1}^rc_{k,\Theta}(F_t)\otimes c_{k,\Theta}(F_t)=C_\Theta(F_t)\).
It follows that \(f_{t+n}\mid F_t\sim\mathcal N(m_\Theta(F_t),C_\Theta(F_t))=Q_\Theta(F_t,\cdot)\).

\paragraph{Context update and recursive generation.}

After generating \(f_{t+n}\), the context window is shifted by removing its oldest entry and appending the newly generated function, giving \(F_{t+1}=(f_{t+1},\ldots,f_{t+n})\).
The updated context is then passed through the same transformer and output heads, and the procedure is repeated recursively. Thus the parameters \(\Theta\) determine a conditional Gaussian distribution for every functional context and, together with the shift-and-append update, specify the complete inference mechanism of the functional autoregressive transformer. The probabilistic properties required for well-posedness and the subsequent Markov-chain analysis are established separately in Section~\ref{sec:General state-space Markov chain framework for functional autoregressive transformers}.

\section{General state-space Markov chain framework for functional autoregressive transformers}\label{sec:General state-space Markov chain framework for functional autoregressive transformers}

In this section, we develop a general state-space Markov-chain formulation of the inference mechanism of the functional autoregressive transformer introduced in Section~\ref{subsec:Functional autoregressive transformer}. The central idea is analogous to the finite-state formulation of autoregressive language models considered by Zekri et al.~\cite{ZekriEtAl2025}: although the sequence of individual outputs is not, in general, a first-order Markov chain, the process becomes Markovian once the state is taken to be the entire context window.

We first formulate the context transition kernel and establish its well-posedness. We then derive a Lyapunov drift condition and an \(n\)-step minorisation, which together yield existence, uniqueness, and uniform geometric convergence to an invariant probability measure. We conclude by discussing the interpretation and modelling implications of the result.

\subsection{Markov chain formulation}\label{subsec:Markov chain formulation}
Let \(H\) be a real separable Hilbert space, fix a context length \(n\in\mathbb N\), and define the context space \(E=H^n\). A first attempt might be to regard the generated sequence of functions \((f_t)_{t\geqslant0}\) itself as a first-order Markov chain on \(H\). This is generally incorrect because the conditional distribution of the next function depends on the full context of the previous \(n\) functions, rather than only on the most recently generated function.

Indeed, if \(F_t=(f_t,\ldots,f_{t+n-1})\in E\), then the functional autoregressive transformer generates the next function according to \(f_{t+n}\mid F_t\sim Q_\Theta(F_t,\cdot)=\mathcal N(m_\Theta(F_t),C_\Theta(F_t))\). In general, this conditional law depends on all components of \(F_t\).
As in finite-context autoregressive language models, the process becomes first-order Markov once the state is taken to be the entire context. We therefore regard \(F_t=(f_t,\ldots,f_{t+n-1})\) as the state of the process at time \(t\). For the moment, suppose that \(Q_\Theta\) is a Markov transition kernel from \(E\) to \(H\).

\begin{definition}[]\label{def:FAT Markov transition kernel}
For \(F=(f_0,\ldots,f_{n-1})\in E\) and \(A\in\mathcal B(E)\), define the \emph{functional autoregressive transformer transition kernel} by
\[P_\Theta(F,A)=\int_H\mathbf 1_A(f_1,\ldots,f_{n-1},g)\,Q_\Theta(F,\mathrm dg).\]
When \(n=1\), this is interpreted as \(P_\Theta(f_0,A)=Q_\Theta(f_0,A)\).
\end{definition}

For a fixed context \(F\), the measure \(P_\Theta(F,\cdot)\) is the law of the next context: the first \(n-1\) coordinates are obtained by deterministically shifting the current context, while the final coordinate is sampled from \(Q_\Theta(F,\cdot)\).

\begin{proposition}[Markov transition kernel]\label{prop:Markov transition kernel}
Suppose that \(Q_\Theta\) is a Markov transition kernel from \(E\) to \(H\). Then \(P_\Theta\) is a Markov transition kernel from \(E\) to \(E\).
\end{proposition}

\begin{proof}
We verify the two kernel properties in Definition \ref{def:Markov transition kernel} separately.

\textit{Step 1: Probability measure for a fixed context.}
Fix \(F=(f_0,\ldots,f_{n-1})\in E\) and define \(T_F:H\to E\) by
\(
T_F(g)=(f_1,\ldots,f_{n-1},g).
\)
Then \(P_\Theta(F,A)=Q_\Theta(F,T_F^{-1}(A))\). Thus \(P_\Theta(F,\cdot)\) is the pushforward of the probability measure \(Q_\Theta(F,\cdot)\) under \(T_F\), and is therefore a probability measure on \(E\).

\textit{Step 2: Mesurability.}
Define \(T:E\times H\to E\) by
\(
T\bigl((f_0,\ldots,f_{n-1}),g\bigr)
=(f_1,\ldots,f_{n-1},g).
\)
This map is continuous. Since \(Q_\Theta\) is a Markov transition kernel, the standard kernel-integration lemma \parencite[Lemma~3.2]{Kallenberg2021} shows that
\(
F\longmapsto\int_H\mathbf 1_A\bigl(T(F,g)\bigr)Q_\Theta(F,\mathrm dg)
\)
is Borel measurable for every \(A\in\mathcal B(E)\). This is precisely \(F\mapsto P_\Theta(F,A)\).
\end{proof}

\begin{theorem}[Markov chain formulation]\label{thm:Markov chain formulation}
Let \(\mu\in\mathcal P(E)\) be an initial distribution and suppose that \(Q_\Theta\) is a Markov transition kernel from \(E\) to \(H\). Then recursive functional autoregressive generation determines a time-homogeneous Markov chain \((F_t)_{t\in\mathbb N_0}\) on \(E=H^n\) with initial distribution \(\mu\) and transition kernel \(P_\Theta\).
\end{theorem}

\begin{proof}
Proposition~\ref{prop:Markov transition kernel} gives a Markov transition kernel on \(E\). Theorem~\ref{thm:Existence of the Markov process} therefore gives a time-homogeneous Markov chain with this kernel and the prescribed initial law. By Definition~\ref{def:FAT Markov transition kernel}, its update is exactly
\(
F_t=(f_t,\ldots,f_{t+n-1})
\longmapsto
F_{t+1}=(f_{t+1},\ldots,f_{t+n}),
\)
where \(f_{t+n}\mid F_t\sim Q_\Theta(F_t,\cdot)\).
\end{proof}

The remaining well-posedness question is whether the coefficient-space Gaussian output of the architecture indeed defines a Markov transition kernel for every context.

\subsection{Well-posedness of the Gaussian output kernel}
\label{subsec:Well-posedness of Markov transition kernels}

We now verify that the output layer introduced in Section~\ref{subsec:Functional autoregressive transformer} defines a well-posed conditional probability distribution for every functional context and that these distributions form a Markov transition kernel.

Recall that the mean is given by \(m_\Theta(F)=\mathsf E_d\mu_\Theta(F)\), where
\(
\mu_\Theta(F)
=
\mu_0+\widetilde{\mathsf T}_m\bigl(h_\Theta(F)\bigr)
\), and the covariance operator is given by
\(
C_\Theta(F)
=
\sigma_0^2\Pi_{S_d}
+
\sum_{k=1}^r
c_{k,\Theta}(F)\otimes c_{k,\Theta}(F),\)
for 
\(
\sigma_0>0,
\)
where
\(c_{k,\Theta}(F)=\mathsf E_d\gamma_{k,\Theta}(F)\)
with
\(
\gamma_{k,\Theta}(F)
=
\gamma_{0,k}
+\widetilde{\mathsf T}_{C,k}\bigl(h_\Theta(F)\bigr).
\)


The first step is to obtain uniform bounds on the quantities predicted by the output heads. These bounds follow from the bounded latent representation \(\|h_\Theta(F)\|_{\mathcal H}\leqslant R\) for every \(F\in E\), together with the Lipschitz assumptions imposed on the coefficient heads.

\begin{proposition}[Uniform bound for the mean]
\label{prop:Uniform bound for the mean}
Suppose that \(\widetilde{\mathsf T}_m:\mathcal H\to\mathbb R^d\) is \(L_m^{\mathrm{head}}\)-Lipschitz. Then
\(
\sup_{F\in E}\|m_\Theta(F)\|_H
\leqslant
M_m,
\)
where
\(
M_m
=
\|\mu_0\|_2
+
\|\widetilde{\mathsf T}_m(0)\|_2
+
L_m^{\mathrm{head}}R
<\infty.
\)
\end{proposition}

\begin{proof}
Fix \(F\in E\). Let us compute
\[
\begin{aligned}
\|m_\Theta(F)\|_H
&=
\|\mu_\Theta(F)\|_2\\
&=
\|
\mu_0
+
\widetilde{\mathsf T}_m
\bigl(h_\Theta(F)\bigr)
\|_2\\
&\leqslant
\|\mu_0\|_2
+
\|
\widetilde{\mathsf T}_m
\bigl(h_\Theta(F)\bigr)
\|_2\\
&\leqslant
\|\mu_0\|_2
+
\|\widetilde{\mathsf T}_m(0)\|_2
+
\|
\widetilde{\mathsf T}_m
\bigl(h_\Theta(F)\bigr)
-
\widetilde{\mathsf T}_m(0)
\|_2
\\
&\leqslant
\|\mu_0\|_2
+
\|\widetilde{\mathsf T}_m(0)\|_2
+
L_m^{\mathrm{head}}
\|h_\Theta(F)\|_{\mathcal H}
\\
&\leqslant
\|\mu_0\|_2
+
\|\widetilde{\mathsf T}_m(0)\|_2
+
L_m^{\mathrm{head}}R
\\
&=:
M_m.
\end{aligned}
\]
The first equality follows from the isometry of \(\mathsf E_d\), see Proposition \ref{prop:Basic identities for the approximation space}. The second equality is the definition of \(\mu_\Theta(F)\). The first inequality follows from the triangle inequality. The second inequality is obtained by adding and subtracting \(\widetilde{\mathsf T}_m(0)\) and applying the triangle inequality. The third inequality uses the \(L_m^{\mathrm{head}}\)-Lipschitz continuity of \(\widetilde{\mathsf T}_m\), while the final inequality uses the uniform latent bound \(\|h_\Theta(F)\|_{\mathcal H}\leqslant R\). Since the resulting bound is independent of \(F\), taking the supremum over \(F\in E\) proves the result.
\end{proof}

\begin{proposition}[Uniform bounds for covariance factors]
\label{prop:Uniform bounds for covariance factors}
Suppose that, for each \(k\in\{1,\ldots,r\}\), the map \(\widetilde{\mathsf T}_{C,k}:\mathcal H\to\mathbb R^d\) is \(L_{C,k}^{\mathrm{head}}\)-Lipschitz. Then
\(
\sup_{F\in E}
\|c_{k,\Theta}(F)\|_H
\leqslant
M_{C,k},
\)
where
\(
M_{C,k}
=
\|\gamma_{0,k}\|_2
+
\|\widetilde{\mathsf T}_{C,k}(0)\|_2
+
L_{C,k}^{\mathrm{head}}R
<\infty.
\)
\end{proposition}

\begin{proof}
Fix \(F\in E\) and \(k\in\{1,\ldots,r\}\). Let us compute
\[
\begin{aligned}
\|c_{k,\Theta}(F)\|_H
&=
\|\gamma_{k,\Theta}(F)\|_2\\
&=
\|
\gamma_{0,k}
+
\widetilde{\mathsf T}_{C,k}
\bigl(h_\Theta(F)\bigr)
\|_2\\
&\leqslant
\|\gamma_{0,k}\|_2
+
\|
\widetilde{\mathsf T}_{C,k}
\bigl(h_\Theta(F)\bigr)
\|_2\\
&\leqslant
\|\gamma_{0,k}\|_2
+
\|\widetilde{\mathsf T}_{C,k}(0)\|_2
+
\|
\widetilde{\mathsf T}_{C,k}
\bigl(h_\Theta(F)\bigr)
-
\widetilde{\mathsf T}_{C,k}(0)
\|_2
\\
&\leqslant
\|\gamma_{0,k}\|_2
+
\|\widetilde{\mathsf T}_{C,k}(0)\|_2
+
L_{C,k}^{\mathrm{head}}
\|h_\Theta(F)\|_{\mathcal H}
\\
&\leqslant
\|\gamma_{0,k}\|_2
+
\|\widetilde{\mathsf T}_{C,k}(0)\|_2
+
L_{C,k}^{\mathrm{head}}R
\\
&=:
M_{C,k}.
\end{aligned}
\]
The first equality follows from the isometry of \(\mathsf E_d\), see Proposition \ref{prop:Basic identities for the approximation space}. The second equality is the definition of \(\gamma_{k,\Theta}(F)\). The first inequality follows from the triangle inequality. The second inequality is obtained by adding and subtracting \(\widetilde{\mathsf T}_{C,k}(0)\) and applying the triangle inequality. The third inequality uses the \(L_{C,k}^{\mathrm{head}}\)-Lipschitz continuity of \(\widetilde{\mathsf T}_{C,k}\), while the final inequality uses the uniform latent bound \(\|h_\Theta(F)\|_{\mathcal H}\leqslant R\). Since the resulting bound is independent of \(F\), taking the supremum over \(F\in E\) proves the result.
\end{proof}

We can now verify that the covariance produced by the architecture is an admissible covariance operator and is uniformly nondegenerate on the approximation space.

\begin{proposition}[Admissibility, spectral bounds, and coefficient representation]
\label{prop:Admissibility of the covariance operator}
Let \(\underline\lambda=\sigma_0^2\) and let
\(
\overline\lambda
=
\sigma_0^2
+
\sum_{k=1}^rM_{C,k}^2.
\)
For every \(F\in E\), the following statements hold.
\begin{enumerate}
    \item \(m_\Theta(F)\in S_d\) and \(C_\Theta(F)\in\mathcal L_1^+(H)\).

    \item The range of \(C_\Theta(F)\) is \(S_d\), and its restriction to \(S_d\) satisfies
    \(
    \underline\lambda I_{S_d}
    \preceq
    C_\Theta(F)|_{S_d}
    \preceq
    \overline\lambda I_{S_d}.
    \)

    \item If
    \(
    G_\Theta(F)
    =
    [
    \gamma_{1,\Theta}(F)
    ,
    \ldots
    ,
    \gamma_{r,\Theta}(F)
    ]
    \in\mathbb R^{d\times r}
    \)
    and \(\Sigma_\Theta(F)=\sigma_0^2I_d+G_\Theta(F)G_\Theta(F)^\top\), then we have
    \(
    C_\Theta(F)
    =
    \mathsf E_d
    \Sigma_\Theta(F)
    \mathsf E_d^*.
    \)

    \item The trace is
    \(
    \operatorname{tr}\bigl(C_\Theta(F)\bigr)
    =
    d\sigma_0^2
    +
    \sum_{k=1}^r
    \|c_{k,\Theta}(F)\|_H^2.
    \)
\end{enumerate}
\end{proposition}

\begin{proof}
Fix \(F\in E\).

\textit{Step 1: Admissibility.}
By construction, \(m_\Theta(F)=\mathsf E_d\mu_\Theta(F)\in S_d\). Proposition~\ref{prop:Basic identities for the approximation space} gives \(\Pi_{S_d}=\mathsf E_d\mathsf E_d^*\), so \(\Pi_{S_d}\) is the orthogonal projection onto \(S_d\). It is therefore finite rank, self-adjoint, and non-negative. For every \(k\), by Proposition \ref{prop:Finite-rank covariance}, the rank-one operator \(c_{k,\Theta}(F)\otimes c_{k,\Theta}(F)\) has the same properties. Hence \(C_\Theta(F)\) is self-adjoint, non-negative, and finite rank, and therefore belongs to \(\mathcal L_1^+(H)\).

\textit{Step 2: Spectral bounds.}
Let \(u\in S_d\). Since \(\Pi_{S_d}u=u\),
\(
\langle C_\Theta(F)u,u\rangle_H
=
\sigma_0^2\|u\|_H^2
+
\sum_{k=1}^r
|\langle u,c_{k,\Theta}(F)\rangle_H|^2.
\)
Dropping the non-negative sum gives \(\langle C_\Theta(F)u,u\rangle_H\geqslant\underline\lambda\|u\|_H^2\), which proves the lower spectral bound.

For the upper bound, Cauchy--Schwarz gives
\(
|\langle u,c_{k,\Theta}(F)\rangle_H|^2
\leqslant
\|u\|_H^2
\|c_{k,\Theta}(F)\|_H^2.
\)
Using Proposition~\ref{prop:Uniform bounds for covariance factors},
\(
\langle C_\Theta(F)u,u\rangle_H
\leqslant
\left(
\sigma_0^2
+
\sum_{k=1}^rM_{C,k}^2
\right)
\|u\|_H^2
=
\overline\lambda\|u\|_H^2.
\)
This proves the upper spectral bound.

Every term in \(C_\Theta(F)\) has range contained in \(S_d\), so \(\operatorname{Ran}(C_\Theta(F))\subseteq S_d\). The lower spectral bound above shows that \(C_\Theta(F)|_{S_d}\) is invertible, and hence \(\operatorname{Ran}(C_\Theta(F))=S_d\).

\textit{Step 3: Coefficient representation.}
Since \(\mathsf E_d\mathsf E_d^*=\Pi_{S_d}\), we have \(\sigma_0^2\mathsf E_dI_d\mathsf E_d^*=\sigma_0^2\Pi_{S_d}\). Moreover, the columns of \(\mathsf E_dG_\Theta(F)\) are \(c_{1,\Theta}(F),\ldots,c_{r,\Theta}(F)\). Therefore,
\(
\mathsf E_d
G_\Theta(F)G_\Theta(F)^\top
\mathsf E_d^*
=
\sum_{k=1}^r
c_{k,\Theta}(F)
\otimes
c_{k,\Theta}(F).
\)
Combining the two identities gives \(C_\Theta(F)=\mathsf E_d\Sigma_\Theta(F)\mathsf E_d^*\).

\textit{Step 4: Trace formula.}
Since \(\operatorname{tr}(\Pi_{S_d})=d\) and \(\operatorname{tr}(c\otimes c)=\|c\|_H^2\), linearity of the trace gives the result.
\end{proof}

The preceding proposition establishes the operator-theoretic properties required of the covariance. We now use them to identify the conditional Gaussian law and its support.

\begin{theorem}[Existence, uniqueness, and support of the Gaussian output measure]
\label{thm:Existence and uniqueness of the Gaussian output measure}
For every \(F\in E\), there exists a unique Borel Gaussian probability measure on \(H\) with mean \(m_\Theta(F)\) and covariance operator \(C_\Theta(F)\). We denote this measure by
\(
Q_\Theta(F,\cdot)
=
\mathcal N
\bigl(
m_\Theta(F),
C_\Theta(F)
\bigr).
\)
Moreover, \(\operatorname{supp}Q_\Theta(F,\cdot)=S_d\) for every \(F\in E\), where the support is taken with respect to the topology inherited from \(H\).
\end{theorem}

\begin{proof}
Fix \(F\in E\). We evaluate each statement in turn.

\textit{Step 1: Existence and uniqueness.}
Proposition~\ref{prop:Admissibility of the covariance operator} gives \(m_\Theta(F)\in H\) and \(C_\Theta(F)\in\mathcal L_1^+(H)\). Theorem~\ref{thm:Existence and uniqueness of Gaussian random elements} therefore gives a unique Borel Gaussian probability measure \(Q_\Theta(F,\cdot)=\mathcal N(m_\Theta(F),C_\Theta(F))\).

\textit{Step 2: Coefficient-space representation.}
Recall that \(\Sigma_\Theta(F)=\sigma_0^2I_d+G_\Theta(F)G_\Theta(F)^\top\). For every \(a\in\mathbb R^d\),
\(
a^\top\Sigma_\Theta(F)a
=
\sigma_0^2\|a\|_2^2
+
\|G_\Theta(F)^\top a\|_2^2
\geqslant
\sigma_0^2\|a\|_2^2
\) and since \(\sigma_0>0\), \(\Sigma_\Theta(F)\) is positive definite.

Let \(Z_F\sim\mathcal N(\mu_\Theta(F),\Sigma_\Theta(F))\) in \(\mathbb R^d\) and define \(Y_F=\mathsf E_dZ_F\). Then \(\mathbb E[Y_F]=\mathsf E_d\mu_\Theta(F)=m_\Theta(F)\), while Proposition~\ref{prop:Admissibility of the covariance operator} gives
\(
\mathsf E_d
\Sigma_\Theta(F)
\mathsf E_d^*
=
C_\Theta(F).
\)
By uniqueness of Gaussian measures with prescribed mean and covariance,
\(
Q_\Theta(F,\cdot)
=
(\mathsf E_d)_\#
\mathcal N
\bigl(
\mu_\Theta(F),
\Sigma_\Theta(F)
\bigr).
\)

\textit{Step 3: Common support.}
Since \(\Sigma_\Theta(F)\) is positive definite, \(Z_F\) has support equal to all of \(\mathbb R^d\). The synthesis operator \(\mathsf E_d:\mathbb R^d\to S_d\) is an isometric isomorphism. Its pushforward therefore has support \(S_d\), and hence
\(
\operatorname{supp}
Q_\Theta(F,\cdot)
=
S_d
\)
for every \(F\in E\).
\end{proof}

The theorem explains why the mean and covariance are parameterised in the same fixed approximation space. All conditional output measures have the same support \(S_d\), which permits the construction of a context-independent minorising measure later in the ergodicity analysis.

To obtain a probability kernel, it remains to verify that the Gaussian output measure depends measurably on the context. We first establish continuity of the model parameter maps, but first a technical lemma.

\begin{lemma}[Trace norm of rank-one operators]
\label{lem:rank-one operations}
For \(u,v\in H\), (i)
\(
\|u\otimes v\|_1
=
\|u\|_H\|v\|_H,
\)
and (ii)
\(
\|u\otimes u-v\otimes v\|_1
\leqslant
\bigl(
\|u\|_H+\|v\|_H
\bigr)
\|u-v\|_H.
\)
\end{lemma}

\begin{proof}
For the first identity, recall that \((u\otimes v)h=\langle h,v\rangle_Hu\) for \(h\in H\), and that \((u\otimes v)^*=v\otimes u\). Hence \((u\otimes v)^*(u\otimes v)=(v\otimes u)(u\otimes v)=\|u\|_H^2(v\otimes v)\). The operator \(v\otimes v\) has the single nonzero eigenvalue \(\|v\|_H^2\), so \((u\otimes v)^*(u\otimes v)\) has the single nonzero eigenvalue \(\|u\|_H^2\|v\|_H^2\). Therefore \(u\otimes v\) has the single nonzero singular value \(\|u\|_H\|v\|_H\), and thus \(\|u\otimes v\|_1=\|u\|_H\|v\|_H\).

For the second identity, write \(u\otimes u-v\otimes v=(u-v)\otimes u+v\otimes(u-v)\). By the triangle inequality and part (i), \(\|u\otimes u-v\otimes v\|_1\leqslant\|(u-v)\otimes u\|_1+\|v\otimes(u-v)\|_1=\|u-v\|_H\|u\|_H+\|v\|_H\|u-v\|_H=(\|u\|_H+\|v\|_H)\|u-v\|_H\), which proves the result.
\end{proof}

\begin{proposition}[Continuity of the Gaussian model parameters]
\label{prop:Continuity of the covariance map}
The parameter maps \(m_\Theta:E\to H\) and \(c_{k,\Theta}:E\to H\), for \(k\in\{1,\ldots,r\}\), are continuous. Moreover,
\(
C_\Theta:E\to\mathcal L_1^+(H)
\)
is continuous when \(\mathcal L_1^+(H)\) is equipped with the trace norm.
\end{proposition}

\begin{proof}
The functional autoregressive transformer consists of continuous operations: bounded linear maps, Lipschitz nonlinearities, finite-dimensional softmax maps, residual connections, and metric projections onto closed balls. Consequently, \(h_\Theta:E\to\mathcal H\) is continuous.

The coefficient heads are globally Lipschitz and hence continuous. Therefore the maps \(F\mapsto\mu_\Theta(F)\) and \(F\mapsto\gamma_{k,\Theta}(F)\) are continuous into \(\mathbb R^d\). Since \(\mathsf E_d:\mathbb R^d\to H\) is bounded, \(m_\Theta(F)=\mathsf E_d\mu_\Theta(F)\) and \(c_{k,\Theta}(F)=\mathsf E_d\gamma_{k,\Theta}(F)\) are continuous as maps into \(H\).

It remains to prove trace-norm continuity of the covariance map. Let \(F_j\to F\) in \(E\). Then we have \(c_{k,\Theta}(F_j)\to c_{k,\Theta}(F)\) in \(H\) for every \(k\). Since the covariance floor is independent of the context,
\[
\begin{aligned}
C_\Theta(F_j)-C_\Theta(F)
=
\sum_{k=1}^r
\Bigl(
c_{k,\Theta}(F_j)\otimes c_{k,\Theta}(F_j)
-
c_{k,\Theta}(F)\otimes c_{k,\Theta}(F)
\Bigr).
\end{aligned}
\]
Lemma~\ref{lem:rank-one operations} gives
\[
\|C_\Theta(F_j)-C_\Theta(F)\|_1
\leqslant
\sum_{k=1}^r
\Bigl(
\|c_{k,\Theta}(F_j)\|_H
+
\|c_{k,\Theta}(F)\|_H
\Bigr)
\|c_{k,\Theta}(F_j)-c_{k,\Theta}(F)\|_H.
\]
For each \(k\), the final factor tends to zero. Since \(c_{k,\Theta}(F_j)\to c_{k,\Theta}(F)\), the first factor remains bounded. The sum is finite, so \(\|C_\Theta(F_j)-C_\Theta(F)\|_1\to0\).
\end{proof}

We now use a standard continuity property of Gaussian measures.

\begin{proposition}[Continuity of the Gaussian parameter map]
\label{prop:Continuity of the Gaussian parameter map}
The map
\(
\mathfrak G:
H\times\mathcal L_1^+(H)
\to
\mathcal P(H)\)
defined by \(
\mathfrak G(m,C)
=
\mathcal N(m,C)
\)
is continuous when \(H\) carries its norm topology, \(\mathcal L_1^+(H)\) the trace-norm topology, and \(\mathcal P(H)\) the topology of weak convergence.
\end{proposition}

\begin{proof}
Suppose that \(m_j\to m\) in \(H\) and \(C_j\to C\) in trace norm. Then, we have the weak convergence
\(
\mathcal N(m_j,C_j)
\Rightarrow
\mathcal N(m,C)
\)
by \cite[Theorem~2.7.21]{Bogachev2018}. Hence \(\mathfrak G\) is sequentially continuous. Since the spaces involved are metrizable under the stated topologies, sequential continuity is equivalent to continuity.
\end{proof}

\begin{theorem}[Gaussian output kernel]
\label{thm:Gaussian output kernel}
The Gaussian output map
\(
Q_\Theta:
E\times\mathcal B(H)
\to
[0,1]\)
given by \(
Q_\Theta(F,A)
=
\mathcal N
\bigl(
m_\Theta(F),
C_\Theta(F)
\bigr)(A),
\)
is a Markov transition kernel from \(E\) to \(H\).
\end{theorem}

\begin{proof}
We verify the two requirements in Definition \ref{def:Markov transition kernel}

\textit{Step 1: Probability measure.}
For every \(F\in E\), Theorem~\ref{thm:Existence and uniqueness of the Gaussian output measure} shows that \(Q_\Theta(F,\cdot)=\mathcal N(m_\Theta(F),C_\Theta(F))\) is a Borel probability measure on \(H\).

\textit{Step 2: Measurable dependence on the context.}
Proposition~\ref{prop:Continuity of the covariance map} gives continuity of \(F\mapsto(m_\Theta(F),C_\Theta(F))\) from \(E\) to \(H\times\mathcal L_1^+(H)\). Proposition~\ref{prop:Continuity of the Gaussian parameter map} therefore implies that \(F\mapsto Q_\Theta(F,\cdot)\) is continuous, and hence Borel measurable, as a map from \(E\) into \(\mathcal P(H)\).

For every \(A\in\mathcal B(H)\), the evaluation map \(\mu\mapsto\mu(A)\) is Borel measurable on \(\mathcal P(H)\). Consequently, \(F\mapsto Q_\Theta(F,A)\) is Borel measurable. Both requirements are satisfied, so \(Q_\Theta\) is a Markov transition kernel from \(E\) to \(H\).
\end{proof}

\begin{corollary}[Well-posed functional autoregressive chain]
\label{cor:Well-posed functional autoregressive chain}
The functional autoregressive transformer transition kernel
\[
P_\Theta(F,A)
=
\int_H
\mathbf 1_A
(f_1,\ldots,f_{n-1},g)
\,Q_\Theta(F,\mathrm dg)
\]
defines a Markov transition kernel on \(E\). Consequently, for every initial probability measure on \(E\), the functional autoregressive transformer defines a time-homogeneous general-state-space Markov chain.
\end{corollary}

\begin{proof}
Theorem~\ref{thm:Gaussian output kernel} establishes that \(Q_\Theta\) is a Markov transition kernel from \(E\) to \(H\). Proposition~\ref{prop:Markov transition kernel} then shows that \(P_\Theta\) is a Markov transition kernel on \(E\). Finally, Theorem~\ref{thm:Markov chain formulation} gives the corresponding time-homogeneous Markov chain.
\end{proof}

\subsection{Drift and minorisation ergodicity}\label{subsec:Drift and minorisation ergodicity}

We now study the long-time behaviour of the recursive inference chain. The architecture provides two key ingredients: the deterministic context shift progressively removes the initial state, while the Gaussian output introduces uniformly nondegenerate randomness on the fixed space \(S_d\).

We first use the uniform conditional second-moment bound to derive a geometric drift inequality in a weighted context norm. We then use the common Gaussian support and covariance bounds to obtain an output minorisation, which becomes a global \(n\)-step minorisation of the context kernel once the initial context has been shifted out. The \(n\)-step Doeblin condition yields existence, uniqueness, and uniform geometric convergence in total variation, while the drift estimate provides moment control.

\begin{definition}[Weighted context norm]\label{def:Weighted context norm}
Let \(\beta>1\). For \(F=(f_0,\ldots,f_{n-1})\in E\), define the \emph{weighted context norm}
\(
\|F\|_\beta^2
=\sum_{j=0}^{n-1}\beta^j\|f_j\|_H^2
\) and its corresponding \emph{Lyapunov function}
\(
V(F)=1+\|F\|_\beta^2.
\)
\end{definition}

Since \(n\) is finite and \(\beta>1\), the weighted norm is equivalent to the usual product norm on \(E=H^n\). Indeed, for every \(F=(f_0,\ldots,f_{n-1})\in E\), \(\|F\|_{H^n}^2\leqslant\|F\|_\beta^2\leqslant\beta^{n-1}\|F\|_{H^n}^2\). Consequently, \((E,\|\cdot\|_\beta)\) has the same topology as \(H^n\). In particular, it remains a complete and separable Hilbert space.

The weighted norm is chosen to reflect the shift structure of the autoregressive update. As context entries move one position to the left, their weights decrease by a factor \(\beta^{-1}\). To see this, let \(F=(f_0,\ldots,f_{n-1})\) and \(G=(g_0,\ldots,g_{n-1})\). After one transition, the deterministic coordinates are shifted one position to the left and their contribution to the squared weighted distance becomes
\[
\sum_{j=1}^{n-1}\beta^{j-1}\|f_j-g_j\|_H^2=\beta^{-1}\sum_{j=1}^{n-1}\beta^j\|f_j-g_j\|_H^2\leqslant\beta^{-1}\|F-G\|_\beta^2.
\]
Thus, the deterministic shift produces a contraction, while the newly generated coordinate will be controlled separately through its conditional second moment. The next step is therefore to control the only genuinely new contribution to the context norm: the generated function. The bounded mean and covariance factors imply that its conditional second moment is uniformly bounded over all possible contexts.


\begin{lemma}[Uniform conditional second moment]\label{lem:Uniform conditional second moment}
For every \(F\in E\),
\(
\int_H\|g\|_H^2Q_\Theta(F,\mathrm dg)
\leqslant M_2,
\)
where
\(
M_2=M_m^2+d\sigma_0^2+\sum_{k=1}^rM_{C,k}^2<\infty.
\)
\end{lemma}

\begin{proof}
By Theorem \ref{thm:Covariance operator of a Gaussian random element}, we have the identity \(
\int_H\|g\|_H^2Q_\Theta(F,\mathrm dg)
=\|m_\Theta(F)\|_H^2+\operatorname{tr}\bigl(C_\Theta(F)\bigr).
\)
By Proposition \ref{prop:Uniform bound for the mean}, the mean term is bounded by \(M_m^2\). For the trace term,
\[
\operatorname{tr}\bigl(C_\Theta(F)\bigr)
=\sigma_0^2\operatorname{tr}(\Pi_{S_d})+\sum_{k=1}^r\operatorname{tr}\!\bigl(c_{k,\Theta}(F)\otimes c_{k,\Theta}(F)\bigr)
=d\sigma_0^2+\sum_{k=1}^r\|c_{k,\Theta}(F)\|_H^2
\leqslant d\sigma_0^2+\sum_{k=1}^rM_{C,k}^2,
\]
where the first equality uses linearity of the trace, the second uses \(\operatorname{tr}(\Pi_{S_d})=d\) and \(\operatorname{tr}(c\otimes c)=\|c\|_H^2\), and the final inequality uses Proposition \ref{prop:Uniform bounds for covariance factors}
\end{proof}

Combining the contraction of the shifted coordinates with the uniform bound on the newly generated coordinate gives the required Lyapunov drift.

\begin{proposition}[Geometric Foster--Lyapunov drift]\label{prop:Geometric Foster--Lyapunov drift}
For every \(F\in E\), \(P_\Theta V(F)\leqslant\beta^{-1}V(F)+b,
\)
where
\(
b=1-\beta^{-1}+\beta^{n-1}M_2<\infty.
\)
\end{proposition}

\begin{proof}
Fix \(F=(f_0,\ldots,f_{n-1})\). Under one transition, the next context is
\(
F'=(f_1,\ldots,f_{n-1},g)\),
where
\(
g\sim Q_\Theta(F,\cdot).
\)
Therefore
\[
\begin{aligned}
P_\Theta V(F)
&=
\int_E V(G)\,P_\Theta(F,\mathrm dG)\\
&=
\int_H V(f_1,\ldots,f_{n-1},g)\,Q_\Theta(F,\mathrm dg)\\
&=1+\sum_{j=1}^{n-1}\beta^{j-1}\|f_j\|_H^2
+\beta^{n-1}\int_H\|g\|_H^2Q_\Theta(F,\mathrm dg)\\
&\leqslant1+\beta^{-1}\sum_{j=1}^{n-1}\beta^j\|f_j\|_H^2
+\beta^{n-1}M_2\\
&\leqslant1+\beta^{-1}\|F\|_\beta^2+\beta^{n-1}M_2\\
&=\beta^{-1}\bigl(1+\|F\|_\beta^2\bigr)
+\bigl(1-\beta^{-1}\bigr)+\beta^{n-1}M_2\\
&=\beta^{-1}V(F)+b.
\end{aligned}
\]
The first inequality uses Lemma~\ref{lem:Uniform conditional second moment}. The second uses the deterministic shift contraction.
\end{proof}

The factor \(\beta^{-1}\) in the drift estimate comes from the deterministic removal of the oldest context component, while the newly generated component contributes the bounded constant \(M_2\). The drift estimate controls the growth of the chain but does not by itself ensure that trajectories starting from different contexts can mix. For this, we use the nondegenerate Gaussian output. Because all conditional laws have the same support \(S_d\), uniformly bounded means, and uniformly bounded covariance eigenvalues, they share a common amount of probability mass on any fixed ball in \(S_d\).

\begin{proposition}[Uniform minorisation of the output kernel]\label{prop:Uniform minorisation of the output kernel}
Let us fix a constant \(R_0>0\), and let \(B_0=\{y\in S_d:\|y\|_H\leqslant R_0\}\) be the closed ball of radius \(R_0\) centred at the zero function in the approximation space \(S_d\). Define \(\nu\in\mathcal P(H)\) by \(\nu(A)=\lambda_{S_d}(A\cap B_0)/\lambda_{S_d}(B_0)\) for every \(A\in\mathcal B(H)\), where \(\lambda_{S_d}\) denotes the \(d\)-dimensional Lebesgue measure on \(S_d\). Let \(c_0=(2\pi)^{-d/2}\overline\lambda^{-d/2}\exp\!\left(-(R_0+M_m)^2/(2\underline\lambda)\right)\) and \(\varepsilon_0=c_0\lambda_{S_d}(B_0)\). Then \(\varepsilon_0\in(0,1]\) and, for every \(F\in E\) and \(A\in\mathcal B(H)\), \(Q_\Theta(F,A)\geqslant\varepsilon_0\nu(A)\).
\end{proposition}

\begin{proof}
The measure \(\nu\) is a normalised Lebesgue measure on the finite-dimensional ball \(B_0\), viewed as a probability measure on \(H\) supported on \(S_d\). Proposition~\ref{prop:Admissibility of the covariance operator} gives that
\(
\|m_\Theta(F)\|_H\leqslant M_m
\) and that
\(
\underline\lambda I_{S_d}\preceq C_\Theta(F)|_{S_d}
\preceq\overline\lambda I_{S_d}.
\)
Lemma~\ref{lem:Uniform lower bound for Gaussian densities on a fixed ball} therefore gives
\(
q_F(y)\geqslant c_0\) for \( y\in B_0,
\)
where \(q_F\) is the density of \(Q_\Theta(F,\cdot)\) relative to \(\lambda_{S_d}\). Hence
\[
\begin{aligned}
Q_\Theta(F,A)
&=\int_{A\cap S_d}q_F(y)\lambda_{S_d}(\mathrm dy)\\
&\geqslant\int_{A\cap B_0}q_F(y)\lambda_{S_d}(\mathrm dy)\\
&\geqslant c_0\lambda_{S_d}(A\cap B_0)\\
&=\varepsilon_0\nu(A).
\end{aligned}
\]
Integrating the inequality over \(B_0\) also shows \(\varepsilon_0\leqslant Q_\Theta(F,B_0)\leqslant1\).
\end{proof}

The preceding proposition is a one-step minorisation of the newly generated function. It is not a one-step minorisation of the full context kernel, since the first \(n-1\) coordinates of \(P_\Theta(F,\cdot)\) are deterministic and depend on \(F\). However, after \(n\) transitions, every coordinate of the context has been generated by the model. Repeatedly applying the same one-step output minorisation at each generation step therefore produces a common product component in the \(n\)-step transition law.

\begin{proposition}[Global \(n\)-step minorisation of the context kernel]\label{prop:Global n-step minorisation}
Let \(\nu\) and \(\varepsilon_0\) be as in Proposition~\ref{prop:Uniform minorisation of the output kernel}. Let \(\nu^{\otimes n}\in\mathcal P(E)\) denote the \(n\)-fold product measure of \(\nu\), supported on \(S_d^n\subset E\), and defined for \(A\in\mathcal B(E)\) by
\(
\nu^{\otimes n}(A)
=
\int_H\cdots\int_H
\mathbf 1_A(g_1,\ldots,g_n)\,
\nu(\mathrm dg_1)\cdots\nu(\mathrm dg_n).
\)
Then, for every \(F\in E\) and \(A\in\mathcal B(E)\),
\(
P_\Theta^n(F,A)
\geqslant
\varepsilon_0^n\nu^{\otimes n}(A).
\)
\end{proposition}

\begin{proof}
We make the sequential conditioning explicit. Starting from \(F\), let \(g_1\) be the first generated function, let \(F^{(1)}\) be the resulting context, and continue recursively. After \(n\) steps all original context coordinates have been removed, and the final context is exactly
\(
(g_1,\ldots,g_n).
\)
Consequently,
\[
P_\Theta^n(F,A)
=\int_{H^n}\mathbf 1_A(g_1,\ldots,g_n)
Q_\Theta(F,\mathrm dg_1)
Q_\Theta(F^{(1)},\mathrm dg_2)\cdots Q_\Theta(F^{(n-1)},\mathrm dg_n),
\]
where \(F^{(j)}\) is a measurable function of \(F,g_1,\ldots,g_j\).

We now apply the output minorisation repeatedly, beginning with the innermost integral. For fixed \(g_1,\ldots,g_{n-1}\),
\[
\int_H\mathbf 1_A(g_1,\ldots,g_n)
Q_\Theta(F^{(n-1)},\mathrm dg_n)
\geqslant
\varepsilon_0\int_H\mathbf 1_A(g_1,\ldots,g_n)\nu(\mathrm dg_n).
\]
The right-hand side is a non-negative measurable function of \(g_1,\ldots,g_{n-1}\). Applying the same argument to the preceding conditional kernel and iterating gives
\[
P_\Theta^n(F,A)
\geqslant\varepsilon_0^n
\int_{H^n}\mathbf 1_A(g_1,\ldots,g_n)
\nu(\mathrm dg_1)\cdots\nu(\mathrm dg_n),
\]
which is \(\varepsilon_0^n\nu^{\otimes n}(A)\).
\end{proof}

The product measure appears because at each conditioning step the same context-independent measure \(\nu\) replaces the conditional output law. The proof would not produce a fixed product measure if the support or the minorising measure varied with the context.

The preceding proposition is the key mixing property. It shows that the entire state space is a small set for the \(n\)-step kernel. Consequently, two arbitrary initial distributions share a common component after every block of \(n\) transitions, which yields a strict contraction in total variation.

We can now combine the global \(n\)-step minorisation with the drift estimate. The minorisation gives existence, uniqueness, and geometric convergence of the invariant law, while the drift inequality provides the required moment control. The resulting theorem therefore describes both the distributional stability and the integrability of the long-run inference process.

\begin{theorem}[Uniform geometric ergodicity of the functional autoregressive transformer]\label{thm:Uniform geometric ergodicity}
Let \(P_\Theta\) be the transition kernel of the functional autoregressive transformer. Then the following statements hold.

\begin{enumerate}
    \item There exists a unique invariant probability measure \(\pi_\Theta\in\mathcal P(E)\) for \(P_\Theta\).

    \item The invariant probability measure is supported on \(S_d^n\).

    \item For every initial probability measure \(\mu\in\mathcal P(E)\) and every \(t\in\mathbb N_0\),
    \(
    d_{\mathrm{TV}}\bigl(\mu P_\Theta^t,\pi_\Theta\bigr)
    \leqslant
    (1-\varepsilon_0^n)^{\lfloor t/n\rfloor}.
    \)
    In particular, \(P_\Theta\) is uniformly geometrically ergodic in total variation. If \(\varepsilon_0<1\), one may take \(q_{\mathrm{TV}}=(1-\varepsilon_0^n)^{1/n}\) and \(M_{\mathrm{TV}}=(1-\varepsilon_0^n)^{-1}\). If \(\varepsilon_0=1\), convergence is exact after at most \(n\) steps.

    \item The invariant probability measure has finite \(V\)-moment and satisfies
    \[
    \int_E V(F)\,\pi_\Theta(\mathrm dF)
    \leqslant
    \frac{b}{1-\beta^{-1}}.
    \]
\end{enumerate}
\end{theorem}

\begin{proof}
The proof is done in a few steps.

\textit{Step 1: Contraction of the \(n\)-step kernel.}
Proposition~\ref{prop:Global n-step minorisation} gives \(P_\Theta^n(F,A)\geqslant\varepsilon_0^n\nu^{\otimes n}(A)\) for every \(F\in E\) and \(A\in\mathcal B(E)\). Applying Lemma~\ref{lem:Doeblin contraction in total variation} to the Markov kernel \(P_\Theta^n\), with minorisation constant \(\varepsilon_0^n\) and reference probability measure \(\nu^{\otimes n}\), yields
\(
d_{\mathrm{TV}}\bigl(\mu_1P_\Theta^n,\mu_2P_\Theta^n\bigr)
\leqslant
(1-\varepsilon_0^n)d_{\mathrm{TV}}(\mu_1,\mu_2)
\)
for all \(\mu_1,\mu_2\in\mathcal P(E)\).

Hence the map \(\mu\mapsto\mu P_\Theta^n\) is a strict contraction on \(\mathcal P(E)\) equipped with total variation distance. The space \(\mathcal P(E)\) is complete in total variation, since it is a closed subset of the Banach space of finite signed measures equipped with the total-variation norm. The Banach fixed-point theorem therefore gives a unique probability measure \(\pi_\Theta\in\mathcal P(E)\) satisfying \(\pi_\Theta P_\Theta^n=\pi_\Theta\).

\textit{Step 2: Invariance and uniqueness for the one-step kernel.}
The measure \(\pi_\Theta P_\Theta\) is also invariant under \(P_\Theta^n\), since
\(
(\pi_\Theta P_\Theta)P_\Theta^n
=
\pi_\Theta P_\Theta^{n+1}
=
(\pi_\Theta P_\Theta^n)P_\Theta
=
\pi_\Theta P_\Theta.
\)
By uniqueness of the \(P_\Theta^n\)-invariant probability measure, it follows that \(\pi_\Theta P_\Theta=\pi_\Theta\). Thus \(\pi_\Theta\) is invariant under \(P_\Theta\).

Conversely, every \(P_\Theta\)-invariant probability measure is also \(P_\Theta^n\)-invariant. Since the latter has a unique invariant probability measure, \(\pi_\Theta\) is the unique invariant probability measure for \(P_\Theta\).

\textit{Step 3: Convergence for all times.}
For \(t\in\mathbb N_0\), write \(t=qn+s\), where \(q=\lfloor t/n\rfloor\) and \(0\leqslant s<n\). Then \(\mu P_\Theta^t=(\mu P_\Theta^s)P_\Theta^{qn}\). Repeated application of the \(n\)-step contraction gives
\[
\begin{aligned}
d_{\mathrm{TV}}\bigl(\mu P_\Theta^t,\pi_\Theta\bigr)
&=
d_{\mathrm{TV}}\bigl((\mu P_\Theta^s)P_\Theta^{qn},
(\pi_\Theta P_\Theta^s)P_\Theta^{qn}\bigr)\\
&\leqslant
(1-\varepsilon_0^n)^q
d_{\mathrm{TV}}\bigl(\mu P_\Theta^s,\pi_\Theta P_\Theta^s\bigr)\\
&\leqslant
(1-\varepsilon_0^n)^q\\
&=
(1-\varepsilon_0^n)^{\lfloor t/n\rfloor},
\end{aligned}
\]
where the second inequality uses the non-expansiveness of Markov kernels in total variation and the invariance \(\pi_\Theta P_\Theta^s=\pi_\Theta\).

If \(\varepsilon_0<1\), define \(q_{\mathrm{TV}}=(1-\varepsilon_0^n)^{1/n}\). Since \(\lfloor t/n\rfloor\geqslant t/n-1\),
\(
(1-\varepsilon_0^n)^{\lfloor t/n\rfloor}
\leqslant
(1-\varepsilon_0^n)^{-1}q_{\mathrm{TV}}^t.
\)
Thus the definition of uniform geometric ergodicity, Definition \ref{def:Uniform geometric ergodicity}, holds with \(M_{\mathrm{TV}}=(1-\varepsilon_0^n)^{-1}\). If \(\varepsilon_0=1\), the \(n\)-step contraction factor is zero, so \(\mu P_\Theta^n=\pi_\Theta\) for every \(\mu\in\mathcal P(E)\), and therefore convergence is exact for every \(t\geqslant n\).

\textit{Step 4: Support of the invariant probability measure.}
After \(n\) transitions, every coordinate of the context is a generated output and therefore belongs to \(S_d\). Hence \(P_\Theta^n(F,S_d^n)=1\) for every \(F\in E\). Using invariance under \(P_\Theta^n\),
\(
\pi_\Theta(S_d^n)
=
(\pi_\Theta P_\Theta^n)(S_d^n)
=
1.
\)
Thus \(\pi_\Theta\) is supported on \(S_d^n\).

\textit{Step 5: Invariant \(V\)-moment.}
Choose the initial probability measure \(\mu_0=\nu^{\otimes n}\). Since \(\nu\) is supported on the bounded set \(B_0\subset S_d\), the product measure \(\mu_0\) is supported on the bounded set \(B_0^n\subset S_d^n\), and hence \(\mu_0(V)<\infty\). Iterating Proposition~\ref{prop:Geometric Foster--Lyapunov drift} gives
\[
\mu_0P_\Theta^t(V)
\leqslant
\beta^{-t}\mu_0(V)
+
b\sum_{j=0}^{t-1}\beta^{-j}
\leqslant
\beta^{-t}\mu_0(V)
+
\frac{b}{1-\beta^{-1}}.
\]
The total-variation convergence established above implies weak convergence \(\mu_0P_\Theta^t\Rightarrow\pi_\Theta\). Since \(V\) is continuous and non-negative, the Portmanteau theorem yields
\[
\pi_\Theta(V)
\leqslant
\liminf_{t\to\infty}\mu_0P_\Theta^t(V)
\leqslant
\frac{b}{1-\beta^{-1}}.
\]
This proves the claimed finite \(V\)-moment bound.
\end{proof}

\begin{remark}[Roles of Doeblin minorisation and drift]\label{rem:Roles of Doeblin minorisation and drift}
The global \(n\)-step Doeblin condition alone already supplies existence, uniqueness, and uniform total-variation convergence. The drift inequality is not needed for those three conclusions in the present globally minorised model. Its additional role is to quantify moment stability and to prove that the invariant law has a finite weighted second moment. It also identifies the part of the argument that would remain useful in future variants with only local small sets, where a full Harris theorem would be required.
\end{remark}

\begin{remark}[Size of the minorisation constant]\label{rem:Size of the minorisation constant}
The result is qualitative rather than a claim of a practically sharp mixing rate. The lower bound contains
\(
\exp\left[-{(R_0+M_m)^2}/{(2\sigma_0^2})\right]
\)
and is then raised to the power \(n\). It may therefore be extremely small when \(\sigma_0\) is small, the approximation dimension is large, or the context is long. The theorem certifies geometric ergodicity but does not imply that the displayed rate is numerically useful.
\end{remark}

\begin{remark}[Inference ergodicity is separate from pretraining dependence]\label{rem:Inference ergodicity separate from pretraining}
The statements in Theorem~\ref{thm:Uniform geometric ergodicity} concern the stochastic chain generated recursively by a trained functional autoregressive transformer. They are not assumptions on the data used to train the model. The fixed-parameter generalisation results in Section~\ref{sec:Fixed-parameter pretraining generalisation} instead impose dependence conditions directly on the training observations. Inference-chain geometric ergodicity is not logically required for those results.
\end{remark}

\subsection{Modelling trade-offs}\label{subsec:Interpretation and modelling trade-offs}

The main modelling restrictions of the functional autoregressive transformer arise from the structure of its output space, conditional distribution, and finite context.

\paragraph{Fixed approximation space.}
All generated functions lie in the fixed finite-dimensional subspace \(S_d\subset H\). This introduces a genuine approximation--resolution trade-off. If dynamically important modes of the target process have non-negligible components in \(S_d^\perp\), then these components cannot be recovered by any choice of transformer parameters, and an irreducible projection error remains. Increasing the dimension \(d\) enlarges the class of functions that can be represented and therefore improves the potential approximation fidelity. However, it also increases the dimension of the conditional Gaussian density and may substantially deteriorate the explicit minorisation constant obtained in the ergodicity analysis. The choice of \(S_d\) must therefore balance representational accuracy against the quantitative strength of the resulting probabilistic bounds.

\paragraph{Nondegenerate covariance floor.}
The covariance operator contains the fixed component
\(
\sigma_0^2\Pi_{S_d}\), \(\sigma_0>0\),
which ensures nondegeneracy in every direction of \(S_d\). This assumption is central to the common-support and minorisation arguments, since it provides a uniform lower spectral bound for the conditional covariance. However, it prevents the conditional distribution from becoming exactly deterministic. For deterministic PDE dynamics, for example, the true conditional law may be a Dirac measure, whereas the model necessarily retains positive variance in every retained mode. Choosing \(\sigma_0\) small allows the conditional law to approximate deterministic transitions closely over finite prediction horizons, but the corresponding Gaussian density lower bound becomes increasingly small. Consequently, improved approximation to deterministic dynamics comes at the cost of a weaker quantitative ergodicity certificate.

\paragraph{Gaussian conditional output law.}
For every context, the conditional distribution of the next function is represented by a single Gaussian measure. This provides a tractable probabilistic model with explicit moments, covariance structure, densities on \(S_d\), and Wasserstein formulae. It nevertheless restricts the class of transition laws that can be represented exactly. In particular, a single Gaussian cannot describe strongly multimodal conditional distributions or distinct transition regimes without averaging them into a unimodal law. Mixtures of Gaussian measures or more expressive generative distributions could address this limitation, but their well-posedness, common-support properties, and minorisation theory would require a separate analysis.

\paragraph{Finite context length.}
The model conditions each prediction on the preceding \(n\) functions, so dependencies beyond this window must be encoded through the current context. Increasing \(n\) allows longer temporal dependence but enlarges the state space. It also weakens the explicit ergodicity bound, since context-level minorisation is obtained only after \(n\) steps and involves the \(n\)-fold product of the one-output minorisation constant.

\paragraph{Finite-horizon prediction and asymptotic inference.}
There is a distinction between finite-horizon predictive behaviour and the asymptotic behaviour of recursive stochastic inference. The conditional mean can approximate periodic, conservative, oscillatory, or weakly dissipative dynamics over the time horizons relevant to prediction. Nevertheless, the repeated injection of nondegenerate Gaussian noise on the fixed space \(S_d\), together with the bounded conditional parameters, yields a uniformly geometrically ergodic context chain.

\section{Fixed-parameter pretraining generalisation bounds}
\label{sec:Fixed-parameter pretraining generalisation}

This section studies pretraining generalisation bounds for a fixed model parameter value \(\Theta\), the aim is to quantify how accurately the transition kernel \(P_\Theta\) approximates a reference transition kernel \(P^*\) when these observations may be dependent.

The dependence assumptions in this section concern the observations used during pretraining. They are logically separate from the recursive-inference chain studied in Section~\ref{sec:General state-space Markov chain framework for functional autoregressive transformers}; see Remark~\ref{rem:Inference ergodicity separate from pretraining}. In particular, the pretraining observations need not be generated by \(P_\Theta\), and the assumptions used to establish geometric ergodicity of the learned inference chain do not enter the generalisation theorem.

The analysis proceeds in four stages. We first define a transition-kernel loss and show that, because the shift part of the autoregressive transition is deterministic, this loss reduces exactly to a Wasserstein discrepancy between the reference and learned conditional output laws. We then derive a generalisation bound and a sample-complexity bound using Marton couplings. Next, we develop a complementary Wasserstein metric approach based on transportation inequalities and Wasserstein dependence. Finally, we compare the sampling regimes under which the two concentration frameworks are most naturally applied.

\subsection{Reference transition kernel, loss, and risk}
\label{subsec:Reference transition kernel loss and risk}

The statistical object to be learned is the one-step transition mechanism of the context process. We therefore begin by defining a reference context-transition kernel \(P^*\), constructed from the reference conditional law \(Q^*\) in exactly the same way that the model transition kernel \(P_\Theta\) is constructed from \(Q_\Theta\).

Let \(\mathsf S=E\times H\), where \(E=H^n\), and let \(S_i=(F_i,Y_i)\in\mathsf S\) denote the \(i\)-th pretraining observation. Since \(E\) and \(H\) are separable Hilbert spaces, \(\mathsf S\) is Polish. We assume that there is a reference conditional probability kernel \(Q^*:E\times\mathcal B(H)\to[0,1]\) such that \(Q^*(F_i,\cdot)\in\mathcal P_2(H)\) for every \(F_i\in E\), and such that, for every \(A\in\mathcal B(H)\), \(Q^*(F_i,A)=\mathbb P(Y_i\in A\mid F_i)\) almost surely. The kernel \(Q^*\) represents the reference conditional law of the next function given the context \(F_i\) and need not coincide with the model conditional law \(Q_\Theta\).

\begin{definition}[Reference context-transition kernel]
\label{def:Reference context transition kernel}
Let \(F=(f_0,\ldots,f_{n-1})\in E\). The reference context-transition kernel \(P^*:E\times\mathcal B(E)\to[0,1]\) induced by \(Q^*\) is defined by
\[
P^*(F,A)
=
\int_H
\mathbf 1_A(f_1,\ldots,f_{n-1},g)\,
Q^*(F,\mathrm dg),
\qquad
A\in\mathcal B(E).
\]
\end{definition}

Thus \(P^*(F,\cdot)\) and \(P_\Theta(F,\cdot)\) have exactly the same deterministic shift in their first \(n-1\) coordinates. They differ only through the conditional law of the newly appended function. The definition of \(P^*\) does not assert that the complete pretraining sequence \((S_i)_{i=1}^N\) is a Markov chain, the dependence between the observations \((S_i)_{i=1}^N\) is treated separately through the assumptions introduced later.
Next, to define the loss, we use the weighted context norm from Definition~\ref{def:Weighted context norm},
\(
\|F\|_\beta^2
=
\sum_{j=0}^{n-1}\beta^j\|f_j\|_H^2,\;
\beta>1,
\)
and denote by \(W_{2,\beta}\) the corresponding \(2\)-Wasserstein distance on \(\mathcal P_2(E)\).

\begin{definition}[Transition loss]
\label{def:Fixed parameter transition loss}
For a model parameter value \(\Theta\), the transition loss at context \(F\in E\) is
\(
\ell_\Theta(F)
=
W_{2,\beta}^2
\bigl(
P^*(F,\cdot),
P_\Theta(F,\cdot)
\bigr).
\)
For a pretraining observation \(S=(F,Y)\in\mathsf S\), we set \(\ell_\Theta(S)=\ell_\Theta(F)\).
\end{definition}

At first sight, Definition~\ref{def:Fixed parameter transition loss} appears to require optimal transport on the product space \(E\). The autoregressive structure makes the situation much simpler. The deterministic shift is identical under \(P^*\) and \(P_\Theta\), so all transport cost is concentrated in the newly appended coordinate. The next proposition makes this reduction exact.

\begin{proposition}[Reduction of the context-transition loss]
\label{prop:Reduction of context transition loss}
For every \(F\in E\),
\[
W_{2,\beta}^2
\bigl(
P^*(F,\cdot),
P_\Theta(F,\cdot)
\bigr)
=
\beta^{n-1}
W_2^2
\bigl(
Q^*(F,\cdot),
Q_\Theta(F,\cdot)
\bigr).
\]
Consequently,
\(
\ell_\Theta(F)
=
\beta^{n-1}
W_2^2
\bigl(
Q^*(F,\cdot),
Q_\Theta(F,\cdot)
\bigr).
\)
\end{proposition}

\begin{proof}
Fix \(F=(f_0,\ldots,f_{n-1})\in E\), and define \(T_F:H\to E\) by
\(
T_F(g)=(f_1,\ldots,f_{n-1},g).
\)
By construction,
\(
P^*(F,\cdot)=(T_F)_\#Q^*(F,\cdot)\) and 
\(
P_\Theta(F,\cdot)=(T_F)_\#Q_\Theta(F,\cdot).
\)

Let \(\pi\) be any coupling of \(Q^*(F,\cdot)\) and \(Q_\Theta(F,\cdot)\). Pushing \(\pi\) forward through \((T_F,T_F)\) gives a coupling of \(P^*(F,\cdot)\) and \(P_\Theta(F,\cdot)\). For \(g,h\in H\),
\(
\|T_F(g)-T_F(h)\|_\beta^2
=
\beta^{n-1}\|g-h\|_H^2,
\)
because the first \(n-1\) coordinates of the two contexts are identical. Hence
\[
W_{2,\beta}^2
\bigl(
P^*(F,\cdot),
P_\Theta(F,\cdot)
\bigr)
\leqslant
\beta^{n-1}
W_2^2
\bigl(
Q^*(F,\cdot),
Q_\Theta(F,\cdot)
\bigr).
\]
For the reverse inequality, let \(\widetilde\pi\) be any coupling of \(P^*(F,\cdot)\) and \(P_\Theta(F,\cdot)\). Both marginals are supported on contexts whose first \(n-1\) coordinates are \((f_1,\ldots,f_{n-1})\). Projecting \(\widetilde\pi\) onto the final coordinates therefore gives a coupling \(\pi\) of \(Q^*(F,\cdot)\) and \(Q_\Theta(F,\cdot)\), and
\[
\int_{E\times E}\|F'-G'\|_\beta^2\,\widetilde\pi(\mathrm dF',\mathrm dG')
=
\beta^{n-1}
\int_{H\times H}\|g-h\|_H^2\,\pi(\mathrm dg,\mathrm dh).
\]
Taking the infimum over all couplings gives the reverse inequality and hence equality.
\end{proof}

The proposition shows that learning the context transition kernel is equivalent, for this loss, to learning the conditional distribution of the newly appended function. The shift component of the transition is known exactly and contributes no approximation error.
Next, we recall the definitions of true and empirical risks.

\begin{definition}[Risk]
\label{def:Empirical and population transition risks}
For a model parameter \(\Theta\) and pretraining observations \(S_1,\ldots,S_N\), define the empirical risk by
\(
\widehat{\mathcal R}_N(\Theta)
=
N^{-1}\sum_{i=1}^N\ell_\Theta(S_i).
\) and the true risk by
\(
\mathcal R_N(\Theta)
=
\mathbb E
\bigl[
\widehat{\mathcal R}_N(\Theta)
\bigr].
\) The generalisation problem consists of bounding \(\ \mathcal R_N(\Theta) - \widehat{\mathcal R}_N(\Theta) \).
\end{definition}

To obtain a concentration inequality from these definitions, the loss must first be uniformly controlled. The model already has a uniform conditional second-moment bound from the inference analysis. We impose the analogous condition on the reference conditional law.

\begin{assumption}[Uniform second moment of the reference conditional law]
\label{ass:Uniform second moment reference law}
There exists \(M_2^*<\infty\) such that
\[
\sup_{F\in E}
\int_H
\|y\|_H^2
Q^*(F,\mathrm dy)
\leqslant
M_2^*.
\]
\end{assumption}

This assumption requires a uniform bound on the conditional energy of the next function. In applications, the assumption is most natural when the admissible training contexts are generated from a bounded-energy family of initial conditions, or when the reference dynamics themselves possess a uniform moment bound. A simple deterministic example is provided by the heat equation.

\begin{example}[Heat-equation data with bounded initial energy]
\label{ex:Heat equation uniform second moment}
Let \(S_t:H\to H\) denote the heat semigroup on \(H=L^2([0,1])\), and suppose that the data are generated from initial conditions satisfying \(\|u_0\|_H\leqslant R_*\) almost surely for some \(R_*<\infty\). Assume that each context \(F=(f_0,\ldots,f_{n-1})\) consists of consecutive states of the heat equation, and consider the heat equation reference conditional law \(Q^*(F,\cdot)=\delta_{S_{\Delta t}f_{n-1}}\), where \(\Delta t>0\) is the sampling interval. Then, Assumption~\ref{ass:Uniform second moment reference law} holds with \(M_2^*=R_*^2\).
\end{example}

\begin{proof}
Let \(\Pi_{R_*}^H:H\to H\) denote the metric projection onto the closed ball \(\{f\in H:\|f\|_H\leqslant R_*\}\). Since Assumption~\ref{ass:Uniform second moment reference law} is formulated on the full context space \(E\), extend the reference kernel away from the heat-equation contexts by defining \(Q^*(F,\cdot)=\delta_{S_{\Delta t}\Pi_{R_*}^H(f_{n-1})}\) for arbitrary \(F=(f_0,\ldots,f_{n-1})\in E\). For a context arising from the stated heat-equation data, the \(L^2\)-contractivity of the heat semigroup gives \(\|f_{n-1}\|_H\leqslant\|u_0\|_H\leqslant R_*\), and hence \(\Pi_{R_*}^H(f_{n-1})=f_{n-1}\). Thus this extension agrees with the reference conditional law specified in the statement of the example on every context generated by the data.

For arbitrary \(F\in E\), since \(Q^*(F,\cdot)\) is a Dirac measure,
\[\int_H\|y\|_H^2Q^*(F,\mathrm dy)=\|S_{\Delta t}\Pi_{R_*}^H(f_{n-1})\|_H^2\leqslant\|\Pi_{R_*}^H(f_{n-1})\|_H^2\leqslant R_*^2,\]
where the first inequality again uses the \(L^2\)-contractivity of the heat semigroup and the second follows from the definition of the metric projection. Taking the supremum over \(F\in E\) gives the result. Hence Assumption~\ref{ass:Uniform second moment reference law} holds with \(M_2^*=R_*^2\).
\end{proof}

Assumption \ref{ass:Uniform second moment reference law} together with the corresponding model bound, is sufficient to obtain a finite deterministic upper bound on the Wasserstein transition loss.

\begin{proposition}[Uniform boundedness of the transition loss]
\label{prop:Uniform boundedness of fixed-parameter loss}
Under Assumption~\ref{ass:Uniform second moment reference law}, the transition loss satisfies \(0\leqslant\ell_\Theta(F)\leqslant B_\Theta\) for every \(F\in E\), where
\(
B_\Theta
=
\beta^{n-1}
\left(
\sqrt{M_2^*}+\sqrt{M_2}
\right)^2,
\)
and \(M_2\) is the model second-moment constant from Lemma~\ref{lem:Uniform conditional second moment}.
\end{proposition}

\begin{proof}
Fix \(F\in E\) and let \(\delta_0\) denote the Dirac measure at the zero function. Then, we compute:
\[
\begin{aligned}
\ell_\Theta(F)
&=
\beta^{n-1}
W_2^2
\bigl(
Q^*(F,\cdot),
Q_\Theta(F,\cdot)
\bigr)
\\
&\leqslant
\beta^{n-1}
\left[
W_2
\bigl(
Q^*(F,\cdot),
\delta_0
\bigr)
+
W_2
\bigl(
\delta_0,
Q_\Theta(F,\cdot)
\bigr)
\right]^2
\\
&=
\beta^{n-1}
\left[
\left(
\int_H \|y\|_H^2\,Q^*(F,\mathrm dy)
\right)^{1/2}
+
\left(
\int_H \|y\|_H^2\,Q_\Theta(F,\mathrm dy)
\right)^{1/2}
\right]^2
\\
&\leqslant
\beta^{n-1}
\left(
\sqrt{M_2^*}+\sqrt{M_2}
\right)^2
\\
&=:
B_\Theta,
\end{aligned}
\]
The first equality follows from Proposition~\ref{prop:Reduction of context transition loss}. The first inequality follows from the triangle inequality for \(W_2\), applied with the intermediate measure \(\delta_0\). The second equality follows from Example~\ref{ex:Dirac optimal transport}. The second inequality follows from Assumption~\ref{ass:Uniform second moment reference law} for \(Q^*\) and Lemma~\ref{lem:Uniform conditional second moment} for \(Q_\Theta\). The final equality is the definition of \(B_\Theta\).
\end{proof}

Proposition \ref{prop:Uniform boundedness of fixed-parameter loss} provides the deterministic bound required by both generalisation arguments developed below. We first derive a concentration and sample-complexity result using Marton couplings. We then develop a metric alternative based on transportation inequalities and Wasserstein dependence. The sampling regimes under which each framework is appropriate are discussed jointly at the end of the section.

\subsection{Marton-coupling generalisation bounds and sample complexity}
\label{subsec:Marton coupling generalisation}

The empirical risk is an average of bounded losses, but the pretraining observations need not be independent. We therefore require a dependence notion that is weak enough to include several practically relevant sampling schemes while still supporting concentration inequalities. Marton coupling provides such a framework by measuring how a perturbation of one coordinate can propagate to later coordinates.

We follow the Hamming-coupling route used structurally by Zekri et al.~\cite{ZekriEtAl2025}. Since the observation space \(\mathsf S\) is Polish, regular conditional distributions exist and the Marton construction applies without discretising the functional observations.
The first definition describes how two versions of the future sequence are coupled after changing one observation in the past.

\begin{definition}[Marton coupling]
\label{def:Marton coupling}
Let \(\mathbf S=(S_1,\ldots,S_N)\) take values in \(\mathsf S^N\). A Marton coupling is a family of coupled pairs
\(
\left(
\mathbf S^{(s_1,\ldots,s_i,s_i')},
\mathbf S'^{(s_1,\ldots,s_i,s_i')}
\right),
\)
indexed by \(i\in[N]\) and admissible \(s_1,\ldots,s_i,s_i'\in\mathsf S\), such that the coupled sequences agree in coordinates \(1,\ldots,i-1\), have \(i\)-th coordinates fixed respectively to \(s_i\) and \(s_i'\), and their future coordinates have the corresponding regular conditional laws. If \(s_i=s_i'\), the two sequences may be chosen equal.
\end{definition}

A Marton coupling is qualitative until the probability of future disagreement is quantified. The mixing matrix records precisely this propagation of dependence: its \((i,j)\)-entry measures how strongly changing coordinate \(i\) can affect coordinate \(j\).

\begin{definition}[Marton mixing matrix]
\label{def:Marton mixing matrix}
The mixing matrix \(\Gamma\in\mathbb R^{N\times N}\) associated with a Marton coupling is the upper-triangular matrix satisfying \(\Gamma_{ii}=1\), \(\Gamma_{ij}=0\) for \(i>j\), and, for \(i<j\),
\[
\Gamma_{ij}
=
\sup_{s_1,\ldots,s_i,s_i'}
\mathbb P
\left(
S_j^{(s_1,\ldots,s_i,s_i')}
\neq
S_j'^{(s_1,\ldots,s_i,s_i')}
\right).
\]
Its Euclidean operator norm is denoted by \(\|\Gamma\|_{\mathrm{op}}\).
\end{definition}

The Marton-coupling assumption is satisfied by several standard dependence structures and is therefore non-vacuous. In the simplest case, independent observations admit a Marton coupling with mixing matrix
\(
\Gamma=I_N,
\)
so that \(\|\Gamma\|_{\mathrm{op}}=1\). Markov-chain observations also fall within the coupling-based concentration framework: Paulin derives a bounded-differences inequality for Markov chains in terms of their mixing time in \cite[Corollary~2.11]{Paulin2015}. In the present functional setting, a concrete instance is provided by independently sampled trajectories. When the initial conditions are independent and each trajectory is generated deterministically from its own initial condition, the resulting trajectory-level observations remain independent and hence have mixing matrix \(I_N\). This is verified explicitly for heat-equation trajectories in Example~\ref{ex:Heat equation Marton coupling} of Section~\ref{subsec:Training-data regimes}.
The next theorem is the concentration tool used in the pretraining analysis. We use the dependent bounded-differences theorem of Paulin \cite[Theorem~2.1]{Paulin2015}, specialised to the trivial partition.

\begin{theorem}[Dependent bounded differences under a Marton coupling]
\label{thm:Dependent bounded differences Marton}
Suppose that \(\mathbf S\) admits a Marton coupling with mixing matrix \(\Gamma\). Let \(f:\mathsf S^N\to\mathbb R\) satisfy
\(
f(s)-f(s')
\leqslant
\sum_{i=1}^N
c_i\mathbf 1_{\{s_i\neq s_i'\}}
\)
for every \(s,s'\in\mathsf S^N\), where \(c_i\geq0\). Then, for every \(t\geq0\),
\[
\mathbb P
\left(
|f(\mathbf S)-\mathbb Ef(\mathbf S)|
\geq t
\right)
\leqslant
2\exp
\left(
-\frac{2t^2}
{\|\Gamma\|_{\mathrm{op}}^2\|c\|_2^2}
\right).
\]
\end{theorem}

The theorem separates the two quantities that determine concentration. The vector \(c\) measures how sensitive the statistic itself is to each observation, while \(\Gamma\) measures how changes in one observation can propagate through the data sequence. For the empirical risk, the boundedness result above makes the first quantity explicit.
We can now apply this abstract concentration theorem to the empirical risk.

\begin{theorem}[Fixed-parameter Marton generalisation bound]
\label{thm:Fixed-parameter Marton generalisation bound}
Fix \(\Theta\), and suppose that \(\mathbf S\) admits a Marton coupling with mixing matrix \(\Gamma\). Under Assumption~\ref{ass:Uniform second moment reference law}, for every \(t\geq0\),
\[
\mathbb P
\left(
\left|
\widehat{\mathcal R}_N(\Theta)
-
\mathcal R_N(\Theta)
\right|
\geq t
\right)
\leqslant
2\exp
\left(
-\frac{2Nt^2}
{B_\Theta^2\|\Gamma\|_{\mathrm{op}}^2}
\right).
\]
Consequently, for every \(\delta\in(0,1)\), with probability at least \(1-\delta\),
\[
\mathcal R_N(\Theta)
\leqslant
\widehat{\mathcal R}_N(\Theta)
+
\frac{
B_\Theta\|\Gamma\|_{\mathrm{op}}
}{
\sqrt{2N}
}
\sqrt{
\log\left(\frac2\delta\right)
}.
\]
\end{theorem}

\begin{proof}
Define
\(
f(s_1,\ldots,s_N)
=
N^{-1}
\sum_{j=1}^N
\ell_\Theta(s_j).
\)
Let \(s=(s_1,\ldots,s_N)\) and \(s'=(s_1',\ldots,s_N')\) differ only in coordinate \(i\). Then
\(
|f(s)-f(s')|
=
N^{-1}
\left|
\ell_\Theta(s_i)-\ell_\Theta(s_i')
\right|
\leqslant{B_\Theta}/{N}
\), where the inequality follows from Proposition~\ref{prop:Uniform boundedness of fixed-parameter loss}, since both \(\ell_\Theta(s_i)\) and \(\ell_\Theta(s_i')\) lie in \([0,B_\Theta]\), and hence their difference has absolute value at most \(B_\Theta\).

Therefore the bounded-differences vector in Theorem~\ref{thm:Dependent bounded differences Marton} may be chosen as \(c=(B_\Theta/N,\ldots,B_\Theta/N)\), with norm \(\|c\|_2
=
{B_\Theta}{/\sqrt N}.
\)

Applying Theorem~\ref{thm:Dependent bounded differences Marton} now gives, for every \(t\geq0\),
\[
\begin{aligned}
\mathbb P
\left(
\left|
f(\mathbf S)-\mathbb Ef(\mathbf S)
\right|
\geq t
\right)
&\leqslant
2\exp
\left(
-\frac{2t^2}
{\|\Gamma\|_{\mathrm{op}}^2\|c\|_2^2}
\right)
\\
&=
2\exp
\left(
-\frac{2t^2}
{\|\Gamma\|_{\mathrm{op}}^2(B_\Theta^2/N)}
\right)
\\
&=
2\exp
\left(
-\frac{2Nt^2}
{B_\Theta^2\|\Gamma\|_{\mathrm{op}}^2}
\right).
\end{aligned}
\]
Noting that \(
f(\mathbf S)=\widehat{\mathcal R}_N(\Theta)\) and 
\(
\mathbb Ef(\mathbf S)=\mathcal R_N(\Theta)
\) concludes the first part.

To obtain the high-probability form, choose \(t=t_\delta\) so that the upper bound equals \(\delta\). Then
\[
\delta
=
2\exp
\left(
-\frac{2Nt_\delta^2}
{B_\Theta^2\|\Gamma\|_{\mathrm{op}}^2}
\right)
\Longleftrightarrow
\log\left(\frac2\delta\right)
=
\frac{2Nt_\delta^2}
{B_\Theta^2\|\Gamma\|_{\mathrm{op}}^2}
\Longleftrightarrow
t_\delta
=
\frac{
B_\Theta\|\Gamma\|_{\mathrm{op}}
}{
\sqrt{2N}
}
\sqrt{
\log\left(\frac2\delta\right)
}.
\]

Hence, with probability at least \(1-\delta\),
\[
\left|
\widehat{\mathcal R}_N(\Theta)
-
\mathcal R_N(\Theta)
\right|
\leqslant
\frac{
B_\Theta\|\Gamma\|_{\mathrm{op}}
}{
\sqrt{2N}
}
\sqrt{
\log\left(\frac2\delta\right)
} \implies
\mathcal R_N(\Theta)
\leqslant
\widehat{\mathcal R}_N(\Theta)
+
\frac{
B_\Theta\|\Gamma\|_{\mathrm{op}}
}{
\sqrt{2N}
}
\sqrt{
\log\left(\frac2\delta\right)
},
\]
which proves the result.
\end{proof}

The theorem gives a direct functional analogue of a dependent bounded-differences generalisation bound. Its rate is \(N^{-1/2}\), while all dependence between the observations appears through \(\|\Gamma\|_{\mathrm{op}}\). In the independent case, this factor is one and the usual bounded-differences scaling is recovered.

The preceding theorem bounds the generalisation gap for a given sample size. The complementary question is how many observations are sufficient to guarantee a prescribed population transition risk. Under the idealised assumption of zero empirical transition risk, this can be obtained directly by inverting the concentration bound.
We first formalise what is meant by pretraining sample complexity.

\begin{definition}[Pretraining sample complexity]
\label{def:Pretraining sample complexity}
For a model parameter \(\Theta\), confidence level \(\delta\in(0,1)\), and target transition risk \(\varepsilon_{\mathrm{pre}}>0\), the pretraining sample complexity is the smallest number of observations sufficient to guarantee
\(
\mathcal R_N(\Theta)
\leqslant
\varepsilon_{\mathrm{pre}}
\)
with probability at least \(1-\delta\), under the stated assumptions on the empirical risk and dependence structure.
\end{definition}

\begin{proposition}[Pretraining sample complexity under zero empirical risk]
\label{prop:Pretraining sample complexity zero empirical risk}
Fix \(\Theta\), and suppose that for every candidate sample size \(N\), the observations admit a Marton coupling with mixing matrix \(\Gamma_N\) satisfying \(\|\Gamma_N\|_{\mathrm{op}}\leqslant\Gamma_\star\) for some finite constant \(\Gamma_\star\). Let \(\varepsilon_{\mathrm{pre}}>0\) and \(\delta\in(0,1)\). If
\(
\widehat{\mathcal R}_{N_{\mathrm{train}}}(\Theta)=0
\)
and
\[
N_{\mathrm{train}}
\geq
\left\lceil
\frac{
B_\Theta^2\Gamma_\star^2
}{
2\varepsilon_{\mathrm{pre}}^2
}
\log
\left(
\frac2\delta
\right)
\right\rceil,
\]
then, with probability at least \(1-\delta\),
\(
\mathcal R_{N_{\mathrm{train}}}(\Theta)
\leqslant
\varepsilon_{\mathrm{pre}}.
\)
\end{proposition}

\begin{proof}
Theorem~\ref{thm:Fixed-parameter Marton generalisation bound}, the zero-empirical-risk assumption, and \(\|\Gamma_{N_{\mathrm{train}}}\|_{\mathrm{op}}\leqslant\Gamma_\star\) give
\[
\mathcal R_{N_{\mathrm{train}}}(\Theta)
\leqslant
\frac{
B_\Theta\Gamma_\star
}{
\sqrt{2N_{\mathrm{train}}}
}
\sqrt{
\log
\left(
\frac2\delta
\right)
}
\]
with probability at least \(1-\delta\). The stated lower bound on \(N_{\mathrm{train}}\) makes the right-hand side at most \(\varepsilon_{\mathrm{pre}}\).
\end{proof}

\begin{remark}[Nonzero empirical risk]
If instead \(\widehat{\mathcal R}_{N_{\mathrm{train}}}(\Theta)\leqslant\varepsilon_{\mathrm{train}}<\varepsilon_{\mathrm{pre}}\), the same calculation applies with required generalisation gap \(\varepsilon_{\mathrm{pre}}-\varepsilon_{\mathrm{train}}\).
\end{remark}

\subsection{Wasserstein metric generalisation bounds and sample complexity}
\label{subsec:Wasserstein and metric dependence}

The Marton-coupling result of the previous subsection measures dependence through the probability that two coupled coordinates disagree exactly. This is natural for discrete observations, but exact equality can be too rigid for function-valued data. Two trajectories may remain very close in the Hilbert norm while being different at every time point.

For example, let \(T:H\to H\) be a deterministic evolution map. If \(Tf\neq Tg\), then \(d_{\mathrm{TV}}(\delta_{Tf},\delta_{Tg})=1\), even when \(\|Tf-Tg\|_H\) is arbitrarily small. By contrast, \(W_1(\delta_{Tf},\delta_{Tg})=\|Tf-Tg\|_H\). Thus a deterministic evolution can be maximally separated in total variation while remaining stable in the natural function-space metric.

The aim of this subsection is to exploit this metric stability directly. We first impose a \(T_1\) transportation inequality on the one-step conditional laws of the data process. This controls concentration of each conditional distribution. We then impose a \(W_1\)-stability condition on the data-generating transition kernel, which controls how perturbations of the current state affect the next conditional distribution. Theorem~2.5 of Djellout, Guillin, and Wu \cite{DjelloutGuillinWu2004} combines these two properties to obtain a transportation inequality for the law of the entire dependent trajectory. This path-space inequality then yields concentration for Lipschitz empirical averages. Finally, we apply that concentration result to the transition loss.

\paragraph{One-step transportation inequality.}

We begin with the one-step concentration property used by Djellout, Guillin, and Wu. Informally, a \(T_1(C)\) inequality states that moving probability mass far from a reference distribution in \(W_1\) requires a correspondingly large relative-entropy cost. The constant \(C\) therefore determines the concentration scale of the measure.

\begin{definition}[\(T_1\) transportation inequality]
\label{def:T1 transportation inequality}
Let \((\mathsf X,d)\) be a Polish metric space. A probability measure \(\mu\in\mathcal P(\mathsf X)\) satisfies \(T_1(C)\) if, for every \(\eta\in\mathcal P(\mathsf X)\),
\(
W_1^d(\eta,\mu)
\leqslant
\sqrt{2C\,H(\eta\mid\mu)},
\)
where \(H(\eta\mid\mu)\) denotes the relative entropy of \(\eta\) with respect to \(\mu\), interpreted as \(+\infty\) when \(\eta\) is not absolutely continuous with respect to \(\mu\).
\end{definition}

Smaller values of \(C\) correspond to stronger concentration. In the dependent setting below, the same finite constant is assumed for the initial law and for every one-step conditional law. This controls the amount of randomness introduced at each step. A separate Wasserstein condition controls the propagation of perturbations between consecutive states.

\paragraph{Path space transportation inequality.}

Let \((X_i)_{i=1}^N\) be a Markov chain on \((\mathsf X,d)\), with initial distribution \(\mu_1\) and transition kernel \(K\). Denote its joint path law by \(\mathbb P_N=\mathcal L(X_1,\ldots,X_N)\), and equip \(\mathsf X^N\) with the product metric \(d_{\ell_1}(x,y)=\sum_{i=1}^N d(x_i,y_i)\).

The next proposition is the bridge between the local assumptions on the data-generating kernel and concentration of the entire dependent sequence.

\begin{proposition}[Path-space transportation inequality]
\label{prop:Path space transportation inequality}
Suppose that \(\mu_1\) and every \(K(x,\cdot)\) satisfy \(T_1(C)\) for the same finite constant \(C\). Assume further that there exists a constant \(\kappa\in[0,1)\) such that 
\(
W_1^d
\bigl(
K(x,\cdot),
K(y,\cdot)
\bigr)
\leqslant
\kappa d(x,y)
\)
for every \(x,y\in\mathsf X\).

Then the joint law \(\mathbb P_N\) satisfies \(T_1(CN/(1-\kappa)^2)\) on \((\mathsf X^N,d_{\ell_1})\). Equivalently, for every probability measure \(\mathbb Q\) on \(\mathsf X^N\),
\[
W_1^{d_{\ell_1}}
\bigl(
\mathbb Q,
\mathbb P_N
\bigr)
\leqslant
\frac{1}{1-\kappa}
\sqrt{
2CN\,H(\mathbb Q\mid\mathbb P_N)
}.
\]
\end{proposition}

\begin{proof}
For \(i\geq2\), the Markov property gives \(P_i(\cdot\mid x^{i-1})=K(x_{i-1},\cdot)\), where \(P_i(\cdot\mid x^{i-1})\) denotes the conditional law of the \(i\)-th coordinate under \(\mathbb P_N\).

Condition \emph{(C1)} of \cite[Theorem~2.5]{DjelloutGuillinWu2004} is verified by choosing \(a_1=\kappa\) and \(a_j=0\) for \(j\geq2\). Indeed,
\[
\begin{aligned}
W_1^d
\left(
P_i(\cdot\mid x^{i-1}),
P_i(\cdot\mid\widetilde x^{i-1})
\right)
&=
W_1^d
\left(
K(x_{i-1},\cdot),
K(\widetilde x_{i-1},\cdot)
\right)
\\
&\leqslant
\kappa
d(x_{i-1},\widetilde x_{i-1})
\\
&=
\sum_{j=1}^{i-1}
a_j
d(x_{i-j},\widetilde x_{i-j}).
\end{aligned}
\]
The first equality follows from the Markov property. The inequality is the assumed Wasserstein stability of \(K\). The final equality follows from the choice \(a_1=\kappa\) and \(a_j=0\) for \(j\geq2\). The dependence coefficient is therefore \(r_1=\sum_{j\geq1}a_j=\kappa<1\).

Taking \(p=1\) in \cite[Theorem~2.5]{DjelloutGuillinWu2004} gives
\[
\begin{aligned}
W_1^{d_{\ell_1}}
\bigl(
\mathbb Q,
\mathbb P_N
\bigr)
&\leqslant
\frac{1}{1-r}
\sqrt{
2CN\,H(\mathbb Q\mid\mathbb P_N)
}
\\
&=
\frac{1}{1-\kappa}
\sqrt{
2CN\,H(\mathbb Q\mid\mathbb P_N)
}
\\
&=
\sqrt{
2
\left(
\frac{CN}{(1-\kappa)^2}
\right)
H(\mathbb Q\mid\mathbb P_N)
}.
\end{aligned}
\]
Comparing the final expression with Definition~\ref{def:T1 transportation inequality} proves that \(\mathbb P_N\) satisfies \(T_1(CN/(1-\kappa)^2)\).
\end{proof}

The two hypotheses play distinct roles. The \(T_1(C)\) assumption controls concentration of each conditional distribution, whereas \(\kappa\) measures sensitivity of the next-step law to the current state. The factor \((1-\kappa)^{-1}\) records the accumulation of this dependence along the trajectory. If \(\kappa\) is small, perturbations are rapidly forgotten; as \(\kappa\) approaches one, the concentration bound deteriorates.

These assumptions are satisfied by concrete classes of stochastic dynamical systems and are therefore not merely abstract regularity conditions. For example, Djellout, Guillin, and Wu study discrete-time stochastic recursions of the form
\(
X_{k+1}=f(X_k)+\sigma(X_k)W_{k+1},
\)
for which the above assumptions are satisfied; see \cite[Section~3.2, Corollary~3.5]{DjelloutGuillinWu2004}. In the present functional setting, an explicit example is given by the heat-equation context process in Example~\ref{ex:Heat equation metric dependence} of Section~\ref{subsec:Training-data regimes}, where the \(L^2\)-contractivity of the heat semigroup provides the Wasserstein dependence coefficient and bounded initial energy yields a common finite \(T_1\) constant.

\paragraph{Concentration of empirical averages.}

Once the path law satisfies a \(T_1\) inequality, the next step is to apply it to an empirical average. Averaging is important because the Lipschitz constant of the empirical statistic is of order \(N^{-1}\), which ultimately produces the familiar \(N^{-1/2}\) concentration rate.

\begin{theorem}[Metric concentration for a Wasserstein-dependent data chain]
\label{thm:Metric concentration Wasserstein data chain}
Let \((X_i)_{i=1}^N\) satisfy the assumptions of Proposition~\ref{prop:Path space transportation inequality}. Let \(\varphi:\mathsf X\to\mathbb R\) be \(L_\varphi\)-Lipschitz, and define \(\widehat r_N=N^{-1}\sum_{i=1}^N\varphi(X_i)\) and \(r_N=\mathbb E[\widehat r_N]\). Then, for every \(t>0\),
\[
\mathbb P
\left(
|\widehat r_N-r_N|
\geq t
\right)
\leqslant
2
\exp
\left(
-\frac{
Nt^2(1-\kappa)^2
}{
2CL_\varphi^2
}
\right).
\]
Consequently, for every \(\delta\in(0,1)\), with probability at least \(1-\delta\),
\[
r_N
\leqslant
\widehat r_N
+
\frac{L_\varphi}{1-\kappa}
\sqrt{
\frac{2C}{N}
\log
\left(
\frac2\delta
\right)
}.
\]
\end{theorem}

\begin{proof}
Define \(f:\mathsf X^N\to\mathbb R\) by \(f(x_1,\ldots,x_N)=N^{-1}\sum_{i=1}^N\varphi(x_i)\). For \(x,y\in\mathsf X^N\),
\[
\begin{aligned}
|f(x)-f(y)|
&=
\left|
\frac1N
\sum_{i=1}^N
\bigl(
\varphi(x_i)-\varphi(y_i)
\bigr)
\right|
\\
&\leqslant
\frac1N
\sum_{i=1}^N
|\varphi(x_i)-\varphi(y_i)|
\\
&\leqslant
\frac{L_\varphi}{N}
\sum_{i=1}^N
d(x_i,y_i)
\\
&=
\frac{L_\varphi}{N}
d_{\ell_1}(x,y).
\end{aligned}
\]
The first inequality is the triangle inequality. The second follows from the \(L_\varphi\)-Lipschitz continuity of \(\varphi\). The final equality is the definition of \(d_{\ell_1}\). Thus \(f\) is \(L_\varphi/N\)-Lipschitz on the path space.

By Proposition~\ref{prop:Path space transportation inequality}, the path law satisfies \(T_1(CN/(1-\kappa)^2)\). Applying the concentration consequence of this transportation inequality, equivalently \cite[Corollary~2.6]{DjelloutGuillinWu2004}, gives
\[
\begin{aligned}
\mathbb P
\left(
f-\mathbb Ef
\geq t
\right)
&\leqslant
\exp
\left(
-\frac{
t^2
}{
2
\left(
CN/(1-\kappa)^2
\right)
\left(
L_\varphi/N
\right)^2
}
\right)
\\
&=
\exp
\left(
-\frac{
Nt^2(1-\kappa)^2
}{
2CL_\varphi^2
}
\right).
\end{aligned}
\]
The inequality is the concentration consequence of the path-space \(T_1\) inequality. 

Applying the same argument to \(-f\) and using a union bound gives
\[
\mathbb P
\left(
|f-\mathbb Ef|
\geq t
\right)
\leqslant
2
\exp
\left(
-\frac{
Nt^2(1-\kappa)^2
}{
2CL_\varphi^2
}
\right).
\]
Since \(f(X_1,\ldots,X_N)=\widehat r_N\) and \(\mathbb Ef=r_N\), this proves the first statement.

For the high-probability form, choose \(t=t_\delta\) such that
\[
2
\exp
\left(
-\frac{
Nt_\delta^2(1-\kappa)^2
}{
2CL_\varphi^2
}
\right)
=
\delta.
\implies
t_\delta
=
\frac{L_\varphi}{1-\kappa}
\sqrt{
\frac{2C}{N}
\log
\left(
\frac2\delta
\right)
}.
\]
\end{proof}

The theorem separates the sources of the statistical constant clearly. The factor \(N^{-1/2}\) comes from averaging \(N\) observations, \(L_\varphi\) measures sensitivity of the statistic being averaged, \(C\) measures one-step concentration, and \((1-\kappa)^{-1}\) measures persistence of dependence along the data trajectory.

\paragraph{Application to the transformer loss.}

The preceding results concern a generic Lipschitz statistic of a metric-dependent data sequence. We now specialise the statistic to the transition loss of the functional autoregressive transformer. In this part of the analysis, \(X_i=F_i\), so the observations are the functional contexts themselves. The empirical function in Theorem~\ref{thm:Metric concentration Wasserstein data chain} is therefore \(\varphi(F)=\ell_\Theta(F)\).

To apply the theorem, it remains to show that the map \(F\mapsto\ell_\Theta(F)\) is Lipschitz. Proposition~\ref{prop:Reduction of context transition loss} gives \(\ell_\Theta(F)=\beta^{n-1}W_2^2(Q^*(F,\cdot),Q_\Theta(F,\cdot))\), so regularity of the transition loss follows from Wasserstein regularity of the reference and learned conditional output laws.

\begin{proposition}[Lipschitz continuity of the transition loss]
\label{prop:Lipschitz continuity squared Wasserstein loss}
Assume that there exist \(L_*\) and \(L_\Theta\) such that
\(
W_2
\bigl(
Q^*(F,\cdot),
Q^*(G,\cdot)
\bigr)
\leqslant
L_*
\|F-G\|_\beta
\)
and
\(
W_2
\bigl(
Q_\Theta(F,\cdot),
Q_\Theta(G,\cdot)
\bigr)
\leqslant
L_\Theta
\|F-G\|_\beta
\)
for every \(F,G\in E\). Under Assumption~\ref{ass:Uniform second moment reference law},
\(
|\ell_\Theta(F)-\ell_\Theta(G)|
\leqslant
L_\ell
\|F-G\|_\beta,
\)
where
\(
L_\ell
=
2
\beta^{(n-1)/2}
\sqrt{B_\Theta}
(L_*+L_\Theta).
\)
\end{proposition}

\begin{proof}
Set \(a=W_2(Q^*(F,\cdot),Q_\Theta(F,\cdot))\) and \(b=W_2(Q^*(G,\cdot),Q_\Theta(G,\cdot))\). Then
\[
\begin{aligned}
|\ell_\Theta(F)-\ell_\Theta(G)|
&=
\beta^{n-1}
|a^2-b^2|
\\
&=
\beta^{n-1}
(a+b)|a-b|
\\
&\leqslant
\beta^{n-1}
(a+b)
\left[
W_2
\bigl(
Q^*(F,\cdot),
Q^*(G,\cdot)
\bigr)
+
W_2
\bigl(
Q_\Theta(F,\cdot),
Q_\Theta(G,\cdot)
\bigr)
\right]
\\
&\leqslant
\beta^{n-1}
(a+b)
(L_*+L_\Theta)
\|F-G\|_\beta
\\
&\leqslant
2
\beta^{(n-1)/2}
\sqrt{B_\Theta}
(L_*+L_\Theta)
\|F-G\|_\beta \\
&= L_\ell\|F-G\|_\beta.
\end{aligned}
\]
The first equality follows from Proposition~\ref{prop:Reduction of context transition loss}. The second equality uses \(a^2-b^2=(a+b)(a-b)\). The first inequality follows from the triangle inequality for \(W_2\). Indeed, inserting \(Q^*(G,\cdot)\) and \(Q_\Theta(G,\cdot)\) as intermediate measures gives an upper bound for \(a-b\); interchanging \(F\) and \(G\) gives the same bound for \(b-a\), combining the two estimates yields
\(
|a-b|
\leqslant
W_2\bigl(Q^*(F,\cdot),Q^*(G,\cdot)\bigr)
+
W_2\bigl(Q_\Theta(F,\cdot),Q_\Theta(G,\cdot)\bigr)
\). The second inequality uses the assumed Wasserstein Lipschitz bounds for \(Q^*\) and \(Q_\Theta\). Finally, Proposition~\ref{prop:Uniform boundedness of fixed-parameter loss} gives \(a,b\leqslant\beta^{-(n-1)/2}\sqrt{B_\Theta}\), yielding the final inequality.
\end{proof}

The constant \(L_*\) describes regularity of the reference conditional law, whereas \(L_\Theta\) describes regularity of the learned transformer conditional law. The latter is supplied by Proposition~\ref{prop:Lipschitz continuity of the Gaussian output kernel} in Appendix~\ref{app:Lipschitz regularity}. These constants determine sensitivity of the loss to a change in context. They are conceptually distinct from \(\kappa\), which measures temporal dependence of the data-generating process.

\begin{example}[Reference-law Lipschitz continuity for the heat equation]
Let \(S_t:H\to H\) denote the homogeneous Dirichlet heat semigroup on \(H=L^2([0,1])\), and suppose the reference next-step law is deterministic, \(Q^*(F,\cdot)=\delta_{S_{\Delta t}f_{n-1}}\). Then, the reference law is Wasserstein-Lipschitz with
\(
L_*
=
e^{-\nu\pi^2\Delta t}
\beta^{-(n-1)/2}.
\)
\end{example}

\begin{proof}
We compute
\[
\begin{aligned}
W_2
\bigl(
Q^*(F,\cdot),
Q^*(G,\cdot)
\bigr)
&=
\|S_{\Delta t}f_{n-1}-S_{\Delta t}g_{n-1}\|_H
\\
&\leqslant
e^{-\nu\pi^2\Delta t}
\|f_{n-1}-g_{n-1}\|_H
\\
&\leqslant
e^{-\nu\pi^2\Delta t}
\beta^{-(n-1)/2}
\|F-G\|_\beta.
\end{aligned}
\]
The first equality follows from the Example \ref{ex:Dirac optimal transport} Wasserstein distance between Dirac measures. The first inequality is the \(L^2\)-contractivity of the heat semigroup \(
\|S_{\Delta t}f-S_{\Delta t}g\|_H
\leqslant
e^{-\nu\pi^2\Delta t}
\|f-g\|_H,
\) and the second follows from the definition of the weighted context norm.
\end{proof}

\paragraph{Metric pretraining generalisation.}

We can now combine the data-process assumptions with the Lipschitz estimate for the transition loss. The constants \(C\) and \(\kappa\) describe the pretraining data, while \(L_\ell\) describes the statistic whose empirical average is being estimated.

\begin{theorem}[Fixed-parameter metric generalisation bound]
\label{thm:Fixed-parameter metric generalisation bound}
Fix \(\Theta\), and suppose that the context observations \((F_i)_{i=1}^N\) form a Markov chain on \((E,\|\cdot\|_\beta)\) satisfying the assumptions of Theorem~\ref{thm:Metric concentration Wasserstein data chain} with constants \(C>0\) and \(\kappa\in[0,1)\). Suppose also that Proposition~\ref{prop:Lipschitz continuity squared Wasserstein loss} holds with Lipschitz constant \(L_\ell<\infty\). Then, for every \(t\geq0\),
\[
\mathbb P
\left(
\left|
\widehat{\mathcal R}_N(\Theta)
-
\mathcal R_N(\Theta)
\right|
\geq t
\right)
\leqslant
2
\exp
\left(
-\frac{
Nt^2(1-\kappa)^2
}{
2CL_\ell^2
}
\right).
\]
Consequently, for every \(\delta\in(0,1)\), with probability at least \(1-\delta\),
\[
\mathcal R_N(\Theta)
\leqslant
\widehat{\mathcal R}_N(\Theta)
+
\frac{L_\ell}{1-\kappa}
\sqrt{
\frac{2C}{N}
\log\left(\frac2\delta\right)
}.
\]
\end{theorem}

\begin{proof}
Apply Theorem~\ref{thm:Metric concentration Wasserstein data chain} with \(\mathsf X=E\), \(d(F,G)=\|F-G\|_\beta\), and \(\varphi=\ell_\Theta\). Proposition~\ref{prop:Lipschitz continuity squared Wasserstein loss} gives \(L_\varphi=L_\ell\). Moreover, by Definition~\ref{def:Empirical and population transition risks}, \(\widehat r_N=\widehat{\mathcal R}_N(\Theta)\) and \(r_N=\mathcal R_N(\Theta)\). Substituting these identities into Theorem~\ref{thm:Metric concentration Wasserstein data chain} gives the stated two-sided concentration inequality.

For the high-probability form, choose \(t=t_\delta\) so that the right-hand side of the concentration inequality is equal to \(\delta\). Then
\[
\delta
=
2
\exp
\left(
-\frac{
Nt_\delta^2(1-\kappa)^2
}{
2CL_\ell^2
}
\right)
\Longleftrightarrow
\log\left(\frac2\delta\right)
=
\frac{
Nt_\delta^2(1-\kappa)^2
}{
2CL_\ell^2
}
\Longleftrightarrow
t_\delta
=
\frac{L_\ell}{1-\kappa}
\sqrt{
\frac{2C}{N}
\log\left(\frac2\delta\right)
}.
\]
Hence, with probability at least \(1-\delta\), we have \(|\widehat{\mathcal R}_N(\Theta)-\mathcal R_N(\Theta)|\leqslant t_\delta\). In particular, \(\mathcal R_N(\Theta)\leqslant\widehat{\mathcal R}_N(\Theta)+t_\delta\), which gives the claimed high-probability generalisation bound.
\end{proof}

The theorem has the same \(N^{-1/2}\) scaling as the Marton result, but the dependence structure is encoded differently. In the Marton theorem it appears through \(\|\Gamma\|_{\mathrm{op}}\), whereas here it appears through the one-step concentration constant \(C\) and the metric-dependence coefficient \(\kappa\).

\begin{proposition}[Metric pretraining sample complexity under zero empirical risk]
\label{prop:Metric pretraining sample complexity zero empirical risk}
Fix \(\Theta\), and suppose that the context observations satisfy the assumptions of Theorem~\ref{thm:Fixed-parameter metric generalisation bound} with constants \(C>0\), \(\kappa\in[0,1)\), and \(L_\ell<\infty\). Let \(\varepsilon_{\mathrm{pre}}>0\) and \(\delta\in(0,1)\). Assume that \(\widehat{\mathcal R}_{N_{\mathrm{train}}}(\Theta)=0\). If
\[
N_{\mathrm{train}}
\geq
\left\lceil
\frac{
2CL_\ell^2
}{
(1-\kappa)^2
\varepsilon_{\mathrm{pre}}^2
}
\log\left(\frac2\delta\right)
\right\rceil,
\]
then, with probability at least \(1-\delta\), \(\mathcal R_{N_{\mathrm{train}}}(\Theta)\leqslant\varepsilon_{\mathrm{pre}}\).
\end{proposition}

\begin{proof}
Theorem~\ref{thm:Fixed-parameter metric generalisation bound} and the zero-empirical-risk assumption give
\[
\mathcal R_{N_{\mathrm{train}}}(\Theta)
\leqslant
\frac{L_\ell}{1-\kappa}
\sqrt{
\frac{2C}{N_{\mathrm{train}}}
\log\left(\frac2\delta\right)
}
\]
with probability at least \(1-\delta\). The stated lower bound on \(N_{\mathrm{train}}\) makes the right-hand side at most \(\varepsilon_{\mathrm{pre}}\).
\end{proof}

The Marton and metric results should be interpreted as complementary concentration principles. The Marton theorem is particularly natural when the sampling design consists of independent blocks or when exact coupling probabilities are informative. The metric theorem is better adapted to a temporally dependent functional process for which nearby states remain close under the dynamics. Which theorem is more useful therefore depends on the sampling regime and the geometry of the data.

\subsection{Training-data regimes and applicability}
\label{subsec:Training-data regimes}

The Marton and metric generalisation frameworks impose different assumptions on the pretraining observations. The Marton result measures dependence through coordinatewise coupling probabilities, whereas the metric result uses the geometry of the functional state space. The appropriate theorem therefore depends both on the dynamics and on how the training observations are constructed. This distinction is particularly important for functional time series, where observations along one trajectory may be strongly dependent even when different trajectories are sampled independently \cite{HormannKokoszka2010,AueEtAl2015}.

\paragraph{Independent trajectories.}

For PDE and operator-learning datasets, a natural statistical unit is an independently generated trajectory. Suppose that \(M\) initial conditions are sampled independently and that one trajectory is generated from each. Any measurable trajectory-level observation constructed separately from these trajectories is then independent across \(m\in[M]\). Consequently, the corresponding Marton mixing matrix is
\(
\Gamma=I_M,
\)
\(
\|\Gamma\|_{\mathrm{op}}=1.
\)
The effective sample size in Theorem~\ref{thm:Fixed-parameter Marton generalisation bound} is therefore the number of independent trajectories or blocks, rather than the total number of time steps stored in the dataset.
This sampling strategy does not require temporal mixing within individual trajectories and therefore remains applicable to deterministic, periodic, or conservative dynamics.

\begin{example}[Independent heat-equation trajectories admit a Marton coupling]
\label{ex:Heat equation Marton coupling}
Let \(S_t:H\to H\) denote the heat semigroup on \(H=L^2([0,1])\). Suppose that \(M\) initial conditions \(u_0^{(1)},\ldots,u_0^{(M)}\) are sampled independently from a probability distribution \(\mu_0\) on \(H\). For each \(m\in[M]\), generate the deterministic trajectory \(u_t^{(m)}=S_tu_0^{(m)}\).
Fix a sampling time \(t_0\geq0\) and a time step \(\Delta t>0\). Define the \(m\)-th context-response observation by \(S_m=(F_m,Y_m)\in E\times H\), where
\(
F_m
=
\left(
u_{t_0}^{(m)},
u_{t_0+\Delta t}^{(m)},
\ldots,
u_{t_0+(n-1)\Delta t}^{(m)}
\right)
\)
and \(Y_m=u_{t_0+n\Delta t}^{(m)}\).
Then \(S_1,\ldots,S_M\) are independent. Consequently, they admit a Marton coupling with mixing matrix \(\Gamma=I_M\), and therefore \(\|\Gamma\|_{\mathrm{op}}=1\).
\end{example}

\begin{proof}
For each \(m\in[M]\), define the deterministic map \(\Phi:H\to E\times H\) by
\[
\Phi(u)
=
\left(
\left(
S_{t_0}u,
S_{t_0+\Delta t}u,
\ldots,
S_{t_0+(n-1)\Delta t}u
\right),
S_{t_0+n\Delta t}u
\right).
\]
Then \(S_m=\Phi(u_0^{(m)})\). Since the initial conditions \(u_0^{(1)},\ldots,u_0^{(M)}\) are independent and \(\Phi\) is deterministic, the observations \(S_1,\ldots,S_M\) are independent.

Fix \(i\in[M]\), a past \(s_1,\ldots,s_{i-1}\), and two possible values \(s_i,s_i'\) of the \(i\)-th observation. By independence, the conditional law of the future observations does not depend on the value assigned to \(S_i\). Hence
\[
\mathcal L
\left(
S_{i+1},\ldots,S_M
\,\middle|\,
S_1=s_1,\ldots,S_i=s_i
\right)
=
\mathcal L
\left(
S_{i+1},\ldots,S_M
\,\middle|\,
S_1=s_1,\ldots,S_{i-1}=s_{i-1},S_i=s_i'
\right).
\]
The two conditional future laws can therefore be coupled identically. Under this coupling, we have \(S_j^{(s_1,\ldots,s_i,s_i')}=S_j'^{(s_1,\ldots,s_i,s_i')}\) almost surely for every \(j>i\). Hence
\(
\mathbb P
\left(
S_j^{(s_1,\ldots,s_i,s_i')}
\neq
S_j'^{(s_1,\ldots,s_i,s_i')}
\right)
=
0
\)
for every \(i<j\).

By Definition~\ref{def:Marton mixing matrix}, all off-diagonal entries of \(\Gamma\) are therefore zero, while \(\Gamma_{ii}=1\). Thus \(\Gamma=I_M\), and consequently \(\|\Gamma\|_{\mathrm{op}}=1\).
\end{proof}

\paragraph{Overlapping windows from a single trajectory.}

Heavily overlapping context windows from one trajectory should not be treated as independent observations merely because they appear as separate training examples. In a deterministic system, all such windows may be functions of the same random initial condition, and adjacent windows additionally share most of their coordinates. The number of stored windows may therefore substantially exceed the number of statistically distinct sampling units.
This issue is standard in dependent functional time-series analysis: temporal dependence must be accounted for explicitly rather than removed by representing the data as separate windows \cite{HormannKokoszka2010}.

\paragraph{Block-dependent and \(m\)-dependent observations.}

When dependence is local in time, observations may instead be grouped into blocks. For \(m\)-dependent or block-dependent sequences, sufficiently separated blocks are independent, and concentration can be formulated in terms of the resulting effective number of blocks. Marton-coupling inequalities naturally accommodate such dependence structures through the mixing matrix or an appropriate partition of the observations \cite{Paulin2015}.

\paragraph{Markov and total-variation mixing data.}

If the pretraining observations form a sufficiently mixing Markov chain, coupling probabilities can be controlled by its total-variation mixing properties. In this setting, the Marton framework remains applicable; see, for example, the Markov-chain bounded-differences inequalities of Paulin \cite{Paulin2015}.
For continuous functional dynamics, however, total variation may be too coarse. If \(T:H\to H\) is deterministic and \(Tf\neq Tg\), then
\(
d_{\mathrm{TV}}(\delta_{Tf},\delta_{Tg})=1
\)
even when \(\|Tf-Tg\|_H\) is arbitrarily small. The metric framework of Section~\ref{subsec:Wasserstein and metric dependence} is designed precisely for such situations.

\paragraph{Metric-dependent functional trajectories.}

The assumptions of the metric theorem arise naturally for dissipative systems. A contractive evolution can control the propagation of perturbations in the Hilbert norm, while a bounded or sufficiently concentrated initial law can provide the required transportation inequality. Transportation inequalities for dependent dynamical systems are developed in this form by Djellout, Guillin, and Wu \cite{DjelloutGuillinWu2004}.
The heat equation gives an explicit example: its semigroup contracts in \(L^2\), which induces a Wasserstein contraction for the corresponding context process, while a bounded initial-condition distribution provides a finite \(T_1\) constant.

\begin{example}[Metric dependence for heat-equation contexts]
\label{ex:Heat equation metric dependence}
Consider the heat equation \(\partial_tu=\nu\partial_{xx}u\) on \([0,1]\), with homogeneous Dirichlet boundary conditions and diffusivity \(\nu>0\). Let \(S_t:H\to H\) denote its semigroup on \(H=L^2([0,1])\), fix a sampling interval \(\Delta t>0\), and set \(\kappa_H=e^{-\nu\pi^2\Delta t}\in(0,1)\).

Suppose that the initial condition satisfies \(\|u_0\|_H\leqslant R_*\) almost surely. For a fixed context length \(n\), let \(F_k=(u_{k\Delta t},\ldots,u_{(k+n-1)\Delta t})\in E\). Define for \(F=(f_0,\ldots,f_{n-1})\) the deterministic context update by \(\Psi(F)=(f_1,\ldots,f_{n-1},S_{\Delta t}f_{n-1})\), and let \(K(F,\cdot)=\delta_{\Psi(F)}\).

Choose \(\beta>1\) such that \(\beta^{-1}+\kappa_H^2<1\), and define \(\kappa_E=(\beta^{-1}+\kappa_H^2)^{1/2}\). Then \(K\) satisfies the inequality
\(
W_1^{\|\cdot\|_\beta}
\bigl(
K(F,\cdot),
K(G,\cdot)
\bigr)
\leqslant
\kappa_E
\|F-G\|_\beta
\)
for every \(F,G\in E\). Moreover, the initial context law and every conditional law \(K(F,\cdot)\) satisfy \(T_1(C_E)\) for a common finite constant \(C_E\). Hence the heat-equation context process satisfies the assumptions of Proposition~\ref{prop:Path space transportation inequality}.
\end{example}

\begin{proof}
We first establish contraction of the heat semigroup. Let \(w(t)=S_tf-S_tg\). Since \(w\) satisfies the homogeneous Dirichlet heat equation,
\[
\begin{aligned}
\frac{\mathrm d}{\mathrm dt}
\|w(t)\|_H^2
&=
2\nu
\int_0^1
w(t,x)\,
\partial_{xx}w(t,x)\,
\mathrm dx
\\
&=
-2\nu
\int_0^1
|\partial_xw(t,x)|^2
\,\mathrm dx
\\
&\leqslant
-2\nu\pi^2
\|w(t)\|_H^2.
\end{aligned}
\]
The first equality follows by differentiating the squared \(L^2\)-norm and using the heat equation. The second follows from integration by parts and the homogeneous Dirichlet boundary conditions. The inequality is Poincaré's inequality. Grönwall's inequality therefore gives \(\|S_tf-S_tg\|_H\leqslant e^{-\nu\pi^2t}\|f-g\|_H\), and in particular \(\|S_{\Delta t}f-S_{\Delta t}g\|_H\leqslant\kappa_H\|f-g\|_H\).

Now let \(F=(f_0,\ldots,f_{n-1})\) and \(G=(g_0,\ldots,g_{n-1})\). Then
\[
\begin{aligned}
\|\Psi(F)-\Psi(G)\|_\beta^2
&=
\sum_{j=1}^{n-1}
\beta^{j-1}
\|f_j-g_j\|_H^2
+
\beta^{n-1}
\|S_{\Delta t}f_{n-1}-S_{\Delta t}g_{n-1}\|_H^2
\\
&\leqslant
\beta^{-1}
\sum_{j=1}^{n-1}
\beta^j
\|f_j-g_j\|_H^2
+
\beta^{n-1}
\kappa_H^2
\|f_{n-1}-g_{n-1}\|_H^2
\\
&\leqslant
\left(
\beta^{-1}+\kappa_H^2
\right)
\|F-G\|_\beta^2
\\
&=
\kappa_E^2
\|F-G\|_\beta^2.
\end{aligned}
\]
The first equality follows from the definitions of \(\Psi\) and \(\|\cdot\|_\beta\). The first inequality uses the heat-semigroup estimate. The second bounds each contribution by the complete weighted context norm. The final equality is the definition of \(\kappa_E\).

Since \(K(F,\cdot)=\delta_{\Psi(F)}\), Example \ref{ex:Dirac optimal transport} gives \(W_1^{\|\cdot\|_\beta}(K(F,\cdot),K(G,\cdot))=\|\Psi(F)-\Psi(G)\|_\beta\leqslant\kappa_E\|F-G\|_\beta\). By construction, \(\kappa_E<1\).

It remains to verify the \(T_1\) condition. Let \(\mu_F\) be the law of the initial \(F_0=(u_0,S_{\Delta t}u_0,\ldots,S_{(n-1)\Delta t}u_0)\) context. The semigroup estimate and the bound \(\|u_0\|_H\leqslant R_*\) give
\[
\|F_0\|_\beta^2
=
\sum_{j=0}^{n-1}
\beta^j
\|S_{j\Delta t}u_0\|_H^2
\leqslant
R_*^2
\sum_{j=0}^{n-1}
\beta^j
\kappa_H^{2j}
=:R_E^2.
\]
Thus \(\mu_F\) is supported on a ball of radius \(R_E\), whose diameter is at most \(2R_E\).
For any probability measure \(\eta\) with finite \(H(\eta\mid\mu_F)\),
\[
\begin{aligned}
W_1^{\|\cdot\|_\beta}
(\eta,\mu_F)
&\leqslant
2R_E
d_{\mathrm{TV}}(\eta,\mu_F)
\\
&\leqslant
2R_E
\sqrt{
\frac12
H(\eta\mid\mu_F)
}
\\
&=
\sqrt{
2R_E^2
H(\eta\mid\mu_F)
}.
\end{aligned}
\]
The first inequality uses the diameter bound on the common support and the second is Pinsker's inequality. Hence \(\mu_F\) satisfies \(T_1(R_E^2)\). Finally, every conditional law \(K(F,\cdot)\) is a Dirac measure. A Dirac measure satisfies \(T_1(C)\) for every \(C>0\): finite relative entropy with respect to \(\delta_{\Psi(F)}\) forces the comparison measure to equal \(\delta_{\Psi(F)}\), in which case the Wasserstein distance is zero. Thus all conditional laws satisfy \(T_1(R_E^2)\) as well. Therefore the assumptions of Proposition~\ref{prop:Path space transportation inequality} hold with \(C_E=R_E^2\) and dependence coefficient \(\kappa_E<1\).
\end{proof}

The two heat-equation examples illustrate that the same physical system can satisfy the two statistical frameworks for different reasons. With independently sampled trajectories, the Marton result uses independence across initial conditions. Along a single dependent trajectory, the metric result instead uses dissipativity of the dynamics.

\paragraph{Nonmixing dynamics.}

Not every functional system satisfies the metric condition with \(\kappa<1\). Periodic advection may preserve \(L^2\) distances, while undamped wave dynamics can preserve an energy norm. Such systems are therefore not natural candidates for the Wasserstein-contraction data theorem.
They can nevertheless be handled by the Marton result when training consists of independently sampled trajectories. In that case, independence comes from the sampling design rather than from temporal mixing of the dynamics.

The Marton and metric results should therefore be viewed as complementary. The Marton framework is particularly natural when the sampling design produces independent or weakly dependent blocks, whereas the metric framework is better adapted to temporally dependent functional processes whose dynamics are stable in the underlying function-space metric. The appropriate generalisation theorem is therefore determined jointly by the sampling scheme and the dependence structure of the data-generating dynamics.


\section{Conclusion and future work}\label{sec:Conclusion and future work}

\paragraph{Conclusion.}
This thesis has developed a rigorous probabilistic framework for finite-context autoregressive transformers acting on sequences of functions. The proposed architecture combines a bounded function-space transformer with coefficient-valued output heads in a fixed approximation space, producing a context-dependent Gaussian law with non-diagonal covariance and a positive covariance floor. Taking the complete functional context as the state turns recursive shift-and-append generation into a time-homogeneous general-state-space Markov chain. The conditional Gaussian laws are shown to be well defined, to have common support, and to depend measurably on the context. Their common support and covariance bounds yield a global \(n\)-step Doeblin minorisation, which gives a unique invariant distribution and uniform geometric convergence in total variation from every initial distribution; the Foster--Lyapunov drift estimate separately provides finite weighted second-moment control. For a fixed parameter value, we also establish two dependent-data pretraining guarantees: a Marton-coupling bound for bounded Wasserstein transition losses and a transportation-based bound for Lipschitz losses under metric dependence, both with explicit sample-complexity statements under zero empirical risk. The analysis therefore includes probability-kernel well-posedness, approximation restrictions, long-time inference behaviour, and fixed-parameter pretraining generalisation.

The in-context analysis extends these concentration tools to fixed pretrained models evaluated on functional task trajectories. A model-comparison lemma for squared Wasserstein risk gives simultaneous comparator bounds, explicit mixing-time and metric dependence factors, and sample complexities above a residual task error. Heat-equation trajectories and a recursion with infinite-rank functional noise verify the assumptions in concrete settings. Projection error and the positive covariance floor quantify restrictions on attainable accuracy, while the distinction between evaluated windows and visible context length separates risk concentration from improvement with longer prompts.

\paragraph{Future work.}
Several directions follow from the results of this thesis. First, the approximation--stability trade-offs associated with the approximation dimension, covariance floor, context length, and conditional output law should be studied theoretically and empirically. Adaptive or increasing approximation spaces could reduce projection error, while Gaussian mixtures or conditional function-space diffusion mechanisms may provide more expressive transition laws. Second, the global boundedness, common-support, and covariance-floor conditions used to obtain the Doeblin minorisation could be relaxed. Models admitting only local small sets may instead require Harris-type drift conditions, while sharper arguments would be needed to obtain quantitatively useful convergence rates. Third, empirical experiments should assess recursive generation, sensitivity to the architectural parameters appearing in the bounds, and the extent to which the predicted long-run behaviour is observed. On the statistical side, an important extension is to move beyond fixed-parameter concentration and control parameters selected from the same dependent data, for example through uniform or algorithm-dependent arguments adapted to functional observations. Finally, combining temporal memory with functional demonstrations would connect the present autoregressive setting with in-context operator learning, allowing an inferred operator to be used for uncertainty-aware recursive generation.


\clearpage

\printbibliography

\clearpage

\appendix

\section{Lipschitz regularity of the functional autoregressive transformer}
\label{app:Lipschitz regularity}

This appendix derives finite global Lipschitz constants for the function-space transformer, its coefficient heads, and the resulting Gaussian output kernel. These estimates formalise the regularity supplied by the bounded projections, finite operator-norm bounds, and Lipschitz nonlinearities in the architecture. In particular, they provide the model constant used in the metric-dependence analysis of the pretraining loss.
Throughout the appendix, the context space is equipped with the weighted norm from Definition~\ref{def:Weighted context norm}.

\paragraph{Wasserstein stability of the Gaussian output.}

We first control the Wasserstein distance between the Gaussian output laws associated with two contexts. The covariance floor is independent of the context, so it cancels when the corresponding covariance factors are compared.

\begin{lemma}[]\label{lem:Hilbert--Schmidt distance between covariance factors}
Let \(c_1,\ldots,c_r,\widetilde c_1,\ldots,\widetilde c_r\in H\), and define \(B_1,B_2:\mathbb R^r\to H\) by \(B_1a=\sum_{k=1}^r a_kc_k\) and \(B_2a=\sum_{k=1}^r a_k\widetilde c_k\). Then \(\|B_1-B_2\|_{\mathrm{HS}}^2=\sum_{k=1}^r\|c_k-\widetilde c_k\|_H^2\).
\end{lemma}

\begin{proof}
Let \(e_1,\ldots,e_r\) denote the standard orthonormal basis of \(\mathbb R^r\). Then, we compute
\[
\|B_1-B_2\|_{\mathrm{HS}}^2 = \sum_{k=1}^r\|(B_1-B_2)e_k\|_H^2 =\sum_{k=1}^r\|c_k-\widetilde c_k\|_H^2,
\]
where the first equality uses Definition \ref{def:Hilbert--Schmidt norm} and the second equality follows because \(B_1e_k=c_k\) and \(B_2e_k=\widetilde c_k\).
\end{proof}

\begin{proposition}[Wasserstein stability of the Gaussian output]\label{prop:Wasserstein stability of the Gaussian output}
For every \(F,G\in E\),
\[
W_2^2\bigl(Q_\Theta(F,\cdot),Q_\Theta(G,\cdot)\bigr)
\leqslant
\|m_\Theta(F)-m_\Theta(G)\|_H^2
+
\sum_{k=1}^r\|c_{k,\Theta}(F)-c_{k,\Theta}(G)\|_H^2.
\]
\end{proposition}

\begin{proof}
We start by defining an augmented covariance factor \(\overline B_\Theta(F):\mathbb R^{d+r}\to H\) by
\(
\overline B_\Theta(F)(z,a)
=\sigma_0\mathsf E_dz+\sum_{k=1}^ra_kc_{k,\Theta}(F).
\)
Then
\(
C_\Theta(F)=\overline B_\Theta(F)\overline B_\Theta(F)^*.
\)
The first \(d\) columns of \(\overline B_\Theta(F)\) are the context-independent functions \(\sigma_0e_1,\ldots,\sigma_0e_d\). Therefore they cancel when two factors are subtracted, and
\(
\|\overline B_\Theta(F)-\overline B_\Theta(G)\|_{\mathrm{HS}}^2
=\sum_{k=1}^r\|c_{k,\Theta}(F)-c_{k,\Theta}(G)\|_H^2.
\)
The Gaussian Wasserstein formula and Proposition~\ref{prop:Coupling bound for covariance factors} now give
\[
\begin{aligned}
W_2^2\bigl(Q_\Theta(F,\cdot),Q_\Theta(G,\cdot)\bigr)
&=\|m_\Theta(F)-m_\Theta(G)\|_H^2
+d_{\mathrm{BW}}^2\bigl(C_\Theta(F),C_\Theta(G)\bigr)\\
&\leqslant\|m_\Theta(F)-m_\Theta(G)\|_H^2
+\|\overline B_\Theta(F)-\overline B_\Theta(G)\|_{\mathrm{HS}}^2,
\end{aligned}
\]
which is the required estimate.
\end{proof}



This estimate reduces the Wasserstein stability of the Gaussian output distribution to the stability of the mean and covariance-factor maps produced by the transformer.

\paragraph{Architectural Lipschitz bound.}

We now derive a global Lipschitz bound for the mean and covariance-factor maps from the constraints built into the functional transformer architecture; see Section~\ref{subsec:Functional autoregressive transformer}. The argument proceeds in stages. We first show that the bounded projection used throughout the architecture is non-expansive. We then control the nonlinear attention mechanism by combining the boundedness of the latent tokens with the Lipschitz continuity of softmax, first for a single attention head and then for multi-head attention. Next, we establish a Lipschitz bound for the feed-forward block. These estimates are then propagated through the residual connections, projections, and successive transformer layers to obtain a global Lipschitz constant for the final latent representation. Finally, we pass this bound through the mean and covariance-factor heads to obtain explicit Lipschitz constants for the Gaussian parameters.

Recall that the main role of the bounded projection is to ensure that every token entering an attention layer belongs to a fixed ball in the latent Hilbert space \(\mathcal H=L^2(U;\mathbb R^{d_{\mathrm{model}}})\). This uniform boundedness is particularly important for the attention mechanism because the query--key scores are bilinear in the latent tokens.

Let \(\mathcal H^n\) be the product space equipped with the product norm \(\|V\|_{\mathcal H^n}^2=\sum_{j=0}^{n-1}\|v_j\|_{\mathcal H}^2\), where \(V=(v_0,\ldots,v_{n-1})\). Define \(\mathcal B_R=\{V=(v_0,\ldots,v_{n-1})\in\mathcal H^n:\|v_j\|_{\mathcal H}\leqslant R \text{ for all }j\}\). By construction of the functional transformer, \(V^{(\ell)}\in\mathcal B_R\) for every transformer layer \(\ell\).

\begin{lemma}[Non-expansiveness of the bounded projection]
\label{lem:Non-expansiveness of the bounded projection}
For every \(u,v\in\mathcal H\), the projection \(\Pi_R:\mathcal H\to\mathcal H\) satisfies \(\|\Pi_R(u)-\Pi_R(v)\|_{\mathcal H}\leqslant\|u-v\|_{\mathcal H}\). Consequently, the componentwise projection \(\Pi_R^{(n)}:\mathcal H^n\to\mathcal H^n\) is also non-expansive.
\end{lemma}

\begin{proof}
Clearly, the closed ball \(\{w\in\mathcal H:\|w\|_{\mathcal H}\leqslant R\}\) is nonempty, closed and convex. By the characterization of metric projections onto closed convex subsets of Hilbert spaces in Brezis \cite[Theorem~5.2]{Brezis2011}, if \(x=\Pi_R(u)\) and \(y=\Pi_R(v)\), then \(\langle u-x,y-x\rangle_{\mathcal H}\leqslant0\) and \(\langle v-y,x-y\rangle_{\mathcal H}\leqslant0\).

Rearranging these inequalities gives \(\langle u-x,x-y\rangle_{\mathcal H}\geq0\) and \(\langle v-y,x-y\rangle_{\mathcal H}\leqslant0\). Subtracting the second inequality from the first yields \(\|x-y\|_{\mathcal H}^2\leqslant\langle u-v,x-y\rangle_{\mathcal H}\).

By the Cauchy--Schwarz inequality, \(\|x-y\|_{\mathcal H}^2\leqslant\|u-v\|_{\mathcal H}\|x-y\|_{\mathcal H}\). If \(x=y\), the result is immediate. Otherwise, dividing by \(\|x-y\|_{\mathcal H}\) gives \(\|\Pi_R(u)-\Pi_R(v)\|_{\mathcal H}\leqslant\|u-v\|_{\mathcal H}\).

For \(U=(u_0,\ldots,u_{n-1})\) and \(V=(v_0,\ldots,v_{n-1})\), applying this inequality componentwise gives \(\|\Pi_R^{(n)}(U)-\Pi_R^{(n)}(V)\|_{\mathcal H^n}^2\leqslant\sum_{j=0}^{n-1}\|u_j-v_j\|_{\mathcal H}^2=\|U-V\|_{\mathcal H^n}^2\). Hence \(\Pi_R^{(n)}\) is also non-expansive.
\end{proof}

The preceding lemma shows that the projections used throughout the architecture do not amplify perturbations between latent representations. Their role is therefore particularly convenient: they keep all tokens inside the bounded set \(\mathcal B_R\), while contributing no additional factor to the Lipschitz constants derived below.

The next source of nonlinearity is the attention mechanism. There are two components to control. First, the attention scores depend bilinearly on the query and key functions. Their inputs are uniformly bounded because they belong to \(\mathcal B_R\). Second, the scores are transformed into attention weights through the softmax map. We therefore begin by recalling the Lipschitz continuity of softmax.

\begin{lemma}[Lipschitz continuity of softmax]
\label{lem:Lipschitz continuity of softmax}
Let \(\operatorname{softmax}:\mathbb R^n\to\mathbb R^n\) be defined by \(\operatorname{softmax}(a)_j=\exp(a_j)/\sum_{k=0}^{n-1}\exp(a_k)\), for \(j\in\{0,\ldots,n-1\}\). We equip \(\mathbb R^n\) with the Euclidean norm \(\|a\|_2^2=\sum_{j=1}^n|a_j|^2\). Then, for every \(a,b\in\mathbb R^n\), \(\|\operatorname{softmax}(a)-\operatorname{softmax}(b)\|_2\leqslant\frac12\|a-b\|_2\).
\end{lemma}

\begin{proof}
This is the \(p=2\) case of the global Lipschitz bound for the softmax operator established in \cite[Theorem~1]{Nair2026}. In particular, the Jacobian of the softmax map is \(J_{\operatorname{softmax}}(a)=\operatorname{Diag}(p)-pp^\top\), where \(p=\operatorname{softmax}(a)\), and its induced Euclidean operator norm is bounded above by \(1/2\). Hence the softmax map is globally \(1/2\)-Lipschitz with respect to \(\|\cdot\|_2\).
\end{proof}

We can now combine the boundedness of the functional tokens, the operator-norm constraints on the query, key, and value operators, and the Lipschitz continuity of softmax to control a complete attention head. The main point is that a perturbation of the input affects the attention output in two ways: it changes the attention weights through the query--key scores, and it changes the value functions being averaged. The following lemma controls these two effects separately.

\begin{lemma}[Lipschitz continuity of a functional attention head]
\label{lem:functional-attention-Lipschitz}
Consider one attention head in layer \(\ell\), with operators satisfying \(\|\mathsf Q_h^{(\ell)}\|_{\mathrm{op}}\leqslant B_Q^{(\ell,h)}\), \(\|\mathsf K_h^{(\ell)}\|_{\mathrm{op}}\leqslant B_K^{(\ell,h)}\), and \(\|\mathsf V_h^{(\ell)}\|_{\mathrm{op}}\leqslant B_V^{(\ell,h)}\). Then, the corresponding attention operator is Lipschitz: for any \(U,V \in \mathcal B_R\), \(\|\mathcal A_h^{(\ell)}(U)-\mathcal A_h^{(\ell)}(V)\|_{\mathcal H_v^n}\leqslant L_{\mathrm{att}}^{(\ell,h)}\|U-V\|_{\mathcal H^n}\), where \(L_{\mathrm{att}}^{(\ell,h)}=\sqrt n\,B_V^{(\ell,h)}\left[1+{B_Q^{(\ell,h)}B_K^{(\ell,h)}R^2\sqrt n(\sqrt n+1)}/{(2\tau)}\right]\).
\end{lemma}

\begin{proof}
Let \(U=(u_0,\ldots,u_{n-1})\in\mathcal B_R\) and \(V=(v_0,\ldots,v_{n-1})\in\mathcal B_R\). We first bound the change in the attention weights and then we compute the bound of the resulting attention output.

\textit{Step 1: Bound for the attention scores.}
Fix a query position \(i\). Then, for the \(j\)-th attention score,
\[
\begin{aligned}
|a_{ij}(U)-a_{ij}(V)|
&=
\frac{1}{\tau}
\left|
\left\langle \mathsf Q_h^{(\ell)}u_i,\mathsf K_h^{(\ell)}u_j\right\rangle_{\mathcal H_q}
-
\left\langle \mathsf Q_h^{(\ell)}v_i,\mathsf K_h^{(\ell)}v_j\right\rangle_{\mathcal H_q}
\right|\\
&\leqslant
\frac{1}{\tau}
\left|
\left\langle
\mathsf Q_h^{(\ell)}(u_i-v_i),
\mathsf K_h^{(\ell)}u_j
\right\rangle_{\mathcal H_q}
\right|
+
\frac{1}{\tau}
\left|
\left\langle
\mathsf Q_h^{(\ell)}v_i,
\mathsf K_h^{(\ell)}(u_j-v_j)
\right\rangle_{\mathcal H_q}
\right|\\
&\leqslant
\frac{1}{\tau}
\|\mathsf Q_h^{(\ell)}\|_{\mathrm{op}}
\|u_i-v_i\|_{\mathcal H}
\|\mathsf K_h^{(\ell)}\|_{\mathrm{op}}
\|u_j\|_{\mathcal H} +
\frac{1}{\tau}
\|\mathsf Q_h^{(\ell)}\|_{\mathrm{op}}
\|v_i\|_{\mathcal H}
\|\mathsf K_h^{(\ell)}\|_{\mathrm{op}}
\|u_j-v_j\|_{\mathcal H}\\
&\leqslant
\frac{B_Q^{(\ell,h)}B_K^{(\ell,h)}R}{\tau}
\left(
\|u_i-v_i\|_{\mathcal H}
+
\|u_j-v_j\|_{\mathcal H}
\right).
\end{aligned}
\]
The first equality follows from the definition of the attention scores. The first inequality is obtained by adding and subtracting \(\langle \mathsf Q_h^{(\ell)}v_i,\mathsf K_h^{(\ell)}u_j\rangle_{\mathcal H_q}\) and applying the triangle inequality. The second inequality follows from the Cauchy--Schwarz inequality and definition of operator norm. The final inequality uses the architectural bounds \(\|\mathsf Q_h^{(\ell)}\|_{\mathrm{op}}\leqslant B_Q^{(\ell,h)}\), \(\|\mathsf K_h^{(\ell)}\|_{\mathrm{op}}\leqslant B_K^{(\ell,h)}\), and the fact that \(U,V\in\mathcal B_R\), so that \(\|u_j\|_{\mathcal H}\leqslant R\) and \(\|v_i\|_{\mathcal H}\leqslant R\).

For fixed \(i\), define the score vector \(a_i(U)=(a_{i0}(U),\ldots,a_{i,n-1}(U))\in\mathbb R^n\). Let \(c={B_Q^{(\ell,h)}B_K^{(\ell,h)}R}/{\tau}\) and using the previous estimate:
\[
\begin{aligned}
\|a_i(U)-a_i(V)\|_2
&\leqslant c\left[\sum_{j=0}^{n-1}\left(\|u_i-v_i\|_{\mathcal H}+\|u_j-v_j\|_{\mathcal H}\right)^2\right]^{1/2}\\
&\leqslant c\left[\left(\sum_{j=0}^{n-1}\|u_i-v_i\|_{\mathcal H}^2\right)^{1/2}+\left(\sum_{j=0}^{n-1}\|u_j-v_j\|_{\mathcal H}^2\right)^{1/2}\right]\\
&=c\left[\sqrt n\,\|u_i-v_i\|_{\mathcal H}+\|U-V\|_{\mathcal H^n}\right]\\
&\leqslant c(\sqrt n+1)\|U-V\|_{\mathcal H^n}.
\end{aligned}
\]
The first inequality applies the componentwise score estimate. The second is the triangle inequality for the Euclidean norm, applied to the vectors whose \(j\)-th components are respectively \(\|u_i-v_i\|_{\mathcal H}\) and \(\|u_j-v_j\|_{\mathcal H}\). The equality uses \(\sum_{j=0}^{n-1}\|u_i-v_i\|_{\mathcal H}^2=n\|u_i-v_i\|_{\mathcal H}^2\) and the definition of the product norm on \(\mathcal H^n\). The final inequality follows because \(\|u_i-v_i\|_{\mathcal H}\leqslant\|U-V\|_{\mathcal H^n}\).

By Lemma \ref{lem:Lipschitz continuity of softmax}, the corresponding attention-weight vectors satisfy \(\|\alpha_i(U)-\alpha_i(V)\|_2\leqslant\frac{B_Q^{(\ell,h)}B_K^{(\ell,h)}R}{2\tau}(\sqrt n+1)\|U-V\|_{\mathcal H^n}\).

\textit{Step 2: Bound for the attention output.} For position \(i\), we have \(\mathcal A_h^{(\ell)}(U)_i=\sum_{j=0}^{n-1}\alpha_{ij}(U)\mathsf V_h^{(\ell)}u_j\). For each \(j\), we add and subtract the intermediate term \(\alpha_{ij}(V)\mathsf V_h^{(\ell)}u_j\). This gives
\[
\begin{aligned}
\mathcal A_h^{(\ell)}(U)_i-\mathcal A_h^{(\ell)}(V)_i
&=
\sum_{j=0}^{n-1}
\left[
\alpha_{ij}(U)\mathsf V_h^{(\ell)}u_j
-
\alpha_{ij}(V)\mathsf V_h^{(\ell)}v_j
\right]\\
&=
\sum_{j=0}^{n-1}
\left[
\alpha_{ij}(U)\mathsf V_h^{(\ell)}u_j
-
\alpha_{ij}(V)\mathsf V_h^{(\ell)}u_j
+
\alpha_{ij}(V)\mathsf V_h^{(\ell)}u_j
-
\alpha_{ij}(V)\mathsf V_h^{(\ell)}v_j
\right]\\
&=
\sum_{j=0}^{n-1}
\bigl(
\alpha_{ij}(U)-\alpha_{ij}(V)
\bigr)
\mathsf V_h^{(\ell)}u_j
+
\sum_{j=0}^{n-1}
\alpha_{ij}(V)
\mathsf V_h^{(\ell)}(u_j-v_j).
\end{aligned}
\]
The first term isolates the change in the attention weights, while the second isolates the change in the value functions.

Let us estimate the first term:
\[
\begin{aligned}
\left\|
\sum_{j=0}^{n-1}
\bigl(\alpha_{ij}(U)-\alpha_{ij}(V)\bigr)
\mathsf V_h^{(\ell)}u_j
\right\|_{\mathcal H_v}
&\leqslant
\sum_{j=0}^{n-1}
\left|
\alpha_{ij}(U)-\alpha_{ij}(V)
\right|
\left\|
\mathsf V_h^{(\ell)}u_j
\right\|_{\mathcal H_v}\\
&\leqslant
B_V^{(\ell,h)}
\sum_{j=0}^{n-1}
\left|
\alpha_{ij}(U)-\alpha_{ij}(V)
\right|
\|u_j\|_{\mathcal H}\\
&\leqslant
B_V^{(\ell,h)}R
\sum_{j=0}^{n-1}
\left|
\alpha_{ij}(U)-\alpha_{ij}(V)
\right|\\
&\leqslant
B_V^{(\ell,h)}R\sqrt n\,
\|\alpha_i(U)-\alpha_i(V)\|_2.
\end{aligned}
\]

The first inequality follows from the triangle inequality in \(\mathcal H_v\). The second uses the architecture constraint \(\|\mathsf V_h^{(\ell)}\|_{\mathrm{op}}\leqslant B_V^{(\ell,h)}\). The third uses the bounded-token property \(\|u_j\|_{\mathcal H}\leqslant R\). The final inequality is the Cauchy--Schwarz inequality in \(\mathbb R^n\), \(\sum_{j=0}^{n-1}|x_j|\leqslant\sqrt n\,(\sum_{j=0}^{n-1}|x_j|^2)^{1/2}\).

Using the previously established estimate for the attention weights, this gives
\[
\left\|
\sum_{j=0}^{n-1}
\bigl(\alpha_{ij}(U)-\alpha_{ij}(V)\bigr)
\mathsf V_h^{(\ell)}u_j
\right\|_{\mathcal H_v}
\leqslant
\frac{
B_V^{(\ell,h)}
B_Q^{(\ell,h)}
B_K^{(\ell,h)}
R^2
\sqrt n(\sqrt n+1)
}{
2\tau
}
\|U-V\|_{\mathcal H^n}.
\]
For the second term
\[
\begin{aligned}
\left\|
\sum_{j=0}^{n-1}
\alpha_{ij}(V)
\mathsf V_h^{(\ell)}(u_j-v_j)
\right\|_{\mathcal H_v}
&\leqslant
\sum_{j=0}^{n-1}
\alpha_{ij}(V)
\left\|
\mathsf V_h^{(\ell)}(u_j-v_j)
\right\|_{\mathcal H_v}\\
&\leqslant
B_V^{(\ell,h)}
\sum_{j=0}^{n-1}
\alpha_{ij}(V)
\|u_j-v_j\|_{\mathcal H}\\
&\leqslant
B_V^{(\ell,h)}
\left(
\sum_{j=0}^{n-1}
\alpha_{ij}(V)^2
\right)^{1/2}
\left(
\sum_{j=0}^{n-1}
\|u_j-v_j\|_{\mathcal H}^2
\right)^{1/2}\\
&\leqslant
B_V^{(\ell,h)}
\|U-V\|_{\mathcal H^n}.
\end{aligned}
\]
The first inequality is again the triangle inequality in \(\mathcal H_v\), while the second follows from the operator-norm bound on \(\mathsf V_h^{(\ell)}\). The third inequality is Cauchy--Schwarz in \(\mathbb R^n\). For the final inequality, observe that the attention coefficients form a probability vector, so \(\|\alpha_i(V)\|_2\leqslant\|\alpha_i(V)\|_1=1\), while the second factor is exactly \(\|U-V\|_{\mathcal H^n}\).

Combining the two terms and applying the triangle inequality on the decomposition, we obtain
\[
\|\mathcal A_h^{(\ell)}(U)_i-\mathcal A_h^{(\ell)}(V)_i\|_{\mathcal H_v}
\leqslant
B_V^{(\ell,h)}
\left(
1+
\frac{
B_Q^{(\ell,h)}
B_K^{(\ell,h)}
R^2
\sqrt n(\sqrt n+1)
}{
2\tau
}
\right)
\|U-V\|_{\mathcal H^n}.
\]
This bound holds for every output position \(i\). Therefore, using the product norm on \(\mathcal H_v^n\),
\[
\begin{aligned}
\|\mathcal A_h^{(\ell)}(U)-\mathcal A_h^{(\ell)}(V)\|_{\mathcal H_v^n}^2
&=
\sum_{i=0}^{n-1}
\|\mathcal A_h^{(\ell)}(U)_i-\mathcal A_h^{(\ell)}(V)_i\|_{\mathcal H_v}^2\\
&\leqslant
n
\left[
B_V^{(\ell,h)}
\left(
1+
\frac{
B_Q^{(\ell,h)}
B_K^{(\ell,h)}
R^2
\sqrt n(\sqrt n+1)
}{
2\tau
}
\right)
\right]^2
\|U-V\|_{\mathcal H^n}^2.
\end{aligned}
\]

Taking square roots concludes the proof.
\end{proof}

The constant \(L_{\mathrm{att}}^{(\ell,h)}\) therefore accounts for both sources of sensitivity in a single attention head. The term involving \(B_Q^{(\ell,h)}B_K^{(\ell,h)}R^2/\tau\) arises from the dependence of the attention weights on the query and key functions, while the remaining contribution arises from the change in the value functions themselves.

A multi-head attention layer applies several such heads in parallel and subsequently combines their outputs through the linear operator \(\mathsf W_O^{(\ell)}\). The single-head bounds can therefore be combined to obtain a Lipschitz constant for the complete multi-head attention operator.

\begin{corollary}[Lipschitz continuity of multi-head attention]
\label{cor:Lipschitz continuity of multi-head attention}
Consider the multi-head attention in layer \(\ell\). Then, the corresponding multi-head attention operator is Lipschitz: for any \(U,V \in \mathcal B_R\), \(\|\mathcal A_{\mathrm{MHA}}^{(\ell)}(U)-\mathcal A_{\mathrm{MHA}}^{(\ell)}(V)\|_{\mathcal H^n}\leqslant L_{\mathrm{MHA}}^{(\ell)}\|U-V\|_{\mathcal H^n}\), where \(L_{\mathrm{MHA}}^{(\ell)}=B_O^{(\ell)}\left(\sum_{h=1}^{N_h}(L_{\mathrm{att}}^{(\ell,h)})^2\right)^{1/2}\).
\end{corollary}

\begin{proof}
For each context position \(i\), the multi-head attention output is obtained by concatenating the outputs of the \(N_h\) attention heads and then applying the output operator \(\mathsf W_O^{(\ell)}\). Hence
\[
\begin{aligned}
&\|\mathcal A_{\mathrm{MHA}}^{(\ell)}(U)_i-\mathcal A_{\mathrm{MHA}}^{(\ell)}(V)_i\|_{\mathcal H}\\
&\qquad=
\left\|
\mathsf W_O^{(\ell)}
\left(
\operatorname{Concat}
\bigl(
\mathcal A_1^{(\ell)}(U)_i,\ldots,\mathcal A_{N_h}^{(\ell)}(U)_i
\bigr)
-
\operatorname{Concat}
\bigl(
\mathcal A_1^{(\ell)}(V)_i,\ldots,\mathcal A_{N_h}^{(\ell)}(V)_i
\bigr)
\right)
\right\|_{\mathcal H}\\
&\qquad\leqslant
\|\mathsf W_O^{(\ell)}\|_{\mathrm{op}}
\left[
\sum_{h=1}^{N_h}
\|\mathcal A_h^{(\ell)}(U)_i-\mathcal A_h^{(\ell)}(V)_i\|_{\mathcal H_v}^2
\right]^{1/2}\\
&\qquad\leqslant
B_O^{(\ell)}
\left[
\sum_{h=1}^{N_h}
\|\mathcal A_h^{(\ell)}(U)_i-\mathcal A_h^{(\ell)}(V)_i\|_{\mathcal H_v}^2
\right]^{1/2}.
\end{aligned}
\]

The first inequality follows from the definition of the operator norm. The norm of the concatenated vector is the square root of the sum of the squared norms of its head components, since concatenation places the head outputs in separate channel coordinates. The second inequality uses the architectural bound \(\|\mathsf W_O^{(\ell)}\|_{\mathrm{op}}\leqslant B_O^{(\ell)}\).

We now sum over the \(n\) output positions:
\[
\begin{aligned}
\|\mathcal A_{\mathrm{MHA}}^{(\ell)}(U)-\mathcal A_{\mathrm{MHA}}^{(\ell)}(V)\|_{\mathcal H^n}^2&=
\sum_{i=0}^{n-1}
\|\mathcal A_{\mathrm{MHA}}^{(\ell)}(U)_i-\mathcal A_{\mathrm{MHA}}^{(\ell)}(V)_i\|_{\mathcal H}^2\\
&\leqslant
\left(B_O^{(\ell)}\right)^2
\sum_{i=0}^{n-1}
\sum_{h=1}^{N_h}
\|\mathcal A_h^{(\ell)}(U)_i-\mathcal A_h^{(\ell)}(V)_i\|_{\mathcal H_v}^2\\
&=
\left(B_O^{(\ell)}\right)^2
\sum_{h=1}^{N_h}
\|\mathcal A_h^{(\ell)}(U)-\mathcal A_h^{(\ell)}(V)\|_{\mathcal H_v^n}^2\\
&\leqslant
\left(B_O^{(\ell)}\right)^2
\sum_{h=1}^{N_h}
\left(L_{\mathrm{att}}^{(\ell,h)}\right)^2
\|U-V\|_{\mathcal H^n}^2.
\end{aligned}
\]
The final inequality follows by applying Lemma \ref{lem:functional-attention-Lipschitz} to each attention head and taking square roots concludes the proof.
\end{proof}

We have now controlled the attention component of a transformer layer. The remaining nonlinear component is the feed-forward network. Since it acts independently at each context position, its Lipschitz constant can be obtained first at the level of an individual functional token and then lifted directly to the product space \(\mathcal H^n\).

\begin{lemma}[Lipschitz continuity of the feed-forward block]
\label{lem:Lipschitz continuity of the feed-forward block}
Suppose that the activation function \(\sigma\) is \(L_\sigma\)-Lipschitz. Then, the corresponding feed-forward operator is Lipschitz: for any \(U,V\in\mathcal H^n\), \(\|\mathcal F_\ell(U)-\mathcal F_\ell(V)\|_{\mathcal H^n} \leqslant L_{\mathrm{FF}}^{(\ell)} \|U-V\|_{\mathcal H^n}\), where \(L_{\mathrm{FF}}^{(\ell)}=B_2^{(\ell)}L_\sigma B_1^{(\ell)}\).
\end{lemma}

\begin{proof}
Let \(U=(u_0,\ldots,u_{n-1})\in \mathcal H^n\) and \(V=(v_0,\ldots,v_{n-1})\in\mathcal H^n\). Since the feed-forward network acts independently on each context position, its \(j\)-th component is
\(
\mathcal F_\ell(u_j)(x)
=
\mathsf W_2^{(\ell)}
\sigma\bigl(
\mathsf W_1^{(\ell)}u_j(x)+b_1^{(\ell)}
\bigr)
+
b_2^{(\ell)}
\).

We first estimate the difference at a single context position:
\[
\begin{aligned}
\|\mathcal F_\ell(u_j)-\mathcal F_\ell(v_j)\|_{\mathcal H}
&=
\left\|
\mathsf W_2^{(\ell)}
\left[
\sigma\bigl(\mathsf W_1^{(\ell)}u_j+b_1^{(\ell)}\bigr)
-
\sigma\bigl(\mathsf W_1^{(\ell)}v_j+b_1^{(\ell)}\bigr)
\right]
\right\|_{\mathcal H}\\
&\leqslant
\|\mathsf W_2^{(\ell)}\|_{\mathrm{op}}
\left\|
\sigma\bigl(\mathsf W_1^{(\ell)}u_j+b_1^{(\ell)}\bigr)
-
\sigma\bigl(\mathsf W_1^{(\ell)}v_j+b_1^{(\ell)}\bigr)
\right\|_{\mathcal H}\\
&\leqslant
\|\mathsf W_2^{(\ell)}\|_{\mathrm{op}}
L_\sigma
\|\mathsf W_1^{(\ell)}(u_j-v_j)\|_{\mathcal H}\\
&\leqslant
B_2^{(\ell)}
L_\sigma
B_1^{(\ell)}
\|u_j-v_j\|_{\mathcal H}\\
&=
L_{\mathrm{FF}}^{(\ell)}
\|u_j-v_j\|_{\mathcal H}.
\end{aligned}
\]
In the first equality \(b_2^{(\ell)}\) cancels when the two feed-forward outputs are subtracted. The first inequality follows from the definition of the operator norm of \(\mathsf W_2^{(\ell)}\). The second inequality uses the \(L_\sigma\)-Lipschitz continuity of \(\sigma\). Since the same input bias \(b_1^{(\ell)}\) appears in both arguments, their difference is \(\mathsf W_1^{(\ell)}(u_j-v_j)\). The third uses the architectural bounds \(\|\mathsf W_1^{(\ell)}\|_{\mathrm{op}}\leqslant B_1^{(\ell)}\) and \(\|\mathsf W_2^{(\ell)}\|_{\mathrm{op}}\leqslant B_2^{(\ell)}\).

We now combine these componentwise estimates using the product norm on \(\mathcal H^n\):
\[
\begin{aligned}
\|\mathcal F_\ell(U)-\mathcal F_\ell(V)\|_{\mathcal H^n}^2
&=
\sum_{j=0}^{n-1}
\|\mathcal F_\ell(u_j)-\mathcal F_\ell(v_j)\|_{\mathcal H}^2\\
&\leqslant
\sum_{j=0}^{n-1}
\left(
L_{\mathrm{FF}}^{(\ell)}
\|u_j-v_j\|_{\mathcal H}
\right)^2\\
&=
\left(
L_{\mathrm{FF}}^{(\ell)}
\right)^2
\sum_{j=0}^{n-1}
\|u_j-v_j\|_{\mathcal H}^2\\
&=
\left(
L_{\mathrm{FF}}^{(\ell)}
\right)^2
\|U-V\|_{\mathcal H^n}^2.
\end{aligned}
\]
Taking square roots concludes the proof.
\end{proof}

The attention and feed-forward components of each transformer layer therefore have explicit Lipschitz bounds. It remains to account for the residual connections and projections that combine these components. Residual connections add the identity map to the corresponding attention or feed-forward transformation, while the projections do not increase distances by Lemma~\ref{lem:Non-expansiveness of the bounded projection}. This allows the Lipschitz constants to be propagated through each complete transformer layer and subsequently through the full depth of the network.

\begin{proposition}[Lipschitz continuity of the functional transformer]
\label{prop:Lipschitz continuity of the functional transformer}
The latent next-function representation operator \(h_\Theta:E\to\mathcal H\) is Lipschitz: for any \(F,G \in E\), \(\|h_\Theta(F)-h_\Theta(G)\|_{\mathcal H}\leqslant L_{\mathcal T}\|F-G\|_{H^n}\), where \(L_{\mathcal T}=B_{\mathrm{in}}\prod_{\ell=0}^{L-1}L_\ell\) with \(L_\ell=\left(1+\xi_\ell L_{\mathrm{FF}}^{(\ell)}\right)\left(1+\eta_\ell L_{\mathrm{MHA}}^{(\ell)}\right)\). 
\end{proposition}

\begin{proof}
Let \(F=(f_0,\ldots,f_{n-1})\in E\) and \(G=(g_0,\ldots,g_{n-1}) \in E\) be two functional contexts. We propagate their difference through each stage of the transformer architecture.

\textit{Step 1: Functional embedding and initial projection.} For each context position \(j\), the unprojected embeddings are \(\widetilde v_j^{(0)}(F)=\mathsf T_{\mathrm{in}}f_j+p_j\) and \(\widetilde v_j^{(0)}(G)=\mathsf T_{\mathrm{in}}g_j+p_j\). Using the product norm on \(\mathcal H^n\),
\[
\begin{aligned}
\|V^{(0)}(F)-V^{(0)}(G)\|_{\mathcal H^n}^2
&=
\sum_{j=0}^{n-1}
\left\|
\Pi_R\bigl(\mathsf T_{\mathrm{in}}f_j+p_j\bigr)
-
\Pi_R\bigl(\mathsf T_{\mathrm{in}}g_j+p_j\bigr)
\right\|_{\mathcal H}^2\\
&\leqslant
\sum_{j=0}^{n-1}
\left\|
\mathsf T_{\mathrm{in}}(f_j-g_j)
\right\|_{\mathcal H}^2\\
&\leqslant
B_{\mathrm{in}}^2
\sum_{j=0}^{n-1}
\|f_j-g_j\|_H^2\\
&=
B_{\mathrm{in}}^2
\|F-G\|_{H^n}^2.
\end{aligned}
\]

The first equality follows from the definition of the product norm on \(\mathcal H^n\). The first inequality follows from Lemma \ref{lem:Non-expansiveness of the bounded projection}, together with the cancellation of the positional embeddings. The second inequality uses the operator-norm bound \(\|\mathsf T_{\mathrm{in}}\|_{\mathrm{op}}\leqslant B_{\mathrm{in}}\). The final equality is the definition of the product norm on \(H^n\).

\textit{Step 2: Attention residual block.} Let \(U,V\in\mathcal B_R\) be two inputs to layer \(\ell\). The attention residual update is \(Z(U)=\Pi_R^{(n)}\left(U+\eta_\ell\mathcal A_{\mathrm{MHA}}^{(\ell)}(U)\right)\), hence
\[
\begin{aligned}
\|Z(U)-Z(V)\|_{\mathcal H^n}
&\leqslant
\left\|
(U-V)
+
\eta_\ell
\left(
\mathcal A_{\mathrm{MHA}}^{(\ell)}(U)
-
\mathcal A_{\mathrm{MHA}}^{(\ell)}(V)
\right)
\right\|_{\mathcal H^n}\\
&\leqslant
\|U-V\|_{\mathcal H^n}
+
\eta_\ell
\left\|
\mathcal A_{\mathrm{MHA}}^{(\ell)}(U)
-
\mathcal A_{\mathrm{MHA}}^{(\ell)}(V)
\right\|_{\mathcal H^n}\\
&\leqslant
\left(
1+\eta_\ell L_{\mathrm{MHA}}^{(\ell)}
\right)
\|U-V\|_{\mathcal H^n}.
\end{aligned}
\]
The first inequality follows from Lemma \ref{lem:Non-expansiveness of the bounded projection}. The second inequality is the triangle inequality in \(\mathcal H^n\). The third inequality follows from the Lipschitz continuity of multi-head attention established in Corollary~\ref{cor:Lipschitz continuity of multi-head attention}.

\textit{Step 3: Feed-forward residual block.} The feed-forward residual update is \(V^{+}(Z)=\Pi_R^{(n)}\left(Z+\xi_\ell\mathcal F_\ell(Z)\right)\). For \(Z_1,Z_2\in\mathcal H^n\),
\[
\begin{aligned}
\|V^{+}(Z_1)-V^{+}(Z_2)\|_{\mathcal H^n}
&\leqslant
\left\|
(Z_1-Z_2)
+
\xi_\ell
\left(
\mathcal F_\ell(Z_1)
-
\mathcal F_\ell(Z_2)
\right)
\right\|_{\mathcal H^n}\\
&\leqslant
\|Z_1-Z_2\|_{\mathcal H^n}
+
\xi_\ell
\|\mathcal F_\ell(Z_1)-\mathcal F_\ell(Z_2)\|_{\mathcal H^n}\\
&\leqslant
\left(
1+\xi_\ell L_{\mathrm{FF}}^{(\ell)}
\right)
\|Z_1-Z_2\|_{\mathcal H^n}.
\end{aligned}
\]
The first inequality again follows from Lemma \ref{lem:Non-expansiveness of the bounded projection}. The second is the triangle inequality in \(\mathcal H^n\). The final inequality follows from the Lipschitz continuity of the feed-forward block established in Lemma~\ref{lem:Lipschitz continuity of the feed-forward block}.

Combining the attention and feed-forward residual blocks, we obtain
\[
\begin{aligned}
\|V^{(\ell+1)}(U)-V^{(\ell+1)}(V)\|_{\mathcal H^n}
&\leqslant
\left(
1+\xi_\ell L_{\mathrm{FF}}^{(\ell)}
\right)
\|Z(U)-Z(V)\|_{\mathcal H^n}\\
&\leqslant
\left(
1+\xi_\ell L_{\mathrm{FF}}^{(\ell)}
\right)
\left(
1+\eta_\ell L_{\mathrm{MHA}}^{(\ell)}
\right)
\|U-V\|_{\mathcal H^n}\\
&=
L_\ell\|U-V\|_{\mathcal H^n}.
\end{aligned}
\]
The first inequality applies the Lipschitz bound for the feed-forward residual block to the two attention outputs \(Z(U)\) and \(Z(V)\). The second inequality then applies the Lipschitz bound for the attention residual block. The final equality follows from the definition \(L_\ell=\left(1+\xi_\ell L_{\mathrm{FF}}^{(\ell)}\right)\left(1+\eta_\ell L_{\mathrm{MHA}}^{(\ell)}\right)\).

\textit{Step 4: Composition of the transformer layers.} Applying the layerwise estimate recursively gives
\[
\begin{aligned}
\|V^{(L)}(F)-V^{(L)}(G)\|_{\mathcal H^n}
&\leqslant
L_{L-1}
\|V^{(L-1)}(F)-V^{(L-1)}(G)\|_{\mathcal H^n}\\
&\leqslant
L_{L-1}L_{L-2}
\|V^{(L-2)}(F)-V^{(L-2)}(G)\|_{\mathcal H^n}\\
&\;\;\vdots\\
&\leqslant
\left(
\prod_{\ell=0}^{L-1}L_\ell
\right)
\|V^{(0)}(F)-V^{(0)}(G)\|_{\mathcal H^n}\\
&\leqslant
B_{\mathrm{in}}
\left(
\prod_{\ell=0}^{L-1}L_\ell
\right)
\|F-G\|_{H^n}\\
&=
L_{\mathcal T}\|F-G\|_{H^n}.
\end{aligned}
\]

The first three inequalities repeatedly apply the Lipschitz bound \(L_\ell\) for each transformer layer. The penultimate inequality uses the initial embedding bound established in Step 1. The final equality follows from the definition \(L_{\mathcal T}=B_{\mathrm{in}}\prod_{\ell=0}^{L-1}L_\ell\).

\textit{Step 5: Readout.} The final representation is obtained by applying the readout operator \(\mathsf R_{\mathrm{out}}:\mathcal H^n\to\mathcal H\), defined by \(\mathsf R_{\mathrm{out}}(v_0,\ldots,v_{n-1})=v_{n-1}\). Therefore
\[
\begin{aligned}
\|h_\Theta(F)-h_\Theta(G)\|_{\mathcal H}
&=
\left\|
\mathsf R_{\mathrm{out}}\bigl(V^{(L)}(F)-V^{(L)}(G)\bigr)
\right\|_{\mathcal H}\\
&\leqslant
\|\mathsf R_{\mathrm{out}}\|_{\mathrm{op}}
\|V^{(L)}(F)-V^{(L)}(G)\|_{\mathcal H^n}\\
&=
\|V^{(L)}(F)-V^{(L)}(G)\|_{\mathcal H^n}\\
&\leqslant
L_{\mathcal T}\|F-G\|_{H^n}.
\end{aligned}
\]

The first equality follows from the definition \(h_\Theta(F)=\mathsf R_{\mathrm{out}}(V^{(L)}(F))\) and linearity of \(\mathsf R_{\mathrm{out}}\). The first inequality is the definition of the operator norm. The second equality uses \(\|\mathsf R_{\mathrm{out}}\|_{\mathrm{op}}=1\), since \(\mathsf R_{\mathrm{out}}\) is the projection onto the final coordinate. The final inequality follows from Step 4.

Therefore \(h_\Theta:E\to\mathcal H\) is globally Lipschitz with Lipschitz constant at most \(L_{\mathcal T}\).
\end{proof}

Proposition~\ref{prop:Lipschitz continuity of the functional transformer} provides a finite constant \(L_{\mathcal T}\) controlling the sensitivity of the final latent representation. We now propagate this estimate through the coefficient heads and use the isometry of \(\mathsf E_d\).

\begin{corollary}[Lipschitz continuity of the Gaussian parameter maps]
\label{cor:Lipschitz continuity of the Gaussian parameter maps}
Let \(L_m^{\mathrm{head}}\) be a Lipschitz constant for
\(\widetilde{\mathsf T}_m:\mathcal H\to\mathbb R^d\), and let
\(L_{C,k}^{\mathrm{head}}\) be a Lipschitz constant for
\(\widetilde{\mathsf T}_{C,k}:\mathcal H\to\mathbb R^d\).
Then, the Gaussian parameter maps are Lipschitz:
\(\|m_\Theta(F)-m_\Theta(G)\|_H
\leq L_m\|F-G\|_{H^n}\),
where \(L_m=L_m^{\mathrm{head}}L_{\mathcal T}\), and
\(\|c_{k,\Theta}(F)-c_{k,\Theta}(G)\|_H
\leq L_{C,k}\|F-G\|_{H^n}\),
where \(L_{C,k}=L_{C,k}^{\mathrm{head}}L_{\mathcal T}\).
\end{corollary}

\begin{proof}
For the mean map, the isometry of \(\mathsf E_d\) gives
\[
\begin{aligned}
\|m_\Theta(F)-m_\Theta(G)\|_H
&=\|\mu_\Theta(F)-\mu_\Theta(G)\|_2\\
&\leqslant L_m^{\mathrm{head}}\|h_\Theta(F)-h_\Theta(G)\|_{\mathcal H}\\
&\leqslant L_m^{\mathrm{head}}L_{\mathcal T}\|F-G\|_{H^n}.
\end{aligned}
\]
The covariance-factor estimate is identical:
\[
\begin{aligned}
\|c_{k,\Theta}(F)-c_{k,\Theta}(G)\|_H
&=\|\gamma_{k,\Theta}(F)-\gamma_{k,\Theta}(G)\|_2\\
&\leqslant L_{C,k}^{\mathrm{head}}L_{\mathcal T}\|F-G\|_{H^n}.
\end{aligned}
\]
\end{proof}

We have therefore obtained explicit Lipschitz bounds for all of the quantities that determine the conditional Gaussian distribution produced by the transformer. The constants \(L_m\) and \(L_{C,k}\) are finite consequences of the bounded projections, finite operator-norm constraints, and Lipschitz nonlinearities.

The preceding corollary gives explicit control of the mean and covariance-factor maps in terms of the architecture. We can now translate these bounds into stability of the Gaussian output distribution in Wasserstein distance.

\begin{proposition}[Lipschitz continuity of the Gaussian output kernel]
\label{prop:Lipschitz continuity of the Gaussian output kernel}
For any \(F,G\in E\), we have the estimate \( W_2\bigl(Q_\Theta(F,\cdot),Q_\Theta(G,\cdot)\bigr)\leqslant L_\Theta\|F-G\|_\beta\), where \(L_\Theta^2=L_m^2+\sum_{k=1}^rL_{C,k}^2\).
\end{proposition}

\begin{proof}
By Proposition~\ref{prop:Wasserstein stability of the Gaussian output}, \(W_2^2\bigl(Q_\Theta(F,\cdot),Q_\Theta(G,\cdot)\bigr)\leqslant
\|m_\Theta(F)-m_\Theta(G)\|_H^2+\sum_{k=1}^r
\|c_{k,\Theta}(F)-c_{k,\Theta}(G)\|_H^2
\). Then,
\[
\begin{aligned}
W_2^2\bigl(Q_\Theta(F,\cdot),Q_\Theta(G,\cdot)\bigr)
&\leqslant
L_m^2\|F-G\|_{H^n}^2
+
\sum_{k=1}^r
L_{C,k}^2\|F-G\|_{H^n}^2\\
&=
\left(
L_m^2+\sum_{k=1}^rL_{C,k}^2
\right)
\|F-G\|_{H^n}^2\\
&=
L_\Theta^2\|F-G\|_{H^n}^2\\
&\leqslant
L_\Theta^2\|F-G\|_{\beta}^2.
\end{aligned}
\]
The first inequality applied Corollary~\ref{cor:Lipschitz continuity of the Gaussian parameter maps} to each term. The first equality factors out the common term \(\|F-G\|_{H^n}^2\), while the second follows from the definition \(L_\Theta^2=L_m^2+\sum_{k=1}^rL_{C,k}^2\). The last inequality follows from \(\beta>1\) implying \(\|F-G\|_{H^n}^2
=
\sum_{j=0}^{n-1}\|f_j-g_j\|_H^2
\leqslant
\sum_{j=0}^{n-1}\beta^j\|f_j-g_j\|_H^2
=
\|F-G\|_\beta^2.
\)
Taking square roots concludes the proof.
\end{proof}

The constant \(L_\Theta\) measures the sensitivity of the conditional Gaussian law to changes in the functional context. It is used in Section~\ref{subsec:Wasserstein and metric dependence} to obtain a Lipschitz constant for the pretraining loss and is also useful for perturbation estimates.


\end{document}
